% Generated mechanically from the frozen properties.tex target.
% Do not edit this generated file; edit the bound locale source instead.

% OK, start here.
%

\setcounter{chapter}{27}
\chapter{概形的性质}



\phantomsection
\label{properties-section-phantom}


\section{引言}
\label{properties-section-introduction}

\noindent
本章介绍概形的一些绝对性质。
基础参考文献是 \cite{EGA}。




\section{可构造集}
\label{properties-section-constructible}

\noindent
可构造集与局部可构造集在《拓扑》第
\ref{topology-section-constructible} 节中引入。
概形的局部可构造子集可刻画如下。

\begin{lemma}
\label{properties-lemma-locally-constructible}
令 $X$ 为概形。
子集 $E$ 在 $X$ 中局部可构造，当且仅当在 $X$ 上，
$E \cap U$ 在 $U$ 中可构造，其中 $U$ 取遍 $X$ 的仿射开集。
\end{lemma}

\begin{proof}
假设 $E$ 局部可构造。则存在开覆盖
$X = \bigcup U_i$，使得 $E \cap U_i$ 在 $U_i$ 中可构造，而这对每个 $i$ 成立。
令 $V \subset X$ 为任意仿射开集。可以找到有限仿射开覆盖
$V = V_1 \cup \ldots \cup V_m$，使得对每个 $j$，都有
$V_j \subset U_i$，其中 $i = i(j)$。由《拓扑》引理
\ref{topology-lemma-open-immersion-constructible-inverse-image}，
每个 $E \cap V_j$ 在 $V_j$ 中可构造。由于包含映射
$V_j \to V$ 拟紧（见《概形》引理
\ref{schemes-lemma-quasi-compact-affine}），由《拓扑》引理
\ref{topology-lemma-collate-constructible} 可知，$E \cap V$ 在 $V$ 中可构造。
逆向蕴含是立即的。
\end{proof}

\begin{lemma}
\label{properties-lemma-generic-point-in-constructible}
令 $X$ 为概形，$E \subset X$ 为局部可构造子集。
令 $\xi \in X$ 为 $X$ 的一个不可约分支的泛点。
\begin{enumerate}
\item 若 $\xi \in E$，则 $\xi$ 的某个开邻域包含于 $E$。
\item 若 $\xi \not \in E$，则 $\xi$ 的某个开邻域与 $E$ 不相交。
\end{enumerate}
\end{lemma}

\begin{proof}
局部可构造子集的补集仍局部可构造，故只需证明 (2)。可以假设 $X$
仿射，因而 $E$ 可构造（引理 \ref{properties-lemma-locally-constructible}）。
在这种情形下，$X$ 是谱空间
（《代数》引理 \ref{algebra-lemma-spec-spectral}）。于是，
$\xi \not \in E$ 蕴含 $\xi \not \in \overline{E}$；这是由《拓扑》引理
\ref{topology-lemma-constructible-stable-specialization-closed}
以及如下事实得出的：$X$ 中没有异于 $\xi$ 且特化到 $\xi$ 的点。
\end{proof}

\begin{lemma}
\label{properties-lemma-quasi-separated-quasi-compact-open-retrocompact}
令 $X$ 为拟分离概形。$X$ 的任意两个拟紧开集之交是 $X$ 的拟紧开集。
$X$ 的每个拟紧开集都在 $X$ 中逆拟紧。
\end{lemma}

\begin{proof}
若 $U$ 与 $V$ 是拟紧开集，则
$U \cap V = \Delta^{-1}(U \times V)$，其中 $\Delta : X \to X \times X$
是对角映射。由于 $X$ 拟分离，$\Delta$ 拟紧。可知 $U \cap V$ 拟紧，
因为 $U \times V$ 拟紧（略去细节；使用《概形》引理
\ref{schemes-lemma-affine-covering-fibre-product}
可知 $U \times V$ 是有限多个仿射概形之并）。
其余断言由第一个断言及《拓扑》引理
\ref{topology-lemma-topology-quasi-separated-scheme} 得出。
\end{proof}

\begin{lemma}
\label{properties-lemma-quasi-compact-quasi-separated-spectral}
令 $X$ 为拟紧且拟分离的概形。
则 $X$ 的底拓扑空间是谱空间。
\end{lemma}

\begin{proof}
根据《拓扑》定义 \ref{topology-definition-spectral-space}，
需要检验 $X$ 清醒、拟紧，具有由拟紧开集组成的拓扑基，并且任意两个拟紧开集
之交仍拟紧。这由《概形》引理 \ref{schemes-lemma-scheme-sober}、
\ref{schemes-lemma-basis-affine-opens}
以及上面的引理
\ref{properties-lemma-quasi-separated-quasi-compact-open-retrocompact}
得出。
\end{proof}

\begin{lemma}
\label{properties-lemma-constructible-quasi-compact-quasi-separated}
令 $X$ 为拟紧且拟分离的概形。
$X$ 的任意局部可构造子集都是可构造的。
\end{lemma}

\begin{proof}
由于 $X$ 拟紧，可以选取有限仿射开覆盖
$X = V_1 \cup \ldots \cup V_m$。由于 $X$ 拟分离，根据引理
\ref{properties-lemma-quasi-separated-quasi-compact-open-retrocompact}，每个 $V_i$
都在 $X$ 中逆拟紧。因此，由《拓扑》引理
\ref{topology-lemma-collate-constructible} 可知，$E \subset X$ 在 $X$
中可构造，当且仅当各个 $E \cap V_j$ 在 $V_j$ 中可构造。
于是由引理 \ref{properties-lemma-locally-constructible} 得到结论。
\end{proof}

\begin{lemma}
\label{properties-lemma-retrocompact}
令 $X$ 为概形。子集 $E$ 在 $X$ 中逆拟紧，当且仅当在 $X$ 上，
$E \cap U$ 拟紧，其中 $U$ 取遍 $X$ 的仿射开集。
\end{lemma}

\begin{proof}
这立即来自如下事实：$X$ 的每个拟紧开集都是有限多个仿射开集之并。
\end{proof}

\begin{lemma}
\label{properties-lemma-stratification-locally-finite-constructible}
一个分割 $X = \coprod_{i \in I} X_i$，其中 $X$ 是概形，若各部分均逆拟紧，
则该分割局部有限，当且仅当各部分均局部可构造。
\end{lemma}

\begin{proof}
关于逆拟紧、分割和局部有限的定义，见《拓扑》的定义
\ref{topology-definition-quasi-compact}、
\ref{topology-definition-paritition} 和
\ref{topology-definition-locally-finite}。

\medskip\noindent
若该分割局部有限，且 $U \subset X$ 是仿射开集，则
$U = \coprod_{i \in I} U \cap X_i$ 是有限分割（更准确地说，除有限多个部分外，
其余部分均为空）。因此 $U \cap X_i$ 拟紧，且它在 $U$ 中的补集作为有限多个
逆拟紧部分之并，仍逆拟紧。于是由《拓扑》引理
\ref{topology-lemma-locally-closed-constructible-image}，$U \cap X_i$
可构造。再由引理 \ref{properties-lemma-locally-constructible}，$X_i$ 局部可构造。

\medskip\noindent
假设各部分均局部可构造。则对任意仿射开集 $U \subset X$，得到由可构造子集
组成的覆盖 $U = \coprod X_i \cap U$。可构造拓扑是拟紧的，见《拓扑》引理
\ref{topology-lemma-constructible-hausdorff-quasi-compact}，
故此覆盖具有有限加细，即该分割局部有限。
\end{proof}




\section{整概形、不可约概形与约化概形}
\label{properties-section-integral}

\begin{definition}
\label{properties-definition-integral}
令 $X$ 为概形。若 $X$ 非空，并且对每个非空仿射开集
$\Spec(R) = U \subset X$，环 $R$ 都是整环，则称它是{it 整的}。
\end{definition}

\begin{lemma}
\label{properties-lemma-characterize-reduced}
令 $X$ 为概形。
下列条件等价。
\begin{enumerate}
\item 概形 $X$ 是约化的，见《概形》定义
\ref{schemes-definition-reduced}。
\item 存在仿射开覆盖 $X = \bigcup U_i$，使得每个
$\Gamma(U_i, \mathcal{O}_X)$ 都约化。
\item 对每个仿射开集 $U \subset X$，环
$\mathcal{O}_X(U)$ 都约化。
\item 对每个开集 $U \subset X$，环 $\mathcal{O}_X(U)$ 都约化。
\end{enumerate}
\end{lemma}

\begin{proof}
见《概形》引理 \ref{schemes-lemma-reduced} 和
\ref{schemes-lemma-affine-reduced}。
\end{proof}

\begin{lemma}
\label{properties-lemma-characterize-irreducible}
令 $X$ 为概形。
下列条件等价。
\begin{enumerate}
\item 概形 $X$ 不可约。
\item 存在仿射开覆盖 $X = \bigcup_{i \in I} U_i$，使得 $I$ 非空，
$U_i$ 对所有 $i \in I$ 都不可约，并且
$U_i \cap U_j \not = \emptyset$ 对所有 $i, j \in I$ 都成立。
\item 概形 $X$ 非空，并且每个非空仿射开集 $U \subset X$ 都不可约。
\end{enumerate}
\end{lemma}

\begin{proof}
假设 (1)。由《概形》引理 \ref{schemes-lemma-scheme-sober}，可知
$X$ 有唯一泛点 $\eta$。于是
$X = \overline{\{\eta\}}$。故 $\eta$ 属于每个非空仿射开集
$U \subset X$。这蕴含 $\eta \in U$ 稠密，因而 $U$ 不可约。
它还蕴含任意两个非空仿射开集相交。
因此 (1) 同时蕴含 (2) 与 (3)。

\medskip\noindent
假设 (2)。设 $X = Z_1 \cup Z_2$ 是两个闭子集之并。
对每个 $i$，要么 $U_i \subset Z_1$，要么 $U_i \subset Z_2$。
选取某个 $i \in I$，并假设 $U_i \subset Z_1$（必要时将
$Z_1$、$Z_2$ 重新编号）。对任意 $j \in I$，开子集
$U_i \cap U_j$ 在 $U_j$ 中稠密，并包含于闭子集 $Z_1 \cap U_j$。
由此可知 $U_j \subset Z_1$。故如所需，$X = Z_1$。

\medskip\noindent
假设 (3)。选取仿射开覆盖 $X = \bigcup_{i \in I} U_i$。
可以假设每个 $U_i$ 都非空。
由于 $X$ 非空，可知 $I$ 非空。
由假设，每个 $U_i$ 都不可约。
假设 $U_i \cap U_j = \emptyset$ 对某对 $i, j \in I$ 成立。
则开集 $U_i \amalg  U_j = U_i \cup U_j$ 是仿射的，见《概形》引理
\ref{schemes-lemma-disjoint-union-affines}。
故由假设它不可约，这与上式矛盾。因此 (3) 蕴含 (2)。引理得证。
\end{proof}

\begin{lemma}
\label{properties-lemma-characterize-integral}
概形 $X$ 是整的，当且仅当它既约化又不可约。
\end{lemma}

\begin{proof}
若 $X$ 不可约，则每个仿射开集 $\Spec(R) = U \subset X$
都不可约。若 $X$ 约化，则由上面的引理 \ref{properties-lemma-characterize-reduced}，
$R$ 约化。因此 $R$ 约化且 $(0)$ 是素理想，即 $R$ 是整环。

\medskip\noindent
若 $X$ 是整的，则对每个非空仿射开集
$\Spec(R) = U \subset X$，环 $R$ 都约化，
因而由引理 \ref{properties-lemma-characterize-reduced}，$X$ 约化。
此外，每个非空仿射开集都不可约。
因此 $X$ 不可约，见引理 \ref{properties-lemma-characterize-irreducible}。
\end{proof}

\noindent
在《例》第
\ref{examples-section-connected-locally-integral-not-integral} 节中，
我们构造了一个连通仿射概形，它的所有局部环都是整环，但它并不是整概形。










\section{由环的性质定义的概形类型}
\label{properties-section-properties-rings}

\noindent
本节研究环的哪些性质能够用来定义概形的局部性质。

\begin{definition}
\label{properties-definition-property-local}
令 $P$ 为环的一种性质。
若下列条件成立，则称 $P$ 是{it 局部的}：
\begin{enumerate}
\item 对任意环 $R$ 及任意 $f \in R$，有
$P(R) \Rightarrow P(R_f)$。
\item 对任意环 $R$ 及满足下列条件的 $f_i \in R$：
$(f_1, \ldots, f_n) = R$，有
$\forall i, P(R_{f_i}) \Rightarrow P(R)$。
\end{enumerate}
\end{definition}

\begin{definition}
\label{properties-definition-locally-P}
令 $P$ 为环的一种性质。令 $X$ 为概形。
称 $X${\it 局部具有 $P$ 性质}，如果对任意 $x \in X$，存在仿射开邻域
$U$，它是 $x$ 在 $X$ 中的邻域，并且 $\mathcal{O}_X(U)$ 具有性质 $P$。
\end{definition}

\noindent
只有当该性质是局部性质时，这才是一个合适的概念。
即使 $P$ 是局部性质，我们也不会自动使用这个定义，说某个概形
“局部具有 $P$ 性质”，除非我们还在别处明确给出该定义。

\begin{lemma}
\label{properties-lemma-locally-P}
令 $X$ 为概形。令 $P$ 为环的一种局部性质。
下列条件等价：
\begin{enumerate}
\item 概形 $X$ 局部具有 $P$ 性质。
\item 对每个仿射开集 $U \subset X$，性质
$P(\mathcal{O}_X(U))$ 成立。
\item 存在仿射开覆盖 $X = \bigcup U_i$，使得每个
$\mathcal{O}_X(U_i)$ 都满足 $P$。
\item 存在开覆盖 $X = \bigcup X_j$，使得每个开子概形 $X_j$
都局部具有 $P$ 性质。
\end{enumerate}
此外，若 $X$ 局部具有 $P$ 性质，则每个开子概形也局部具有 $P$ 性质。
\end{lemma}

\begin{proof}
显然 (1) $\Leftrightarrow$ (3)，且 (2) $\Rightarrow$ (1)。
若 (3) $\Rightarrow$ (2)，则引理的最后一个陈述成立，并容易推出 (4)
也等价于 (1)。因此只需证明 (3) $\Rightarrow$ (2)。

\medskip\noindent
令 $X = \bigcup U_i$ 为仿射开覆盖，记
$U_i = \Spec(R_i)$。假设 $P(R_i)$。
令 $\Spec(R) = U \subset X$ 为任意仿射开集。
由《概形》引理 \ref{schemes-lemma-good-subcover}，存在 $U = \Spec(R)$
的一个标准覆盖，由标准开集 $D(f_j)$ 组成，使得每个环 $R_{f_j}$
都是某个环 $R_i$ 的主局部化。由定义
\ref{properties-definition-property-local} (1)，得到 $P(R_{f_j})$。
进而由定义 \ref{properties-definition-property-local} (2)，得到 $P(R)$。
\end{proof}

\noindent
下面给出一个应用示例。

\begin{lemma}
\label{properties-lemma-reduced-is-locally-reduced}
令 $X$ 为概形。则 $X$ 约化，当且仅当 $X$ 在定义
\ref{properties-definition-locally-P} 的意义下“局部约化”。
\end{lemma}

\begin{proof}
这由引理 \ref{properties-lemma-characterize-reduced} 立即可见。
\end{proof}

\begin{lemma}
\label{properties-lemma-properties-local}
环 $R$ 的下列性质是局部性质。
\begin{enumerate}
\item （Cohen--Macaulay。）
环 $R$ 是 Noether 环且为 CM，见《代数》定义
\ref{algebra-definition-ring-CM}。
\item （正则。）
环 $R$ 是 Noether 环且正则，见《代数》定义
\ref{algebra-definition-regular}。
\item （绝对 Noether。）
环 $R$ 在 $Z$ 上有限型。
\item 可按需补充更多性质。\footnote{但这里只列出那些我们尚未在别处分别处理的性质。}
\end{enumerate}
\end{lemma}

\begin{proof}
略。
\end{proof}















\section{Noether 概形}
\label{properties-section-noetherian}

\noindent
回顾：若环 $R$ 满足理想的升链条件，则称它是{it Noether 环}。
等价地，$R$ 的每个理想都有限生成。

\begin{definition}
\label{properties-definition-noetherian}
令 $X$ 为概形。
\begin{enumerate}
\item 称 $X$ 是{it 局部 Noether 的}，如果每个 $x \in X$ 都有一个仿射开邻域
$\Spec(R) = U \subset X$，使得环 $R$ 是 Noether 环。
\item 若 $X$ 局部 Noether 且拟紧，则称 $X$ 是{it Noether 的}。
\end{enumerate}
\end{definition}

\noindent
下面是刻画局部 Noether 概形的标准结果。

\begin{lemma}
\label{properties-lemma-locally-Noetherian}
令 $X$ 为概形。下列条件等价：
\begin{enumerate}
\item 概形 $X$ 局部 Noether。
\item 对每个仿射开集 $U \subset X$，环 $\mathcal{O}_X(U)$
都是 Noether 环。
\item 存在仿射开覆盖 $X = \bigcup U_i$，使得每个
$\mathcal{O}_X(U_i)$ 都是 Noether 环。
\item 存在开覆盖 $X = \bigcup X_j$，使得每个开子概形 $X_j$
都局部 Noether。
\end{enumerate}
此外，若 $X$ 局部 Noether，则每个开子概形也局部 Noether。
\end{lemma}

\begin{proof}
只需证明 Noether 性是环的一种局部性质，见引理 \ref{properties-lemma-locally-P}。
Noether 环的任意局部化仍是 Noether 环，见《代数》引理
\ref{algebra-lemma-Noetherian-permanence}。
由《代数》引理 \ref{algebra-lemma-cover}，可知定义
\ref{properties-definition-property-local} 中的第二个性质成立。
\end{proof}

\begin{lemma}
\label{properties-lemma-immersion-into-noetherian}
若浸入 $Z \to X$ 的目标 $X$ 局部 Noether，则该浸入拟紧。
\end{lemma}

\begin{proof}
闭浸入显然拟紧。
拟紧态射的复合仍拟紧，见《拓扑》引理
\ref{topology-lemma-composition-quasi-compact}。
因此只需证明到局部 Noether 概形中的开浸入拟紧。
利用《概形》引理 \ref{schemes-lemma-quasi-compact-affine}，
可化约到 $X$ 仿射的情形。
Noether 环的谱的任意开子集都拟紧（例如，将《代数》引理
\ref{algebra-lemma-Noetherian-topology} 与《拓扑》引理
\ref{topology-lemma-Noetherian}、
\ref{topology-lemma-Noetherian-quasi-compact} 结合即可）。
\end{proof}

\begin{lemma}
\label{properties-lemma-locally-Noetherian-quasi-separated}
局部 Noether 概形是拟分离的。
\end{lemma}

\begin{proof}
由《概形》引理 \ref{schemes-lemma-characterize-quasi-separated}，
需要证明两个仿射开集之交 $U \cap V$ 在 $X$ 中拟紧。
例如，考虑开浸入 $U \cap V \to U$，结论就由上面的引理
\ref{properties-lemma-immersion-into-noetherian} 得出。
（其实这只不过是因为 Noether 环的谱的任意开集都拟紧。）
\end{proof}

\begin{lemma}
\label{properties-lemma-Noetherian-topology}
一个（局部）Noether 概形的底拓扑空间是（局部）Noether 的，
见《拓扑》定义 \ref{topology-definition-noetherian}。
\end{lemma}

\begin{proof}
这是因为 Noether 概形是有限多个 Noether 环的谱之并，并可应用
《代数》引理 \ref{algebra-lemma-Noetherian-topology} 与
《拓扑》引理 \ref{topology-lemma-finite-union-Noetherian}。
\end{proof}

\begin{lemma}
\label{properties-lemma-locally-closed-in-Noetherian}
一个（局部）Noether 概形的任意局部闭子概形都是（局部）Noether 的。
\end{lemma}

\begin{proof}
略。提示：Noether 环的任意商环与任意局部化都是 Noether 环。
对于 Noether 情形，再次使用如下事实：Noether 空间的任意子集在诱导拓扑下
都是 Noether 空间。
\end{proof}

\begin{lemma}
\label{properties-lemma-Noetherian-irreducible-components}
Noether 概形只有有限多个不可约分支。
\end{lemma}

\begin{proof}
Noether 概形的底拓扑空间是 Noether 的
（引理 \ref{properties-lemma-Noetherian-topology}），而 Noether 拓扑空间只有有限多个
不可约分支（《拓扑》引理 \ref{topology-lemma-Noetherian}），故结论成立。
\end{proof}

\begin{lemma}
\label{properties-lemma-morphism-Noetherian-schemes-quasi-compact}
概形态射 $f : X \to Y$ 若具有 Noether 源概形 $X$，则它拟紧。
\end{lemma}

\begin{proof}
使用引理 \ref{properties-lemma-Noetherian-topology}，并使用 Noether 拓扑空间的任意子集
都是拟紧的这一事实（见《拓扑》引理
\ref{topology-lemma-Noetherian} 和
\ref{topology-lemma-Noetherian-quasi-compact}）。
\end{proof}

\noindent
下面是一个有趣的引理。
它说明每个局部 Noether 概形都有大量闭点（每个闭子集中至少有一个）。

\begin{lemma}
\label{properties-lemma-locally-Noetherian-closed-point}
任意非空局部 Noether 概形都有闭点。
局部 Noether 概形的任意非空闭子集都有闭点。
等价地，局部 Noether 概形的任意一点都特化到一个闭点。
\end{lemma}

\begin{proof}
第二个断言由第一个断言得出（使用《概形》引理
\ref{schemes-lemma-reduced-closed-subscheme} 与引理
\ref{properties-lemma-locally-closed-in-Noetherian}）。
考虑任意非空仿射开集 $U \subset X$。
令 $x \in U$ 为闭点。若 $x$ 是 $X$ 的闭点，则已经完成。
否则，令 $X_0 \subset X$ 为 $\overline{\{x\}}$ 上的既约诱导闭子概形结构。
于是由《概形》引理 \ref{schemes-lemma-closed-subspace-scheme}，
$U_0 = U \cap X_0$ 是 $X_0$ 的仿射开集，并且
$U_0 = \{x\}$。令 $y \in X_0$、$y \not = x$ 为 $x$ 的一个特化。
考虑局部环 $R = \mathcal{O}_{X_0, y}$。
由于引理 \ref{properties-lemma-locally-closed-in-Noetherian} 表明 $X_0$ 是 Noether 的，
这是一个 Noether 局部环。记 $V \subset \Spec(R)$ 为 $U_0$ 在
$\Spec(R)$ 中通过典范态射 $\Spec(R) \to X_0$ 得到的逆像（见《概形》第
\ref{schemes-section-points} 节）。
按构造，$V$ 是单点集，其唯一点对应于 $x$（使用《概形》引理
\ref{schemes-lemma-specialize-points}）。
由《代数》引理
\ref{algebra-lemma-Noetherian-local-domain-dim-2-infinite-opens}，可知
$\dim(R) = 1$。
换言之，$y$ 是 $x$ 的直接特化（见《拓扑》定义
\ref{topology-definition-dimension-function}）。
再换言之，任意满足 $y \not = x$ 且 $x \leadsto y$ 的点，
都是 $x$ 的直接特化。显然这些点中的每一个都是闭点，如所需。
\end{proof}

\begin{lemma}
\label{properties-lemma-locally-Noetherian-specialization-dvr}
令 $X$ 为局部 Noether 概形。
令 $x' \leadsto x$ 为 $X$ 中两点之间的特化。则
\begin{enumerate}
\item 存在离散赋值环 $R$ 与态射
$f : \Spec(R) \to X$，使得泛点 $\eta$（它属于 $\Spec(R)$）
映到 $x'$，而特殊点映到 $x$；并且
\item 在 $x \not = x'$ 的前提下，给定有限生成域扩张
$K/\kappa(x')$，可以作出安排，使得扩张
$\kappa(\eta)/\kappa(x')$（它由 $f$ 诱导）与给定扩张同构。
\end{enumerate}
\end{lemma}

\begin{proof}
首先假设 $x' \leadsto x$ 是 $X$ 中的一个特化，且
$x \not = x'$；令 $K/\kappa(x')$ 为有限生成域扩张。由《概形》引理
\ref{schemes-lemma-specialize-points} 以及《概形》引理
\ref{schemes-lemma-characterize-points} 后面的讨论，得到环同态
$\mathcal{O}_{X, x} \to \kappa(x') \to K$。
由于 $x \not = x'$，$\mathcal{O}_{X, x}$ 在 $K$ 中的像不是域（略去细节）。
令 $R \subset K$ 为任意分式域为 $K$ 且支配
$\mathcal{O}_{X, x} \to K$ 的像的离散赋值环，见《代数》引理
\ref{algebra-lemma-exists-dvr}。
环同态 $\mathcal{O}_{X, x} \to R$ 诱导态射
$f : \Spec(R) \to X$，见《概形》引理
\ref{schemes-lemma-morphism-from-spec-local-ring}。
按构造，该态射具有所需的全部性质。
若 $x = x'$，则可以取 $R = \kappa(x)[t]_{(t)}$。
\end{proof}

\begin{lemma}
\label{properties-lemma-thin-infinite-sequence}
令 $S$ 为 Noether 概形。令 $T \subset S$ 为无限子集。
则存在无限子集 $T' \subset T$，使得 $T'$ 的各点之间不存在非平凡特化。
\end{lemma}

\begin{proof}
令 $T_0 \subset T$ 为由 $t \in T$ 组成的集合，其中这些点不特化到 $T$ 中另一点。
若 $T_0$ 无限，则取 $T' = T_0$ 即可。因此可以且确实假设 $T_0$ 有限。
归纳地，对 $i > 0$，考虑集合 $T_i \subset T$，它由满足下列条件的
$t \in T$ 组成：
\begin{enumerate}
\item $t \not \in T_{i - 1} \cup T_{i - 2} \cup \ldots \cup T_0$；
\item 存在非平凡特化 $t \leadsto t'$，且 $t' \in T_{i - 1}$；并且
\item 对任意非平凡特化 $t \leadsto t'$，若 $t' \in T$，则有
$t' \in T_{i - 1} \cup T_{i - 2} \cup \ldots \cup T_0$。
\end{enumerate}
同样，若 $T_i$ 无限，则取 $T' = T_i$ 即可。
令 $d$ 为各局部环 $\mathcal{O}_{S, t}$（其中 $t \in T_0$）维数的最大值；
那么 $d$ 是整数，因为 $T_0$ 有限，并且由《代数》命题
\ref{algebra-proposition-dimension} 可知这些局部环的维数有限。
于是，$T_i = \emptyset$ 对 $i > d$ 成立。
事实上，若 $t \in T_i$，则可以找到非平凡特化序列
$t = t_i \leadsto t_{i - 1} \leadsto \ldots \leadsto t_0$，
其中 $t_0 \in T_0$。由于各点
$t = t_i, t_{i - 1}, \ldots, t_0$ 均属于
$\Spec(\mathcal{O}_{S, t_0})$
（《概形》引理 \ref{schemes-lemma-specialize-points}），可知 $i \leq d$。
因此 $\bigcup T_i = T_d \cup \ldots \cup T_0$ 是 $T$ 的有限子集。

\medskip\noindent
假设 $t \in T$ 不属于 $\bigcup T_i$。则必存在特化
$t \leadsto t'$，其中 $t' \in T$ 且 $t' \not \in \bigcup T_i$。
（事实上，若 $t$ 的每个特化都属于有限集 $T_d \cup \ldots \cup T_0$，
则存在最大的 $i$，使得有某个特化 $t \leadsto t'$ 满足 $t' \in T_i$；
此时按构造 $t \in T_{i + 1}$。）
于是得到无限序列
$$
t \leadsto t' \leadsto t'' \leadsto \ldots
$$
，它由 $T \setminus \bigcup T_i$ 中各点之间的非平凡特化组成。
这是不可能的，因为由引理 \ref{properties-lemma-locally-Noetherian-quasi-separated}，
$S$ 的底拓扑空间是 Noether 的。
\end{proof}

\begin{lemma}
\label{properties-lemma-maximal-points}
令 $S$ 为 Noether 概形。令 $T \subset S$ 为子集。令
$T_0 \subset T$ 为由如下 $t \in T$ 组成的集合：不存在满足
非平凡特化 $t' \leadsto t$，其中 $t' \in T'$。则：(a) $T_0$ 的各点之间
不存在特化；(b) $T$ 的每一点都是 $T_0$ 中某一点的特化；(c) $T$ 与
$T_0$ 的闭包相同。
\end{lemma}

\begin{proof}
回顾：$\dim(\mathcal{O}_{S, s}) < \infty$ 对任意 $s \in S$ 成立，
见《代数》命题 \ref{algebra-proposition-dimension}。令 $t \in T$。
若 $t' \leadsto t$，则由维数理论，
$\dim(\mathcal{O}_{S, t'}) \leq \dim(\mathcal{O}_{S, t})$，
且等号成立当且仅当 $t' = t$。因此，若选取 $t' \leadsto t$，使得
$\dim(\mathcal{O}_{T, t'})$ 最小，则 $t' \in T_0$。
换言之，每个 $t \in T$ 都是 $T_0$ 中某元素的特化。
\end{proof}

\begin{lemma}
\label{properties-lemma-countable-dense-subset}
令 $S$ 为 Noether 概形。令 $T \subset S$ 为无限稠密子集。
则存在可数子集 $E \subset T$，它在 $S$ 中稠密。
\end{lemma}

\begin{proof}
令 $T'$ 为由如下 $s \in S$ 组成的集合：$\overline{\{s\}} \cap T$
含有一个可数子集，其闭包是 $\overline{\{s\}}$。
由于有限集可数，有 $T \subset T'$。
对 $s \in T'$，选取这样的可数子集
$E_s \subset \overline{\{s\}} \cap T$。
令 $E' = \{s_1, s_2, s_3, \ldots\} \subset T'$ 为可数子集。
则 $E'$ 在 $S$ 中的闭包，就是可数子集 $\bigcup_n E_{s_n}$ 的闭包，
而该子集包含于 $T$。
由此可知，若 $Z$ 是 $E'$ 的闭包的一个不可约分支，则 $Z$ 的泛点属于 $T'$。

\medskip\noindent
记 $T'_0 \subset T'$ 为由如下 $t \in T'$ 组成的子集：不存在满足
非平凡特化 $t' \leadsto t$，其中 $t' \in T'$；这正如引理
\ref{properties-lemma-maximal-points}，下文将直接使用该引理的结果。
若 $T'_0$ 无限，则选取可数子集 $E' \subset T'_0$。
由第一段的论证，$E'$ 的闭包的各个不可约分支的泛点都属于 $T'$。
然而，由于其中一点特化到 $E' \subset T'_0$ 中无限多个互不相同的元素，
这产生矛盾。因此 $T'_0$ 有限，设
$T'_0 = \{s_1, \ldots, s_m\}$。于是，$S$ 是 $T$ 的闭包，且包含于
$\{s_1, \ldots, s_m\}$ 的闭包；后者又包含于可数子集
$E_{s_1} \cup \ldots \cup E_{s_m} \subset T$ 的闭包，如所需。
\end{proof}








\section{Jacobson 概形}
\label{properties-section-jacobson}

\noindent
回顾：若闭点在每个闭子集中都稠密，则称一个空间是{it Jacobson 的}，
见《拓扑》第 \ref{topology-section-space-jacobson} 节。

\begin{definition}
\label{properties-definition-jacobson}
称概形 $S$ 是{it Jacobson 概形}，如果它的底拓扑空间是 Jacobson 的。
\end{definition}

\noindent
回顾：若环 $R$ 的每个根理想都是极大理想之交，则称 $R$ 是 Jacobson 环，
见《代数》定义 \ref{algebra-definition-ring-jacobson}。

\begin{lemma}
\label{properties-lemma-affine-jacobson}
仿射概形 $\Spec(R)$ 是 Jacobson 概形，当且仅当环 $R$ 是 Jacobson 环。
\end{lemma}

\begin{proof}
这就是《代数》引理 \ref{algebra-lemma-jacobson}。
\end{proof}

\noindent
下面是刻画 Jacobson 概形的标准结果。
直观地说，它断言 Jacobson $\Leftrightarrow$ 局部 Jacobson。

\begin{lemma}
\label{properties-lemma-locally-jacobson}
令 $X$ 为概形。下列条件等价：
\begin{enumerate}
\item 概形 $X$ 是 Jacobson 概形。
\item 概形 $X$ 在定义 \ref{properties-definition-locally-P} 的意义下“局部 Jacobson”。
\item 对每个仿射开集 $U \subset X$，环 $\mathcal{O}_X(U)$ 是 Jacobson 环。
\item 存在仿射开覆盖 $X = \bigcup U_i$，使得每个 $\mathcal{O}_X(U_i)$
都是 Jacobson 环。
\item 存在开覆盖 $X = \bigcup X_j$，使得每个开子概形 $X_j$
都是 Jacobson 概形。
\end{enumerate}
此外，若 $X$ 是 Jacobson 概形，则每个开子概形也都是 Jacobson 概形。
\end{lemma}

\begin{proof}
引理的最后一个断言由《拓扑》引理
\ref{topology-lemma-jacobson-inherited} 得出。
(5) 与 (1) 的等价性就是《拓扑》引理
\ref{topology-lemma-jacobson-local}。
因此，使用引理 \ref{properties-lemma-affine-jacobson}，可知
(1) $\Leftrightarrow$ (2)。
要完成引理的证明，只需证明“Jacobson”是环的一种局部性质，见引理
\ref{properties-lemma-locally-P}。
Jacobson 环关于一个元素的任意局部化仍为 Jacobson 环，见《代数》引理
\ref{algebra-lemma-Jacobson-invert-element}。
设 $R$ 为环，$f_1, \ldots, f_n \in R$ 生成单位理想，并且每个
$R_{f_i}$ 都是 Jacobson 环。于是
$\Spec(R) = \bigcup D(f_i)$ 是若干开子集之并，这些开子集全都是 Jacobson 的；
因此再次由《拓扑》引理 \ref{topology-lemma-jacobson-local}，
$\Spec(R)$ 是 Jacobson 概形。
这证明了定义 \ref{properties-definition-property-local} 的第二个性质。
\end{proof}

\noindent
代数几何中常用的许多概形都是 Jacobson 概形，见《态射》引理
\ref{morphisms-lemma-ubiquity-Jacobson-schemes}。
这里指出以下有趣的情形。

\begin{lemma}
\label{properties-lemma-complement-closed-point-Jacobson}
以下是 Noether Jacobson 概形的例子。
\begin{enumerate}
\item 若 $(R, \mathfrak m)$ 是 Noether 局部环，则穿孔谱
$\Spec(R) \setminus \{\mathfrak m\}$ 是 Jacobson 概形。
\item 若 $R$ 是 Jacobson 根为 $\text{rad}(R)$ 的 Noether 环，
则 $\Spec(R) \setminus V(\text{rad}(R))$ 是 Jacobson 概形。
\item 若 $(R, I)$ 是 Zariski 对（《更多代数》定义
\ref{more-algebra-definition-zariski-pair}），且 $R$ 是 Noether 环，
则 $\Spec(R) \setminus V(I)$ 是 Jacobson 概形。
\end{enumerate}
\end{lemma}

\begin{proof}
证明 (3)。注意，$\Spec(R) - V(I)$ 由仿射开集
$\Spec(R_f)$（其中 $f \in I$）覆盖。由《更多代数》引理
\ref{more-algebra-lemma-noetherian-zariski-jacobson-complement}，
环 $R_f$ 是 Jacobson 环。
因此由引理 \ref{properties-lemma-locally-jacobson}，
$\Spec(R) \setminus V(I)$ 是 Jacobson 概形。
(1) 与 (2) 是 (3) 的特殊情形。

\medskip\noindent
下面直接证明情形 (1)。
由于 $\Spec(R)$ 是 Noether 概形，
$S$ 是 Noether 概形（引理 \ref{properties-lemma-locally-closed-in-Noetherian}）。
因此 $S$ 是清醒的 Noether 拓扑空间（使用《概形》引理
\ref{schemes-lemma-scheme-sober}）。
假设 $S$ 不是 Jacobson 概形，以导出矛盾。
由《拓扑》引理
\ref{topology-lemma-non-jacobson-Noetherian-characterize}，
存在某个非闭点 $\xi \in S$，使得 $\{\xi\}$ 局部闭。
这对应于素理想 $\mathfrak p \subset R$，并且：(1) 存在素理想
$\mathfrak q$，满足 $\mathfrak p \subset \mathfrak q \subset \mathfrak m$，
两个包含关系都是严格的；(2) $\{\mathfrak p\}$ 在
$\Spec(R/\mathfrak p)$ 中是开的。这与《代数》引理
\ref{algebra-lemma-Noetherian-local-domain-dim-2-infinite-opens} 矛盾。
\end{proof}








\section{正规概形}
\label{properties-section-normal}

\noindent
回顾：若环 $R$ 的所有局部环都是正规整环，则称该环正规，
见《代数》定义 \ref{algebra-definition-ring-normal}。
正规整环是指在其分式域中整闭的整环，见《代数》定义
\ref{algebra-definition-domain-normal}。
因此可以如下定义正规概形。

\begin{definition}
\label{properties-definition-normal}
概形 $X$ 正规，当且仅当对所有 $x \in X$，局部环
$\mathcal{O}_{X, x}$ 都是正规整环。
\end{definition}

\noindent
这似乎是 EGA 所采用的定义，见 \cite[0, 4.1.4]{EGA}。
假设 $X = \Spec(A)$，且 $A$ 约化。此时，说 $X$ 正规并不等价于说
$A$ 在其总分式环中整闭。然而，若 $A$ 是 Noether 环，则二者等价
（见《代数》引理 \ref{algebra-lemma-characterize-reduced-ring-normal}）。

\begin{lemma}
\label{properties-lemma-locally-normal}
令 $X$ 为概形。下列条件等价：
\begin{enumerate}
\item 概形 $X$ 正规。
\item 对每个仿射开集 $U \subset X$，环 $\mathcal{O}_X(U)$ 正规。
\item 存在仿射开覆盖 $X = \bigcup U_i$，使得每个 $\mathcal{O}_X(U_i)$
都正规。
\item 存在开覆盖 $X = \bigcup X_j$，使得每个开子概形 $X_j$ 都正规。
\end{enumerate}
此外，若 $X$ 正规，则每个开子概形也正规。
\end{lemma}

\begin{proof}
这由定义立即可见。
\end{proof}

\begin{lemma}
\label{properties-lemma-normal-reduced}
正规概形是约化的。
\end{lemma}

\begin{proof}
这立即来自定义。
\end{proof}

\begin{lemma}
\label{properties-lemma-integral-normal}
令 $X$ 为整概形。
则 $X$ 正规，当且仅当对每个非空仿射开集
$U \subset X$，环 $\mathcal{O}_X(U)$ 都是正规整环。
\end{lemma}

\begin{proof}
这由《代数》引理 \ref{algebra-lemma-normality-is-local} 得出。
\end{proof}

\begin{lemma}
\label{properties-lemma-normal-locally-finite-nr-irreducibles}
令 $X$ 为概形，并假设任意拟紧开集都只有有限多个不可约分支。
下列条件等价：
\begin{enumerate}
\item $X$ 正规；
\item $X$ 是若干正规整概形的不交并。
\end{enumerate}
\end{lemma}

\begin{proof}
由定义立即可见，(2) 蕴含 (1)。
令 $X$ 为正规概形，并假设每个拟紧开集都只有有限多个不可约分支。
若 $X$ 仿射，则由《代数》引理
\ref{algebra-lemma-characterize-reduced-ring-normal}，$X$ 满足 (2)。
对于一般的 $X$，令 $X = \bigcup X_i$ 为仿射开覆盖。
注意，每个 $X_i$ 也只有有限多个不可约分支，并且该引理对每个 $X_i$ 成立。
令 $T \subset X$ 为一个不可约分支。由仿射情形，每个交
$T \cap X_i$ 在 $X_i$ 中是开的，并且是正规整概形。
因此 $T \subset X$ 是开子概形，并且是正规整概形。
这证明 $X$ 是其不可约分支的不交并，而这些分支都是正规整概形。
\end{proof}

\begin{lemma}
\label{properties-lemma-normal-Noetherian}
令 $X$ 为 Noether 概形。
下列条件等价：
\begin{enumerate}
\item $X$ 正规；
\item $X$ 是有限多个正规整概形的不交并。
\end{enumerate}
\end{lemma}

\begin{proof}
这是引理 \ref{properties-lemma-normal-locally-finite-nr-irreducibles} 的特殊情形，
因为 Noether 概形具有 Noether 底拓扑空间（引理
\ref{properties-lemma-Noetherian-topology} 以及《拓扑》引理
\ref{topology-lemma-Noetherian}）。
\end{proof}

\begin{lemma}
\label{properties-lemma-normal-locally-Noetherian}
令 $X$ 为局部 Noether 概形。
下列条件等价：
\begin{enumerate}
\item $X$ 正规；
\item $X$ 是若干正规整概形的不交并。
\end{enumerate}
\end{lemma}

\begin{proof}
略。提示：由引理 \ref{properties-lemma-normal-Noetherian}，这纯粹是一个拓扑事实。
\end{proof}

\begin{remark}
\label{properties-remark-normal-connected-irreducible}
令 $X$ 为正规概形。若 $X$ 局部 Noether，则可知
$X$ 是整概形，当且仅当 $X$ 连通，见引理
\ref{properties-lemma-normal-locally-Noetherian}。
但是，存在连通仿射概形 $X$，使得 $\mathcal{O}_{X, x}$ 对所有
$x \in X$ 都是整环，而 $X$ 不可约，见《例》第
\ref{examples-section-connected-locally-integral-not-integral} 节。
这个例子甚至还是正规概形（略去证明），务请注意！
\end{remark}

\begin{lemma}
\label{properties-lemma-normal-integral-sections}
\begin{slogan}
正规概形上的函数环是正规环。
\end{slogan}
令 $X$ 为正规整概形。
则 $\Gamma(X, \mathcal{O}_X)$ 是正规整环。
\end{lemma}

\begin{proof}
置 $R = \Gamma(X, \mathcal{O}_X)$。
显然 $R$ 是整环。
假设 $f = a/b$ 是其分式域中的元素，并且在 $R$ 上整。设有
$f^d + \sum_{i = 0, \ldots, d - 1} a_i f^i = 0$，其中
$a_i \in R$。令 $U \subset X$ 为非空仿射开集。
由于 $b \in R$ 非零，并且 $X$ 是整概形，可知
$b|_U \in \mathcal{O}_X(U)$ 也非零。
因此 $a/b$ 是 $\mathcal{O}_X(U)$ 的分式域中的元素，并且在
$\mathcal{O}_X(U)$ 上整（因为可以使用同一个多项式
$f^d + \sum_{i = 0, \ldots, d - 1} a_i|_U f^i = 0$，在 $U$ 上亦然）。
由于 $\mathcal{O}_X(U)$ 是正规整环
（引理 \ref{properties-lemma-integral-normal}），可知
$f_U = (a|_U)/(b|_U) \in \mathcal{O}_X(U)$。显然，
$f_U|_V = f_V$ 只要 $V \subset U \subset X$ 是非空仿射开集。
因此各局部截面 $f_U$ 粘合成一个元素
$g \in R = \Gamma(X, \mathcal{O}_X)$。于是 $bg$ 与 $a$
限制为 $\mathcal{O}_X(U)$ 的同一元素，对所有上述 $U$ 都如此，故
$bg = a$；换言之，$g$ 映到 $f$，而这是在 $R$ 的分式域中。
\end{proof}


\section{Cohen-Macaulay 概形}
\label{properties-section-Cohen-Macaulay}

\noindent
回顾《代数》定义 \ref{algebra-definition-local-ring-CM}：若 Noether 局部环
$(R, \mathfrak m)$ 满足
$\text{depth}_{\mathfrak m}(R) = \dim(R)$，则称其为 Cohen-Macaulay 环。
又回顾：若 Noether 环 $R$ 的每个局部环 $R_{\mathfrak p}$ 都是
Cohen-Macaulay 环，则称 $R$ 为 Cohen-Macaulay 环，见《代数》定义
\ref{algebra-definition-ring-CM}。

\begin{definition}
\label{properties-definition-Cohen-Macaulay}
令 $X$ 为概形。称 $X$ 是{\it Cohen-Macaulay 概形}，如果对每个
$x \in X$，都存在仿射开邻域 $U \subset X$ 包含 $x$，使得环
$\mathcal{O}_X(U)$ 是 Noether 环且为 Cohen-Macaulay 环。
\end{definition}

\begin{lemma}
\label{properties-lemma-characterize-Cohen-Macaulay}
令 $X$ 为概形。下列条件等价：
\begin{enumerate}
\item $X$ 是 Cohen-Macaulay 概形；
\item $X$ 局部 Noether，且它的所有局部环都是 Cohen-Macaulay 环；
以及
\item $X$ 局部 Noether，且对任意闭点 $x \in X$，局部环
$\mathcal{O}_{X, x}$ 都是 Cohen-Macaulay 环。
\end{enumerate}
\end{lemma}

\begin{proof}
《代数》引理 \ref{algebra-lemma-localize-CM} 说明，Cohen-Macaulay
局部环的局部化仍是 Cohen-Macaulay 环。将此与引理
\ref{properties-lemma-locally-Noetherian}、局部 Noether 概形上闭点的存在性
（引理 \ref{properties-lemma-locally-Noetherian-closed-point}）以及各定义结合，
即得本引理。
\end{proof}

\begin{lemma}
\label{properties-lemma-locally-Cohen-Macaulay}
令 $X$ 为概形。下列条件等价：
\begin{enumerate}
\item 概形 $X$ 是 Cohen-Macaulay 概形。
\item 对每个仿射开集 $U \subset X$，环 $\mathcal{O}_X(U)$
都是 Noether 环且为 Cohen-Macaulay 环。
\item 存在仿射开覆盖 $X = \bigcup U_i$，使得每个
$\mathcal{O}_X(U_i)$ 都是 Noether 环且为 Cohen-Macaulay 环。
\item 存在开覆盖 $X = \bigcup X_j$，使得每个开子概形 $X_j$
都是 Cohen-Macaulay 概形。
\end{enumerate}
此外，若 $X$ 是 Cohen-Macaulay 概形，则每个开子概形也都是
Cohen-Macaulay 概形。
\end{lemma}

\begin{proof}
结合引理 \ref{properties-lemma-locally-Noetherian}
和 \ref{properties-lemma-characterize-Cohen-Macaulay}。
\end{proof}

\noindent
关于 Cohen-Macaulay 概形和深度的更多信息，见《概形上同调》第
\ref{coherent-section-depth} 节。








\section{正则概形}
\label{properties-section-regular}

\noindent
回顾《代数》定义 \ref{algebra-definition-regular-local}：若 Noether
局部环 $(R, \mathfrak m)$ 的理想 $\mathfrak m$ 可由
$\dim(R)$ 个元素生成，则称该局部环{\it 正则}。
又回顾：若 Noether 环 $R$ 的每个局部环 $R_{\mathfrak p}$ 都正则，
则称 $R$ 为{\it 正则环}，见《代数》定义
\ref{algebra-definition-regular}。

\begin{definition}
\label{properties-definition-regular}
令 $X$ 为概形。称 $X$ 是{\it 正则概形}，或称{\it 非奇异概形}，
如果对每个 $x \in X$，都存在仿射开邻域 $U \subset X$ 包含 $x$，
使得环 $\mathcal{O}_X(U)$ 是 Noether 正则环。
\end{definition}

\begin{lemma}
\label{properties-lemma-characterize-regular}
令 $X$ 为概形。下列条件等价：
\begin{enumerate}
\item $X$ 正则；
\item $X$ 局部 Noether，且它的所有局部环都正则；
以及
\item $X$ 局部 Noether，且对任意闭点 $x \in X$，局部环
$\mathcal{O}_{X, x}$ 都正则。
\end{enumerate}
\end{lemma}

\begin{proof}
由《代数》定义 \ref{algebra-definition-regular} 之前的讨论可知，
正则局部环的局部化仍正则。将此与引理
\ref{properties-lemma-locally-Noetherian}、局部 Noether 概形上闭点的存在性
（引理 \ref{properties-lemma-locally-Noetherian-closed-point}）以及各定义结合，
即得本引理。
\end{proof}

\begin{lemma}
\label{properties-lemma-locally-regular}
令 $X$ 为概形。下列条件等价：
\begin{enumerate}
\item 概形 $X$ 正则。
\item 对每个仿射开集 $U \subset X$，环 $\mathcal{O}_X(U)$
都是 Noether 正则环。
\item 存在仿射开覆盖 $X = \bigcup U_i$，使得每个
$\mathcal{O}_X(U_i)$ 都是 Noether 正则环。
\item 存在开覆盖 $X = \bigcup X_j$，使得每个开子概形 $X_j$
都正则。
\end{enumerate}
此外，若 $X$ 正则，则每个开子概形也正则。
\end{lemma}

\begin{proof}
结合引理 \ref{properties-lemma-locally-Noetherian}
和 \ref{properties-lemma-characterize-regular}。
\end{proof}

\begin{lemma}
\label{properties-lemma-regular-normal}
正则概形是正规的。
\end{lemma}

\begin{proof}
见《代数》引理 \ref{algebra-lemma-regular-normal}。
\end{proof}





\section{维数}
\label{properties-section-dimension}

\noindent
概形的维数就是其底拓扑空间的维数。

\begin{definition}
\label{properties-definition-dimension}
令 $X$ 为概形。
\begin{enumerate}
\item $X$ 的{\it 维数}就是 $X$ 作为拓扑空间的维数，见《拓扑》定义
\ref{topology-definition-Krull}。
\item 对 $x \in X$，按照《拓扑》定义
\ref{topology-definition-Krull}，以 $\dim_x(X)$ 表示 $X$ 的底拓扑空间
在 $x$ 处的维数。称 $\dim_x(X)$ 为{\it $X$ 在 $x$ 处的维数}。
\end{enumerate}
\end{definition}

\noindent
由于概形的底拓扑空间是清醒空间
（《概形》引理 \ref{schemes-lemma-scheme-sober}），我们可以把 $X$
的维数算作下列不可约闭子集链的长度 $n$ 的上确界：
$$
T_0 \subset T_1 \subset \ldots \subset T_n
$$
其中各集合都是 $X$ 的不可约闭子集；也可以把它算作下列点的特化链
的长度 $n$ 的上确界：
$$
\xi_n \leadsto \xi_{n - 1} \leadsto \ldots \leadsto \xi_0
$$
其中各点属于 $X$。

\begin{lemma}
\label{properties-lemma-dimension}
令 $X$ 为概形。下列三个量相等：
\begin{enumerate}
\item $X$ 的维数；
\item $X$ 的各局部环维数的上确界；
\item $\dim_x(X)$ 在 $x \in X$ 时的上确界。
\end{enumerate}
\end{lemma}

\begin{proof}
注意，给定 $X$ 中点的特化链
$$
\xi_n \leadsto \xi_{n - 1} \leadsto \ldots \leadsto \xi_0
$$
由《概形》引理 \ref{schemes-lemma-specialize-points}，所有点
$\xi_i$ 都对应于 $X$ 在 $\xi_0$ 处的局部环的素理想。
因此，$X$ 的维数就是其各局部环维数的上确界。特别地，
$\dim_x(X) \geq \dim(\mathcal{O}_{X, x})$，因为 $\dim_x(X)$
是 $x$ 的各开邻域维数的最小值。因此
$\sup_{x \in X} \dim_x(X) \geq \dim(X)$。另一方面，显然
$\sup_{x \in X} \dim_x(X) \leq \dim(X)$，因为
$\dim(U) \leq \dim(X)$ 对 $X$ 的任意开子集都成立。
\end{proof}

\begin{lemma}
\label{properties-lemma-codimension-local-ring}
令 $X$ 为概形，$Y \subset X$ 为不可约闭子集，并令 $\xi \in Y$
为其泛点。则
$$
\text{codim}(Y, X) = \dim(\mathcal{O}_{X, \xi})
$$
其中余维按《拓扑》定义 \ref{topology-definition-codimension} 定义。
\end{lemma}

\begin{proof}
由《拓扑》引理 \ref{topology-lemma-codimension-at-generic-point}，
可以把 $X$ 替换为 $\xi$ 的一个仿射开邻域。在这种情形下，结论容易由
《代数》引理 \ref{algebra-lemma-irreducible-components-containing-x} 得出。
\end{proof}

\begin{lemma}
\label{properties-lemma-generic-point}
令 $X$ 为概形，$x \in X$。则 $x$ 是 $X$ 的某个不可约分支的泛点，
当且仅当 $\dim(\mathcal{O}_{X, x}) = 0$。
\end{lemma}

\begin{proof}
例如，这可由引理 \ref{properties-lemma-codimension-local-ring} 得出。
\end{proof}

\begin{lemma}
\label{properties-lemma-locally-Noetherian-dimension-0}
维数为 $0$ 的局部 Noether 概形是不交并；其每个分支都是一个
Artin 局部环的谱。
\end{lemma}

\begin{proof}
维数为 $0$ 的 Noether 环是有限多个 Artin 局部环的直积，见《代数》命题
\ref{algebra-proposition-dimension-zero-ring}。因此，局部 Noether 概形
$X$ 的维数为 $0$ 时，其任意仿射开集都有离散的底拓扑空间。这说明 $X$
上的拓扑是离散的。本引理容易由这些观察得出。
\end{proof}

\begin{lemma}
\label{properties-lemma-dimension-zero}
\begin{reference}
Ofer Gabber 于 2016 年 6 月 4 日发来的电子邮件
\end{reference}
令 $X$ 为零维概形。下列条件等价：
\begin{enumerate}
\item $X$ 拟分离；
\item $X$ 分离；
\item $X$ 是豪斯多夫空间；
\item 每个仿射开集都是闭集。
\end{enumerate}
在这种情形下，$X$ 的连通分支都是点，并且 $X$ 的每个拟紧开集都仿射。
特别地，若 $X$ 拟紧，则 $X$ 仿射。
\end{lemma}

\begin{proof}
由于 $X$ 的维数为零，对任意仿射开集 $U \subset X$，空间 $U$
都是预有限空间，并且满足我们将在下文自由使用的许多其他性质；见《代数》引理
\ref{algebra-lemma-ring-with-only-minimal-primes}。选择仿射开覆盖
$X = \bigcup U_i$。

\medskip\noindent
若 (4) 成立，则 $U_i \cap U_j$ 是 $U_i$ 的闭子集，因而拟紧；所以由
《概形》引理 \ref{schemes-lemma-characterize-quasi-separated}，$X$
拟分离，即 (1) 成立。

\medskip\noindent
若 (1) 成立，则 $U_i \cap U_j$ 是 $U_i$ 的拟紧开集，因而在 $U_i$
中闭。于是 $U_i \cap U_j \to U_i$ 是像为闭集的开浸入，故它是闭浸入。
特别地，$U_i \cap U_j$ 仿射，并且
$\mathcal{O}(U_i) \to \mathcal{O}_X(U_i \cap U_j)$ 满射。因此由《概形》
引理 \ref{schemes-lemma-characterize-separated}，$X$ 分离，即 (2) 成立。

\medskip\noindent
假设 (2) 成立，并令 $x, y \in X$。设 $x \in U_i$。若 $y \in U_i$，
则可以找到 $x$ 与 $y$ 的不交开邻域，因为 $U_i$ 是豪斯多夫空间。
设 $y \not \in U_i$ 且 $y \in U_j$。则 $y \not \in U_i \cap U_j$；
该交是 $U_j$ 的仿射开集，因而在 $U_j$ 中闭。因此可以找到 $y$ 的一个
不与 $U_i$ 相交的开邻域，从而断定 $X$ 是豪斯多夫空间，即 (3) 成立。

\medskip\noindent
假设 (3) 成立。令 $U \subset X$ 为仿射开集。由《拓扑》引理
\ref{topology-lemma-quasi-compact-in-Hausdorff}，$U$ 在 $X$ 中闭。
这证明 (4) 成立。

\medskip\noindent
假设 $X$ 满足等价条件 (1) -- (4)。我们证明本引理的最后几个断言。
设 $x, y \in X$ 且 $x \not = y$。由于 $y$ 不特化到 $x$，可以选择仿射开集
$U \subset X$，使得 $x \in U$ 且 $y \not \in U$。于是
$X = U \amalg (X \setminus U)$ 是开闭子集的一个分解；这说明 $x$ 与 $y$
不属于 $X$ 的同一个连通分支。其次，假设 $U \subset X$ 是拟紧开集。
将 $U = U_1 \cup \ldots \cup U_n$ 写成仿射开集之并。我们对 $n$ 归纳，
证明 $U$ 仿射。这立即将问题归结到 $n = 2$ 的情形。此时
$U = (U_1 \setminus U_2) \amalg (U_1 \cap U_2) \amalg (U_2 \setminus U_1)$，
而上述论证说明每一部分都仿射。
\end{proof}

\begin{lemma}
\label{properties-lemma-isolated-dim-0-locally-noetherian}
令 $x$ 为局部 Noether 概形 $X$ 的一点。则 $\dim_x(X) = 0$，
当且仅当 $x$ 是 $X$ 的孤立点。
\end{lemma}

\begin{proof}
若 $x$ 是孤立点，则 $\{x\}$ 是 $X$ 的开子集
（《拓扑》定义 \ref{topology-definition-isolated-point}），因而由定义
$\dim_x(X) = \dim(\{x\}) = 0$。反之，若 $\dim_x(X) = 0$，则存在
开邻域 $U \subset X$ 包含 $x$，使得 $\dim(U) = 0$
（《拓扑》定义 \ref{topology-definition-Krull}）。由引理
\ref{properties-lemma-locally-Noetherian-dimension-0}，$U$ 上的拓扑是离散的，故
$x$ 是孤立点。
\end{proof}






\section{链状概形}
\label{properties-section-catenary}

\noindent
回顾：若对拓扑空间 $X$ 的每一对不可约闭子集 $T \subset T'$，都存在
不可约闭子集的极大链
$$
T = T_0 \subset T_1 \subset \ldots \subset T_e = T'
$$
并且每条这样的链都有相同长度，则称该空间是{\it 链状的}。见《拓扑》定义
\ref{topology-definition-catenary}。

\begin{definition}
\label{properties-definition-catenary}
令 $S$ 为概形。若 $S$ 的底拓扑空间是链状的，则称 $S$ 是
{\it 链状概形}。
\end{definition}

\noindent
回顾：若对环 $A$ 的任意一对素理想
$\mathfrak p \subset \mathfrak q$，都存在素理想的极大链
$$
\mathfrak p =
\mathfrak p_0 \subset \ldots \subset \mathfrak p_e
= \mathfrak q
$$
并且所有这样的链都有相同长度，则称该环为{\it 链环}。见《代数》定义
\ref{algebra-definition-catenary}。

\begin{lemma}
\label{properties-lemma-catenary-local}
令 $S$ 为概形。下列条件等价：
\begin{enumerate}
\item $S$ 是链状概形；
\item 存在 $S$ 的开覆盖，其每个成员都是链状概形；
\item 对每个仿射开集 $\Spec(R) = U \subset S$，环 $R$ 都是链环；以及
\item 存在仿射开覆盖 $S = \bigcup U_i$，使得每个 $U_i$ 都是一个链环的谱。
\end{enumerate}
此外，在这种情形下，$S$ 的任意局部闭子概形也都是链状概形。
\end{lemma}

\begin{proof}
结合《拓扑》引理 \ref{topology-lemma-catenary} 与《代数》引理
\ref{algebra-lemma-catenary}。
\end{proof}

\begin{lemma}
\label{properties-lemma-catenary-dimension-function}
令 $S$ 为局部 Noether 概形。下列条件等价：
\begin{enumerate}
\item $S$ 是链状概形；以及
\item 在 Zariski 拓扑下局部地，$S$ 上存在维数函数
（见《拓扑》定义 \ref{topology-definition-dimension-function}）。
\end{enumerate}
\end{lemma}

\begin{proof}
这由《拓扑》引理
\ref{topology-lemma-catenary}、
\ref{topology-lemma-dimension-function-catenary} 和
\ref{topology-lemma-locally-dimension-function}，
《概形》引理 \ref{schemes-lemma-scheme-sober}，以及最后的引理
\ref{properties-lemma-Noetherian-topology} 得出。
\end{proof}

\noindent
可以证明，一个概形是链状概形，当且仅当它的各局部环都是链环。

\begin{lemma}
\label{properties-lemma-catenary-local-rings-catenary}
令 $X$ 为概形。下列条件等价：
\begin{enumerate}
\item $X$ 是链状概形；以及
\item 对任意 $x \in X$，局部环 $\mathcal{O}_{X, x}$ 都是链环。
\end{enumerate}
\end{lemma}

\begin{proof}
假设 $X$ 是链状概形。令 $x \in X$。由引理
\ref{properties-lemma-catenary-local}，可以把 $X$ 替换为 $x$ 的一个仿射开邻域；
此时 $\Gamma(X, \mathcal{O}_X)$ 是链环。由《代数》引理
\ref{algebra-lemma-localization-catenary}，链环的任意局部化仍为链环。
所以 $\mathcal{O}_{X, x}$ 是链环。

\medskip\noindent
反过来，假设 $X$ 的所有局部环都是链环。令 $Y \subset Y'$ 为 $X$
的不可约闭子集之间的包含关系，并令 $\xi \in Y$ 为泛点。令
$\mathfrak p \subset \mathcal{O}_{X, \xi}$ 为对应于 $Y'$ 的泛点的素理想；
见《概形》引理 \ref{schemes-lemma-specialize-points}。由同一引理，
$X$ 中介于 $Y$ 与 $Y'$ 之间的不可约闭子集，对应于素理想
$\mathfrak q \subset \mathcal{O}_{X, \xi}$，这些素理想满足
$\mathfrak p \subset \mathfrak q \subset \mathfrak m_{\xi}$。因为
$\mathcal{O}_{X, \xi}$ 是链环，可见这些素理想的所有极大链都是有限的，
并且具有相同长度。
\end{proof}






\section{Serre 条件}
\label{properties-section-Rk}

\noindent
下面给出两个经常有用的技术概念。另见《概形上同调》第
\ref{coherent-section-depth} 节。

\begin{definition}
\label{properties-definition-Rk}
令 $X$ 为局部 Noether 概形，并令 $k \geq 0$。
\begin{enumerate}
\item 称 $X$ {\it 在余维 $k$ 内正则}，或称 $X$ 具有性质
{\it $(R_k)$}，如果对每个 $x \in X$ 都有
$$
\dim(\mathcal{O}_{X, x}) \leq k
\Rightarrow
\mathcal{O}_{X, x}\text{ 正则}
$$
\item 称 $X$ 具有性质 {\it $(S_k)$}，如果对每个 $x \in X$ 都有
如下不等式：\par\noindent
$\text{depth}(\mathcal{O}_{X, x}) \geq \min(k, \dim(\mathcal{O}_{X, x}))$。
\end{enumerate}
\end{definition}

\noindent
“在余维 $k$ 内正则”这一说法是有意义的，因为我们已在第
\ref{properties-section-catenary} 节看到：若 $Y \subset X$ 是泛点为 $x$ 的不可约
闭子集，则 $\dim(\mathcal{O}_{X, x}) = \text{codim}(Y, X)$。例如，条件
$(R_0)$ 意味着：对每个泛点 $\eta \in X$，其中该点是 $X$ 的一个不可约分支的泛点，
局部环 $\mathcal{O}_{X, \eta}$ 都是域。但是，对一般 Noether 概形，$X$
的正则轨迹可能表现得很糟糕，所以必须谨慎。

\begin{lemma}
\label{properties-lemma-scheme-regular-iff-all-Rk}
令 $X$ 为局部 Noether 概形。则 $X$ 正则，当且仅当 $X$ 具有性质
$(R_k)$，对所有 $k \geq 0$ 都如此。
\end{lemma}

\begin{proof}
这由引理 \ref{properties-lemma-characterize-regular} 及各定义得出。
\end{proof}

\begin{lemma}
\label{properties-lemma-scheme-CM-iff-all-Sk}
令 $X$ 为局部 Noether 概形。则 $X$ 是 Cohen-Macaulay 概形，当且仅当
$X$ 具有性质 $(S_k)$，对所有 $k \geq 0$ 都如此。
\end{lemma}

\begin{proof}
由引理 \ref{properties-lemma-characterize-Cohen-Macaulay}，只须考察局部环。
因此，本引理成立，因为 Noether 局部环是 Cohen-Macaulay 环，当且仅当
其深度等于其维数。
\end{proof}

\begin{lemma}
\label{properties-lemma-criterion-reduced}
令 $X$ 为局部 Noether 概形。则 $X$ 约化，当且仅当 $X$ 具有性质
$(S_1)$ 和 $(R_0)$。
\end{lemma}

\begin{proof}
这就是《代数》引理 \ref{algebra-lemma-criterion-reduced}。
\end{proof}

\begin{lemma}
\label{properties-lemma-criterion-normal}
令 $X$ 为局部 Noether 概形。则 $X$ 正规，当且仅当 $X$ 具有性质
$(S_2)$ 和 $(R_1)$。
\end{lemma}

\begin{proof}
这就是《代数》引理 \ref{algebra-lemma-criterion-normal}。
\end{proof}

\begin{lemma}
\label{properties-lemma-normal-dimension-1-regular}
令 $X$ 为正规局部 Noether 概形，且维数 $\leq 1$。则 $X$ 正则。
\end{lemma}

\begin{proof}
这由引理 \ref{properties-lemma-criterion-normal} 及各定义得出。
\end{proof}

\begin{lemma}
\label{properties-lemma-normal-dimension-2-Cohen-Macaulay}
令 $X$ 为正规局部 Noether 概形，且维数 $\leq 2$。则 $X$ 是
Cohen-Macaulay 概形。
\end{lemma}

\begin{proof}
这由引理 \ref{properties-lemma-criterion-normal} 及各定义得出。
\end{proof}








\section{Japanese 概形与 Nagata 概形}
\label{properties-section-nagata}

\noindent
本节所考虑的概念在 EGA 中并没有得到醒目的定义。EGA 的
\cite[IV Corollary 5.11.4]{EGA} 提到并定义了“泛 Japanese 概形”；其
\cite[IV Remark 10.4.14 (ii)]{EGA} 提到“Japanese 概形”，但没有给出定义。
文献中有少数地方使用下面定义的 Nagata 概形
（例如见 \cite[Definition 8.2.30]{Liu} 和 \cite[Page 142]{Greco}）。

\medskip\noindent
简要回顾：若整环 $R$ 在其分式域的任意有限扩张中的整闭包都是 $R$ 上的
有限模，则称 $R$ 是{\it Japanese 整环}。若环 $R$ 对任意有限型环同态
$R \to S$，只要 $S$ 是整环，$S$ 就是 Japanese 整环，则称该环是
{\it 泛 Japanese 环}。若环 $R$ 是 Noether 环，并且
$R/\mathfrak p$ 对每个素理想 $\mathfrak p$ 都是 Japanese 整环，则称 $R$ 为
{\it Nagata 环}。

\begin{definition}
\label{properties-definition-nagata}
令 $X$ 为概形。
\begin{enumerate}
\item 假设 $X$ 是整概形。称 $X$ 是{\it Japanese 概形}，如果对每个
$x \in X$，都存在仿射开邻域 $x \in U \subset X$，使得环
$\mathcal{O}_X(U)$ 是 Japanese 整环
（见《代数》定义 \ref{algebra-definition-N}）。
\item 称 $X$ 是{\it 泛 Japanese 概形}，如果对每个 $x \in X$，都存在
仿射开邻域 $x \in U \subset X$，使得环 $\mathcal{O}_X(U)$ 是泛
Japanese 环（见《代数》定义 \ref{algebra-definition-nagata}）。
\item 称 $X$ 是{\it Nagata 概形}，如果对每个 $x \in X$，都存在仿射
开邻域 $x \in U \subset X$，使得环 $\mathcal{O}_X(U)$ 是 Nagata 环
（见《代数》定义 \ref{algebra-definition-nagata}）。
\end{enumerate}
\end{definition}

\noindent
成为 Nagata 概形等价于局部 Noether 且为泛 Japanese 概形；见引理
\ref{properties-lemma-nagata-universally-Japanese}。

\begin{remark}
\label{properties-remark-non-integral-Japanese}
在 \cite{Hoobler-finite} 中，称一个（局部 Noether）概形 $X$ 是
Japanese 概形，如果对每个 $x \in X$ 以及每个伴随素理想
$\mathfrak p$（它是 $\mathcal{O}_{X, x}$ 的伴随素理想），环
$\mathcal{O}_{X, x}/\mathfrak p$ 都是
Japanese 整环。我们不使用这个定义，因为存在一维 Noether 整环，其局部环
都是优秀环（特别地，都是 Japanese 整环），但它的正规化不是有限的。见
\cite[Example 1]{Hochster-loci}、\cite{Heinzer-Levy} 或
\cite[Expos\'e XIX]{Traveaux}。另一方面，可以用如下方式避开这个问题：
称概形 $X$ 是 Japanese 概形，如果对每个仿射开集
$\Spec(A) \subset X$，环 $A/\mathfrak p$ 对每个伴随素理想
$\mathfrak p$（属于 $A$）都是 Japanese 整环。
\end{remark}

\begin{lemma}
\label{properties-lemma-nagata-locally-Noetherian}
Nagata 概形是局部 Noether 的。
\end{lemma}

\begin{proof}
这是因为 Nagata 环按定义是 Noether 环。
\end{proof}

\begin{lemma}
\label{properties-lemma-locally-Japanese}
令 $X$ 为整概形。下列条件等价：
\begin{enumerate}
\item 概形 $X$ 是 Japanese 概形。
\item 对每个仿射开集 $U \subset X$，整环 $\mathcal{O}_X(U)$ 都是
Japanese 整环。
\item 存在仿射开覆盖 $X = \bigcup U_i$，使得每个
$\mathcal{O}_X(U_i)$ 都是 Japanese 整环。
\item 存在开覆盖 $X = \bigcup X_j$，使得每个开子概形 $X_j$ 都是
Japanese 概形。
\end{enumerate}
此外，若 $X$ 是 Japanese 概形，则每个开子概形也是 Japanese 概形。
\end{lemma}

\begin{proof}
这由引理 \ref{properties-lemma-locally-P} 以及《代数》引理
\ref{algebra-lemma-localize-N} 和
\ref{algebra-lemma-Japanese-local} 得出。
\end{proof}

\begin{lemma}
\label{properties-lemma-locally-universally-Japanese}
令 $X$ 为概形。下列条件等价：
\begin{enumerate}
\item 概形 $X$ 是泛 Japanese 概形。
\item 对每个仿射开集 $U \subset X$，环 $\mathcal{O}_X(U)$ 都是
泛 Japanese 环。
\item 存在仿射开覆盖 $X = \bigcup U_i$，使得每个
$\mathcal{O}_X(U_i)$ 都是泛 Japanese 环。
\item 存在开覆盖 $X = \bigcup X_j$，使得每个开子概形 $X_j$ 都是
泛 Japanese 概形。
\end{enumerate}
此外，若 $X$ 是泛 Japanese 概形，则每个开子概形也是泛 Japanese 概形。
\end{lemma}

\begin{proof}
这由引理 \ref{properties-lemma-locally-P} 以及《代数》引理
\ref{algebra-lemma-universally-japanese} 和
\ref{algebra-lemma-nagata-local} 得出。
\end{proof}

\begin{lemma}
\label{properties-lemma-locally-nagata}
令 $X$ 为概形。下列条件等价：
\begin{enumerate}
\item 概形 $X$ 是 Nagata 概形。
\item 对每个仿射开集 $U \subset X$，环 $\mathcal{O}_X(U)$ 都是 Nagata 环。
\item 存在仿射开覆盖 $X = \bigcup U_i$，使得每个
$\mathcal{O}_X(U_i)$ 都是 Nagata 环。
\item 存在开覆盖 $X = \bigcup X_j$，使得每个开子概形 $X_j$ 都是
Nagata 概形。
\end{enumerate}
此外，若 $X$ 是 Nagata 概形，则每个开子概形也是 Nagata 概形。
\end{lemma}

\begin{proof}
这由引理 \ref{properties-lemma-locally-P} 以及《代数》引理
\ref{algebra-lemma-nagata-localize} 和
\ref{algebra-lemma-nagata-local} 得出。
\end{proof}

\begin{lemma}
\label{properties-lemma-characterize-nagata}
令 $X$ 为局部 Noether 概形。则 $X$ 是 Nagata 概形，当且仅当每个整的
闭子概形 $Z \subset X$ 都是 Japanese 概形。
\end{lemma}

\begin{proof}
假设 $X$ 是 Nagata 概形。令 $Z \subset X$ 为整的闭子概形。令 $z \in Z$。
取仿射开集 $\Spec(A) = U \subset X$ 包含 $z$，使得 $A$ 是 Nagata 环。
则有 $Z \cap U \cong \Spec(A/\mathfrak p)$，其中 $\mathfrak p$ 为某个素理想；
见《概形》引理
\ref{schemes-lemma-closed-subspace-scheme}
（以及定义 \ref{properties-definition-integral}）。由《代数》定义
\ref{algebra-definition-nagata} 可知 $A/\mathfrak p$ 是 Japanese 整环。
因此，按定义，$Z$ 是 Japanese 概形。

\medskip\noindent
假设 $X$ 的每个整的闭子概形都是 Japanese 概形。任取仿射开集
$\Spec(A) = U \subset X$。因为 $X$ 局部 Noether，由引理
\ref{properties-lemma-locally-Noetherian} 可知 $A$ 是 Noether 环。令
$\mathfrak p \subset A$ 为素理想。我们须证明 $A/\mathfrak p$ 是
Japanese 整环。令 $T \subset U$ 为闭子集
$V(\mathfrak p) \subset \Spec(A)$，并令 $\overline{T} \subset X$ 为其
闭包。由于 $\overline{T}$ 是不可约子集的闭包，故它不可约。因此，由
$\overline{T}$ 定义的约化闭子概形是整的闭子概形（仍称为
$\overline{T}$）；见《概形》引理
\ref{schemes-lemma-reduced-closed-subscheme}。换言之，
$\Spec(A/\mathfrak p)$ 是 $X$ 的一个整的闭子概形的仿射开集。依假设，
这个子概形是 Japanese 概形；再由引理 \ref{properties-lemma-locally-Japanese}，
可知 $A/\mathfrak p$ 是 Japanese 整环。
\end{proof}

\begin{lemma}
\label{properties-lemma-nagata-universally-Japanese}
令 $X$ 为概形。下列条件等价：
\begin{enumerate}
\item $X$ 是 Nagata 概形；以及
\item $X$ 局部 Noether 且为泛 Japanese 概形。
\end{enumerate}
\end{lemma}

\begin{proof}
这就是《代数》命题
\ref{algebra-proposition-nagata-universally-japanese}。
\end{proof}

\noindent
这一讨论将在《态射》第 \ref{morphisms-section-nagata} 节继续。




\section{G-概形}
\label{properties-section-G-scheme}

\noindent
我们指出：环 $R$ 是 G-环，当且仅当 $R$ 是 Noether 环，并且对每个素理想
$\mathfrak p$（属于 $R$），环同态
$R_\mathfrak p \to (R_\mathfrak p)^\wedge$ 都是正则的。

\begin{definition}
\label{properties-definition-G-scheme}
令 $X$ 为概形。称 $X$ 是{\it G-概形}\footnote{这可能不是标准术语。若
$G$ 是有限群或群概形，则 $G$-概形有时指带有 $G$ 作用的概形。}，如果对每个
$x \in X$，都存在仿射开邻域 $x \in U \subset X$，使得环
$\mathcal{O}_X(U)$ 是 G-环（见《代数续篇》定义
\ref{more-algebra-definition-G-ring}）。
\end{definition}

\begin{lemma}
\label{properties-lemma-characterize-G-scheme}
令 $X$ 为概形。下列条件等价：
\begin{enumerate}
\item $X$ 是 G-概形；
\item $X$ 局部 Noether，并且对所有 $x \in X$，环同态
$\mathcal{O}_{X, x} \to \mathcal{O}_{X, x}^\wedge$ 都正则；以及
\item $X$ 局部 Noether，并且对任意闭点 $x \in X$，环同态
$\mathcal{O}_{X, x} \to \mathcal{O}_{X, x}^\wedge$ 都正则。
\end{enumerate}
\end{lemma}

\begin{proof}
显然，(1) 蕴含 (2)，而 (2) 蕴含 (3)。假设 (3)。任取 $x \in X$。
由引理 \ref{properties-lemma-locally-Noetherian-closed-point}，存在特化
$x \leadsto x'$，其中 $x'$ 在 $X$ 中闭。于是，由《代数续篇》引理
\ref{more-algebra-lemma-check-G-ring-maximal-ideals}，
$\mathcal{O}_{X, x'}$ 是 G-环。由于 $\mathcal{O}_{X, x}$ 是
$\mathcal{O}_{X, x'}$ 在一个素理想处的局部化
（《概形》引理 \ref{schemes-lemma-specialize-points}），可知
$\mathcal{O}_{X, x} \to \mathcal{O}_{X, x}^\wedge$ 是正则环同态。
由于 $x \in X$ 任意，可得：对任意仿射开集 $U \subset X$，Noether 环
（引理 \ref{properties-lemma-locally-Noetherian}）$\mathcal{O}_X(U)$ 都是 G-环。
这证明 (1) 成立。
\end{proof}

\begin{lemma}
\label{properties-lemma-locally-G-scheme}
令 $X$ 为概形。下列条件等价：
\begin{enumerate}
\item 概形 $X$ 是 G-概形。
\item 对每个仿射开集 $U \subset X$，环 $\mathcal{O}_X(U)$ 都是 G-环。
\item 存在仿射开覆盖 $X = \bigcup U_i$，使得每个
$\mathcal{O}_X(U_i)$ 都是 G-环。
\item 存在开覆盖 $X = \bigcup X_j$，使得每个开子概形 $X_j$ 都是
G-概形。
\end{enumerate}
此外，若 $X$ 是 G-概形，则每个开子概形也都是 G-概形。
\end{lemma}

\begin{proof}
结合引理 \ref{properties-lemma-locally-Noetherian}
和 \ref{properties-lemma-characterize-G-scheme}。
\end{proof}



\section{奇异轨迹}
\label{properties-section-singular-locus}

\noindent
定义如下。

\begin{definition}
\label{properties-definition-singular-locus}
令 $X$ 为局部 Noether 概形。{\it 正则轨迹} $\text{Reg}(X)$，即 $X$
的所有 $x \in X$ 中满足 $\mathcal{O}_{X, x}$ 为正则局部环的点所成的集合。
{\it 奇异轨迹} $\text{Sing}(X)$ 是补集
$X \setminus \text{Reg}(X)$，即所有 $x \in X$ 中满足
$\mathcal{O}_{X, x}$ 不是正则局部环的点所成的集合。
\end{definition}

\noindent
局部 Noether 概形的正则轨迹在泛化下稳定；见《代数》定义
\ref{algebra-definition-regular} 之前的讨论。然而，对一般局部 Noether
概形，正则轨迹未必是开的。读者可以在《代数续篇》第
\ref{more-algebra-section-singular-locus} 节找到它成为开集的一些判据。
我们将在《态射》第 \ref{morphisms-section-singular-locus} 节进一步讨论。






\section{局部不可约性}
\label{properties-section-unibranch}

\noindent
回顾：我们已在《代数续篇》第 \ref{more-algebra-section-unibranch} 节
引入了（几何）单枝局部环的概念。

\begin{definition}
\label{properties-definition-unibranch}
\begin{reference}
\cite[Chapter IV (6.15.1)]{EGA4}
\end{reference}
令 $X$ 为概形，并令 $x \in X$。称 $X$ {\it 在 $x$ 处单枝}，如果局部环
$\mathcal{O}_{X, x}$ 单枝。称 $X$ {\it 在 $x$ 处几何单枝}，如果局部环
$\mathcal{O}_{X, x}$ 几何单枝。称 $X$ {\it 单枝}，如果 $X$ 在其所有点处
都单枝。称 $X$ {\it 几何单枝}，如果 $X$ 在其所有点处都几何单枝。
\end{definition}

\noindent
确实可能发生如下情形：局部环 $A$ 在《代数续篇》定义
\ref{more-algebra-definition-unibranch} 的意义下几何单枝，但概形
$\Spec(A)$ 在定义 \ref{properties-definition-unibranch} 的意义下并非几何单枝。
例如，当 $A$ 是一条具有普通二重点奇点（结点）的不可约平面曲线之锥
在顶点处的局部环时，就会发生这种情形。

\begin{lemma}
\label{properties-lemma-normal-geometrically-unibranch}
正规概形是几何单枝的。
\end{lemma}

\begin{proof}
这由各定义得出。具体地说，概形正规是指其各局部环都是正规整环；
由《代数续篇》定义 \ref{more-algebra-definition-unibranch} 立即可知，
正规局部整环几何单枝。
\end{proof}

\begin{lemma}
\label{properties-lemma-geometrically-unibranch}
\begin{reference}
与 \cite[Proposition 2.3]{Etale-coverings} 比较
\end{reference}
令 $X$ 为 Noether 概形。下列条件等价：
\begin{enumerate}
\item $X$ 几何单枝（定义 \ref{properties-definition-unibranch}）；
\item 对每个点 $x \in X$，若它不是 $X$ 的某个不可约分支之泛点，则
严格 Hensel 化 $\mathcal{O}_{X, x}^{sh}$ 的穿孔谱都连通。
\end{enumerate}
\end{lemma}

\begin{proof}
《代数续篇》引理 \ref{more-algebra-lemma-geometrically-unibranch}
说明，(1) 蕴含 (2) 中的穿孔谱不可约，因而特别地连通。

\medskip\noindent
假设 (2)。令 $x \in X$。我们须证明 $\mathcal{O}_{X, x}$ 几何单枝。
对 $\dim(\mathcal{O}_{X, x})$ 归纳，可以假设结论对 $x$ 的每个非平凡
泛化都成立。可以把 $X$ 替换为 $\Spec(\mathcal{O}_{X, x})$。换言之，
可以假设 $X = \Spec(A)$，其中 $A$ 是局部环，并且 $A_\mathfrak p$
对每个非极大素理想 $\mathfrak p \subset A$ 都几何单枝。

\medskip\noindent
令 $A^{sh}$ 为 $A$ 的严格 Hensel 化。若
$\mathfrak q \subset A^{sh}$ 是位于 $\mathfrak p \subset A$ 上方的素理想，
则 $A_\mathfrak p \to A^{sh}_\mathfrak q$ 是 étale 代数的滤过余极限。
因此，$A_\mathfrak p$ 与 $A^{sh}_\mathfrak q$ 的严格 Hensel 化同构。
由此及《代数续篇》引理
\ref{more-algebra-lemma-geometrically-unibranch} 可知：
$A^{sh}_\mathfrak q$ 对每个非极大素理想 $\mathfrak q$（属于 $A^{sh}$）
都有唯一极小素理想。

\medskip\noindent
令 $\mathfrak q_1, \ldots, \mathfrak q_r$ 为 $A^{sh}$ 的极小素理想。
我们须证明 $r = 1$。由上文可知，
$V(\mathfrak q_1) \cap V(\mathfrak q_j) = \{\mathfrak m^{sh}\}$ 对
$j = 2, \ldots, r$ 成立。
因此，$V(\mathfrak q_1) \setminus \{\mathfrak m^{sh}\}$ 是 $A^{sh}$ 的
穿孔谱的开闭子集；这与该穿孔谱连通的假设矛盾，除非 $r = 1$。
\end{proof}

\begin{definition}
\label{properties-definition-number-of-branches}
令 $X$ 为概形，并令 $x \in X$。{\it $X$ 在 $x$ 处的分枝数}，是局部环
$\mathcal{O}_{X, x}$ 的分枝数，其定义见《代数续篇》定义
\ref{more-algebra-definition-number-of-branches}。{\it $X$ 在 $x$ 处的
几何分枝数}，是局部环 $\mathcal{O}_{X, x}$ 的几何分枝数，其定义见
《代数续篇》定义 \ref{more-algebra-definition-number-of-branches}。
\end{definition}

\noindent
我们常想把它与完备局部环的分枝作比较，但一般而言，这种比较并不直接；
关于这个问题的一些信息见《代数续篇》第
\ref{more-algebra-section-branches-completion} 节。

\begin{lemma}
\label{properties-lemma-number-of-branches-irreducible-components}
令 $X$ 为概形且 $x \in X$。令 $X_i$（$i \in I$）为 $X$ 的所有经过
$x$ 的不可约分支。则 $X$ 在 $x$ 处的（几何）分枝数，等于对
$i \in I$ 求和所得的 $X_i$ 在 $x$ 处的（几何）分枝数之和。
\end{lemma}

\begin{proof}
把 $X_i$ 看作 $X$ 的整的闭子概形；见《概形》定义
\ref{schemes-definition-reduced-induced-scheme} 和引理
\ref{properties-lemma-characterize-integral}。注意，$X_i$ 在 $x$ 处
的（几何）分枝数至少为 $1$，对所有 $i$ 都如此（本质上由定义）。回顾：
$X_i$ 按 $1$ 对 $1$ 的方式与极小素理想
$\mathfrak p_i \subset \mathcal{O}_{X, x}$ 对应；见《代数》引理
\ref{algebra-lemma-irreducible-components-containing-x}。因此，若 $I$
无限，则 $\mathcal{O}_{X, x}$ 有无限多个极小素理想，从而
$\mathcal{O}_{X, x}^h$ 与 $\mathcal{O}_{X, x}^{sh}$ 都有无限多个极小
素理想（结合《代数》引理
\ref{algebra-lemma-injective-minimal-primes-in-image}、
\ref{algebra-lemma-minimal-prime-image-minimal-prime} 以及同态
$\mathcal{O}_{X, x} \to  \mathcal{O}_{X, x}^h \to \mathcal{O}_{X, x}^{sh}$
的单射性）。在这种情形下，$X$ 在 $x$ 处的（几何）分枝数按定义为
$\infty$，求和所得结果也是如此。因此可假设 $I$ 有限。
令 $A'$ 为 $\mathcal{O}_{X, x}$ 在总分式环 $Q$ 中的整闭包，而该环
是 $(\mathcal{O}_{X, x})_{red}$ 的总分式环。令 $A'_i$ 为
$\mathcal{O}_{X, x}/\mathfrak p_i$ 在总分式环 $Q_i$ 中的整闭包，
而该环是 $\mathcal{O}_{X, x}/\mathfrak p_i$ 的总分式环。由《代数》引理
\ref{algebra-lemma-total-ring-fractions-no-embedded-points}，
$Q = \prod_{i \in I} Q_i$。因此 $A' = \prod A'_i$。本引理的等式随即由
《代数续篇》引理 \ref{more-algebra-lemma-number-of-branches-1} 得出；
该引理用 $A'$ 的极大理想表出（几何）分枝数。
\end{proof}

\begin{lemma}
\label{properties-lemma-number-of-branches-1}
令 $X$ 为概形，并令 $x \in X$。
\begin{enumerate}
\item $X$ 在 $x$ 处的分枝数为 $1$，当且仅当 $X$ 在 $x$ 处单枝。
\item $X$ 在 $x$ 处的几何分枝数为 $1$，当且仅当 $X$ 在 $x$ 处几何单枝。
\end{enumerate}
\end{lemma}

\begin{proof}
本引理由各定义以及环的相应结论立即得出；见《代数续篇》引理
\ref{more-algebra-lemma-number-of-branches-1}。
\end{proof}




\section{有限型与有限表现模的刻画}
\label{properties-section-characterizing-finite-type-presentation}

\noindent
令 $X$ 为概形。
令 $\mathcal{F}$ 为准凝聚 $\mathcal{O}_X$-模。
下述引理表明，$\mathcal{F}$ 是有限型的
（见《模》定义 \ref{modules-definition-finite-type}）
当且仅当 $\mathcal{F}$ 在每个仿射开子集 $\Spec(A) = U \subset X$ 上
都具有形式 $\widetilde M$，即来自某个有限 $A$-模 $M$。
类似地，$\mathcal{F}$ 是有限表现的
（见《模》定义 \ref{modules-definition-finite-presentation}）
当且仅当 $\mathcal{F}$ 在每个仿射开子集 $\Spec(A) = U \subset X$ 上
都具有形式 $\widetilde M$，其中它来自某个有限表现 $A$-模 $M$。

\begin{lemma}
\label{properties-lemma-finite-type-module}
令 $X = \Spec(R)$ 为仿射概形。
准凝聚 $\mathcal{O}_X$-模层
$\widetilde M$ 是有限型 $\mathcal{O}_X$-模，当且仅当
$M$ 是有限 $R$-模。
\end{lemma}

\begin{proof}
假设 $\widetilde M$ 是有限型 $\mathcal{O}_X$-模。
这意味着存在 $X$ 的一个开覆盖，使得 $\widetilde M$ 在该覆盖各成员
上的限制由有限多个截面整体生成。因此还存在一个标准开覆盖
$X = \bigcup_{i = 1, \ldots, n} D(f_i)$，使得 $\widetilde M|_{D(f_i)}$
由有限多个截面生成。于是 $M_{f_i}$ 对每个 $i$ 都有限生成。
故结论由《代数》引理 \ref{algebra-lemma-cover} 得出。
\end{proof}

\begin{lemma}
\label{properties-lemma-finite-presentation-module}
令 $X = \Spec(R)$ 为仿射概形。准凝聚 $\mathcal{O}_X$-模层
$\widetilde M$ 是有限表现 $\mathcal{O}_X$-模，当且仅当
$M$ 是有限表现 $R$-模。
\end{lemma}

\begin{proof}
假设 $\widetilde M$ 是有限表现 $\mathcal{O}_X$-模。
由引理 \ref{properties-lemma-finite-type-module} 可知 $M$ 是有限 $R$-模。
选取一个满射 $R^n \to M$，记其核为 $K$。由《概形》引理
\ref{schemes-lemma-spec-sheaves}，存在短正合列
$$
0 \to \widetilde{K} \to
\bigoplus \mathcal{O}_X^{\oplus n} \to
\widetilde{M} \to 0
$$
由《模》引理
\ref{modules-lemma-kernel-surjection-finite-free-onto-finite-presentation}
可知 $\widetilde{K}$ 是有限型 $\mathcal{O}_X$-模。
再次应用引理 \ref{properties-lemma-finite-type-module}，
可知 $K$ 是有限 $R$-模。因此 $M$ 是有限表现 $R$-模。
\end{proof}

\section{主开集上的截面}
\label{properties-section-principal-opens}

\noindent
下面是这类结果的一个典型例子。在后面的引理中，我们将采用一种更朴素但
也更直接的证明方法。

\begin{lemma}
\label{properties-lemma-invert-f-sections}
\begin{slogan}
准凝聚层的截面在无穷远处只有亚纯奇点。
\end{slogan}
令 $X$ 为概形。令 $f \in \Gamma(X, \mathcal{O}_X)$。
把 $X_f \subset X$ 记为 $f$ 可逆之处的开集，见《概形》引理
\ref{schemes-lemma-f-open}。
若 $X$ 拟紧且拟分离，则典范映射
$$
\Gamma(X, \mathcal{O}_X)_f \longrightarrow \Gamma(X_f, \mathcal{O}_X)
$$
是同构。此外，若 $\mathcal{F}$ 是准凝聚 $\mathcal{O}_X$-模层，则映射
$$
\Gamma(X, \mathcal{F})_f \longrightarrow \Gamma(X_f, \mathcal{F})
$$
是同构。
\end{lemma}

\begin{proof}
记 $R = \Gamma(X, \mathcal{O}_X)$。考虑概形的典范态射
$$
\varphi : X \longrightarrow \Spec(R)
$$
，见《概形》引理
\ref{schemes-lemma-morphism-into-affine}。
右端标准开集 $D(f)$ 的逆像就是左端的 $X_f$。
此外，因为假设 $X$ 拟紧且拟分离，态射 $\varphi$ 也拟紧且拟分离；
见《概形》引理 \ref{schemes-lemma-quasi-compact-affine} 和
\ref{schemes-lemma-compose-after-separated}。因此由《概形》引理
\ref{schemes-lemma-push-forward-quasi-coherent} 可知
$\varphi_*\mathcal{F}$ 是准凝聚的。
于是 $\varphi_*\mathcal{F} = \widetilde M$，其中
$M = \Gamma(X, \mathcal{F})$ 是一个 $R$-模。因此
$$
\Gamma(X_f, \mathcal{F}) =
\Gamma(D(f), \varphi_*\mathcal{F}) =
\Gamma(D(f), \widetilde M) = M_f
$$
，这恰是本引理的内容。令 $\mathcal{F} = \mathcal{O}_X$，
即可得到本引理中第一个陈列同构。
\end{proof}

\noindent
回顾：给定概形 $X$、可逆层 $\mathcal{L}$（在 $X$ 上），以及
$\mathcal{O}_X$-模层 $\mathcal{F}$，可得到分次环
$\Gamma_*(X, \mathcal{L}) =
\bigoplus\nolimits_{n \geq 0} \Gamma(X, \mathcal{L}^{\otimes n})$
以及分次 $\Gamma_*(X, \mathcal{L})$-模
$\Gamma_*(X, \mathcal{L}, \mathcal{F}) =
\bigoplus\nolimits_{n \in \mathbf{Z}}
\Gamma(X, \mathcal{F} \otimes_{\mathcal{O}_X} \mathcal{L}^{\otimes n})$；
见《模》定义 \ref{modules-definition-gamma-star}。
若还给定一个截面 $s \in \Gamma(X, \mathcal{L})$，则得到映射
\begin{equation}
\label{properties-equation-module-invert-s}
\Gamma_*(X, \mathcal{L}, \mathcal{F})_{(s)}
\longrightarrow
\Gamma(X_s, \mathcal{F}|_{X_s})
\end{equation}
，它把 $t/s^n$（其中
$t \in \Gamma(X, \mathcal{F} \otimes_{\mathcal{O}_X} \mathcal{L}^{\otimes n})$）
映到 $t|_{X_s} \otimes s|_{X_s}^{-n}$。这是有意义的，因为按定义
$X_s \subset X$ 是 $s$ 有逆的开集；见《模》引理
\ref{modules-lemma-s-open}。

\begin{lemma}
\label{properties-lemma-invert-s-sections}
\begin{reference}
\cite[Section 6.8]{EGA1-second}
\end{reference}
令 $X$ 为概形。令 $\mathcal{L}$ 为 $X$ 上的可逆层。
令 $s \in \Gamma(X, \mathcal{L})$。令 $\mathcal{F}$ 为准凝聚
$\mathcal{O}_X$-模。
\begin{enumerate}
\item 若 $X$ 拟紧，则 (\ref{properties-equation-module-invert-s}) 是单射；
\item 若 $X$ 拟紧且拟分离，则
(\ref{properties-equation-module-invert-s}) 是同构。
\end{enumerate}
特别地，若 $X$ 拟紧且拟分离，则典范映射
$$
\Gamma_*(X, \mathcal{L})_{(s)}
\longrightarrow
\Gamma(X_s, \mathcal{O}_X),\quad
a/s^n \longmapsto a \otimes s^{-n}
$$
是同构。
\end{lemma}

\begin{proof}
假设 $X$ 拟紧。选取一个有限仿射开覆盖
$X = U_1 \cup \ldots \cup U_m$，其中 $U_j$ 仿射且
$\mathcal{L}|_{U_j} \cong \mathcal{O}_{U_j}$。经此同构，
像 $s|_{U_j}$ 对应于某个
$f_j \in \Gamma(U_j, \mathcal{O}_{U_j})$。于是
$X_s \cap U_j = D(f_j)$。

\medskip\noindent
(1) 的证明。令 $t/s^n$ 为
(\ref{properties-equation-module-invert-s}) 的核中元素。则 $t|_{X_s} = 0$。
故 $(t|_{U_j})|_{D(f_j)} = 0$。由引理
\ref{properties-lemma-invert-f-sections} 可知
$f_j^{e_j} t|_{U_j} = 0$，其中某个 $e_j \geq 0$。
令 $e = \max(e_j)$。于是 $t \otimes s^e$ 在 $U_j$ 上的限制对所有
$j$ 均为零，
故它为零。由于 $t/s^n$ 等于
$t \otimes s^e/s^{n + e}$（作为
$\Gamma_*(X, \mathcal{L}, \mathcal{F})_{(s)}$ 中的元素），
故如所需有 $t/s^n = 0$。

\medskip\noindent
(2) 的证明。假设 $X$ 拟紧且拟分离。
则 $U_j \cap U_{j'}$ 对所有指标对 $j, j'$ 均拟紧；见《概形》引理
\ref{schemes-lemma-characterize-quasi-separated}。
由 (1) 可知 (\ref{properties-equation-module-invert-s}) 是单射。
令 $t' \in \Gamma(X_s, \mathcal{F}|_{X_s})$。对每个 $j$，存在整数
$e_j \geq 0$ 以及
$t'_j \in \Gamma(U_j, \mathcal{F}|_{U_j})$，使得
$t'|_{D(f_j)}$ 经由引理 \ref{properties-lemma-invert-f-sections} 的同构对应于
$t'_j/f_j^{e_j}$。
令 $e = \max(e_j)$，并置
$$
t_j = f_j^{e - e_j} t'_j \otimes q_j^e \in
\Gamma(U_j,
(\mathcal{F} \otimes_{\mathcal{O}_X} \mathcal{L}^{\otimes e})|_{U_j})
$$
，其中 $q_j \in \Gamma(U_j, \mathcal{L}|_{U_j})$ 是来自同构
$\mathcal{L}|_{U_j} \cong \mathcal{O}_{U_j}$ 的平凡化截面。特别地，
$s|_{U_j} = f_j q_j$。利用这一点计算可见，
$t_j|_{U_j \cap U_{j'}}$ 与 $t_{j'}|_{U_j \cap U_{j'}}$
映到 $\mathcal{F}$ 在 $U_j \cap U_{j'} \cap X_s$ 上的同一截面。
由 $U_j \cap U_{j'}$ 的拟紧性及 (1)，存在整数
$e' \geq 0$，使得
$$
t_j|_{U_j \cap U_{j'}} \otimes s^{e'}|_{U_j \cap U_{j'}} =
t_{j'}|_{U_j \cap U_{j'}} \otimes s^{e'}|_{U_j \cap U_{j'}}
$$
作为 $\mathcal{F} \otimes \mathcal{L}^{\otimes e + e'}$ 在
$U_j \cap U_{j'}$ 上的截面相等。
可选择同一个 $e'$ 对所有指标对 $j, j'$ 都适用。于是层条件给出一个截面
$t \in \Gamma(X, \mathcal{F} \otimes \mathcal{L}^{\otimes e + e'})$，
其在 $U_j$ 上的限制为 $t_j \otimes s^{e'}|_{U_j}$。
简单计算表明，$t/s^{e + e'}$ 映到 $t'$，如所需。
\end{proof}

\noindent
令 $X$ 为概形。令 $\mathcal{L}$ 为可逆 $\mathcal{O}_X$-模。
令 $\mathcal{F}$ 与 $\mathcal{G}$ 为准凝聚 $\mathcal{O}_X$-模。
考虑分次 $\Gamma_*(X, \mathcal{L})$-模
$$
M = \bigoplus\nolimits_{n \in \mathbf{Z}} \Hom_{\mathcal{O}_X}(\mathcal{F},
\mathcal{G} \otimes_{\mathcal{O}_X} \mathcal{L}^{\otimes n})
$$
接着，令 $s \in \Gamma(X, \mathcal{L})$ 为截面。于是有典范映射
\begin{equation}
\label{properties-equation-hom-invert-s}
M_{(s)} \longrightarrow
\Hom_{\mathcal{O}_{X_s}}(\mathcal{F}|_{X_s}, \mathcal{G}|_{X_s})
\end{equation}
，它把 $\alpha/s^n$ 映到映射
$\alpha|_{X_s} \otimes s|_{X_s}^{-n}$。
下述引理与引理 \ref{properties-lemma-extend-finite-presentation} 合在一起，大致说明：
若 $X$ 拟紧且拟分离，则有限表现 $\mathcal{O}_{X_s}$-模范畴就是
有限表现 $\mathcal{O}_X$-模范畴把映射的乘法系
$s^n: \mathcal{F} \to
\mathcal{F} \otimes_{\mathcal{O}_X} \mathcal{L}^{\otimes n}$ 中的映射都
变为可逆之后所得的范畴。

\begin{lemma}
\label{properties-lemma-section-maps-backwards}
令 $X$ 为概形。令 $\mathcal{L}$ 为可逆 $\mathcal{O}_X$-模。
令 $s \in \Gamma(X, \mathcal{L})$ 为截面。
令 $\mathcal{F}$、$\mathcal{G}$ 为准凝聚 $\mathcal{O}_X$-模。
\begin{enumerate}
\item 若 $X$ 拟紧且 $\mathcal{F}$ 是有限型的，则
(\ref{properties-equation-hom-invert-s}) 是单射；
\item 若 $X$ 拟紧且拟分离，并且 $\mathcal{F}$ 是有限表现的，则
(\ref{properties-equation-hom-invert-s}) 是双射。
\end{enumerate}
\end{lemma}

\begin{proof}
先证明 $X = \Spec(A)$ 为仿射且
$\mathcal{L} = \mathcal{O}_X$ 的情形。此时 $s$ 对应于某个
$f \in A$。设
$\mathcal{F} = \widetilde{M}$ 且 $\mathcal{G} = \widetilde{N}$，
其中它们是某些 $A$-模 $M$ 和 $N$。于是本引理
（通过引理 \ref{properties-lemma-finite-type-module} 与
\ref{properties-lemma-finite-presentation-module}）转化为下列代数陈述：
\begin{enumerate}
\item 若 $M$ 是有限 $A$-模，且 $\varphi : M \to N$ 是
$A$-模映射，使得诱导映射 $M_f \to N_f$ 为零，
则 $f^n\varphi = 0$ 对某个 $n$ 成立。
\item 若 $M$ 是有限表现 $A$-模，则
$\Hom_A(M, N)_f = \Hom_{A_f}(M_f, N_f)$。
\end{enumerate}
第二个陈述是《代数》引理
\ref{algebra-lemma-hom-from-finitely-presented}；第一个陈述的证明从略。

\medskip\noindent
接着证明一般 $X$ 的 (1)。
假设 $X$ 拟紧，并选取有限仿射开覆盖
$X = U_1 \cup \ldots \cup U_m$，其中 $U_j$ 仿射且
$\mathcal{L}|_{U_j} \cong \mathcal{O}_{U_j}$。经此同构，
像 $s|_{U_j}$ 对应于某个
$f_j \in \Gamma(U_j, \mathcal{O}_{U_j})$。于是
$X_s \cap U_j = D(f_j)$。
令 $\alpha/s^n$ 为
(\ref{properties-equation-hom-invert-s}) 的核中元素。则 $\alpha|_{X_s} = 0$。
故 $(\alpha|_{U_j})|_{D(f_j)} = 0$。由上面处理的仿射情形，
$f_j^{e_j} \alpha|_{U_j} = 0$ 对某个 $e_j \geq 0$ 成立。
令 $e = \max(e_j)$。于是 $\alpha \otimes s^e$ 在 $U_j$ 上的限制对所有
$j$ 均为零，故它为零。由于 $\alpha/s^n$ 等于
$\alpha \otimes s^e/s^{n + e}$（作为 $M_{(s)}$ 中的元素），
故如所需有 $\alpha/s^n = 0$。

\medskip\noindent
(2) 的证明。由于 $\mathcal{F}$ 是有限表现的，层
$\SheafHom_{\mathcal{O}_X}(\mathcal{F}, \mathcal{G})$ 是准凝聚的；
见《概形》第 \ref{schemes-section-quasi-coherent} 节。此外，显然对所有
$n$ 都有
$$
\SheafHom_{\mathcal{O}_X}(\mathcal{F},
\mathcal{G} \otimes_{\mathcal{O}_X} \mathcal{L}^{\otimes n}) =
\SheafHom_{\mathcal{O}_X}(\mathcal{F}, \mathcal{G})
\otimes_{\mathcal{O}_X} \mathcal{L}^{\otimes n}
$$
因此在这种情形下，结论由引理 \ref{properties-lemma-invert-s-sections} 对
$\SheafHom_{\mathcal{O}_X}(\mathcal{F}, \mathcal{G})$ 的应用得出。
\end{proof}

\section{拟仿射概形}
\label{properties-section-quasi-affine}

\begin{definition}
\label{properties-definition-quasi-affine}
若概形 $X$ 拟紧并且同构于某个仿射概形的开子概形，则称它是
{\it 拟仿射的}。
\end{definition}

\begin{lemma}
\label{properties-lemma-quasi-coherent-quasi-affine}
令 $A$ 为环，并令 $U \subset \Spec(A)$ 为拟紧开子概形。
对于准凝聚层 $\mathcal{F}$（在 $U$ 上），典范映射
$$
\widetilde{H^0(U, \mathcal{F})}|_U \to \mathcal{F}
$$
是同构。
\end{lemma}

\begin{proof}
以 $j : U \to \Spec(A)$ 表示包含态射。则
$H^0(U, \mathcal{F}) = H^0(\Spec(A), j_*\mathcal{F})$，并且
$j_*\mathcal{F}$ 由《概形》引理
\ref{schemes-lemma-push-forward-quasi-coherent} 是准凝聚的。
因此，由《概形》引理 \ref{schemes-lemma-equivalence-quasi-coherent}，
$j_*\mathcal{F} = \widetilde{H^0(U, \mathcal{F})}$。
再限制回 $U$ 即得本引理。
\end{proof}

\begin{lemma}
\label{properties-lemma-invert-f-affine}
令 $X$ 为概形。令 $f \in \Gamma(X, \mathcal{O}_X)$。
假设 $X$ 拟紧且拟分离，并假设 $X_f$ 是仿射的。则典范态射
$$
j : X \longrightarrow \Spec(\Gamma(X, \mathcal{O}_X))
$$
（来自《概形》引理 \ref{schemes-lemma-morphism-into-affine}）
诱导一个同构，把 $X_f = j^{-1}(D(f))$ 同构到标准仿射开集
$D(f) \subset \Spec(\Gamma(X, \mathcal{O}_X))$ 上。
\end{lemma}

\begin{proof}
这是显然的，因为由上面的引理 \ref{properties-lemma-invert-f-sections}，
$j$ 诱导环同构
$\Gamma(X, \mathcal{O}_X)_f \to \mathcal{O}_X(X_f)$。
\end{proof}

\begin{lemma}
\label{properties-lemma-quasi-affine}
令 $X$ 为概形。则 $X$ 是拟仿射的，当且仅当来自《概形》引理
\ref{schemes-lemma-morphism-into-affine} 的典范态射
$$
X \longrightarrow \Spec(\Gamma(X, \mathcal{O}_X))
$$
是拟紧开浸入。
\end{lemma}

\begin{proof}
若所陈列的态射是拟紧开浸入，则 $X$ 同构于
$\Spec(\Gamma(X, \mathcal{O}_X))$ 的一个拟紧开子概形，显然
$X$ 是拟仿射的。

\medskip\noindent
假设 $X$ 是拟仿射的，即 $X \subset \Spec(R)$ 是拟紧开集。
特别地，这蕴涵 $X$ 是分离的；见《概形》引理
\ref{schemes-lemma-subscheme-of-separated-scheme}。
令 $A = \Gamma(X, \mathcal{O}_X)$。
考虑环映射 $R \to A$；它来自
$R = \Gamma(\Spec(R), \mathcal{O}_{\Spec(R)})$
以及层 $\mathcal{O}_{\Spec(R)}$ 的限制映射。
由《概形》引理 \ref{schemes-lemma-morphism-into-affine}，得到包含态射的分解：
$$
X \longrightarrow
\Spec(A) \longrightarrow
\Spec(R)
$$
令 $x \in X$。选取 $r \in R$，使得
$x \in D(r)$ 且 $D(r) \subset X$。以 $f \in A$ 表示 $r$ 在 $A$ 中的像。
上面引理 \ref{properties-lemma-invert-f-sections} 的开集 $X_f$ 等于
$D(r) \subset X$，故由该引理的结论有 $A_f \cong R_r$。
因此 $D(r) \to \Spec(A)$ 是到标准仿射开集 $D(f)$（属于
$\Spec(A)$）上的同构。由于 $X$ 可由这样的仿射开集 $D(f)$ 覆盖，结论得证。
\end{proof}

\begin{lemma}
\label{properties-lemma-cartesian-diagram-quasi-affine}
令 $U \to V$ 为拟仿射概形的开浸入。则
$$
\xymatrix{
U \ar[d] \ar[rr]_-j & & \Spec(\Gamma(U, \mathcal{O}_U)) \ar[d] \\
U \ar[r] & V \ar[r]^-{j'} & \Spec(\Gamma(V, \mathcal{O}_V))
}
$$
是笛卡儿图。
\end{lemma}

\begin{proof}
由《概形》引理 \ref{schemes-lemma-morphism-into-affine}，该图交换。
记 $A = \Gamma(U, \mathcal{O}_U)$ 且
$B = \Gamma(V, \mathcal{O}_V)$。令 $g \in B$ 使得 $V_g$ 仿射且包含于
$U$。这意味着，若 $f$ 是 $g$ 在 $A$ 中的像，则 $U_f = V_g$。
由引理 \ref{properties-lemma-invert-f-affine} 可知，$j'$ 诱导一个同构，把
$V_g$ 同构到标准开集 $D(g)$（属于 $\Spec(B)$）上。因此
$V_g \times_{\Spec(B)} \Spec(A) \to \Spec(A)$ 是到
$D(f) \subset \Spec(A)$ 上的同构。再次由引理
\ref{properties-lemma-invert-f-affine}，$j$ 把 $U_f$ 同构地映到 $D(f)$。
因此可见 $U_f = U_f \times_{\Spec(B)} \Spec(A)$。由于由引理
\ref{properties-lemma-quasi-affine}，可覆盖 $U$，所用开集为上述 $V_g = U_f$，故
$U \to U \times_{\Spec(B)} \Spec(A)$ 是同构。
\end{proof}

\begin{lemma}
\label{properties-lemma-quasi-affine-presentation}
令 $X$ 为拟仿射概形。存在整数 $n \geq 0$、仿射概形 $T$ 以及态射
$T \to X$，使得对每个态射 $X' \to X$，若 $X'$ 仿射，则纤维积
$X' \times_X T$ 同构于 $\mathbf{A}^n_{X'}$，并且这是在 $X'$ 上的同构。
\end{lemma}

\begin{proof}
由定义，存在环 $A$，使得 $X$ 同构于拟紧开子概形
$U \subset \Spec(A)$。回顾：标准开集 $D(f) \subset \Spec(A)$ 构成拓扑的
一组基；见《代数》第 \ref{algebra-section-spectrum-ring} 节。因为 $U$
拟紧，可选取 $f_1, \ldots, f_n \in A$，使得
$U = D(f_1) \cup \ldots \cup D(f_n)$。因此不妨设
$X = \Spec(A) \setminus V(I)$，其中 $I = (f_1, \ldots, f_n)$。置
$$
T = \Spec(A[t, x_1, \ldots, x_n]/(f_1 x_1 + \ldots + f_n x_n - 1))
$$
结构态射 $T \to \Spec(A)$ 通过开集 $X$ 分解，从而给出态射
$T \to X$。若 $X' = \Spec(A')$，且态射 $X' \to X$ 对应于环映射
$A \to A'$，则 $f'_1, \ldots, f'_n \in A'$（即
$f_1, \ldots, f_n$ 的像）生成 $A'$ 中的单位理想。
设 $1 = f'_1 a'_1 + \ldots + f'_n a'_n$。
基变换 $X' \times_X T$ 是下列环的谱：
$A'[t, x_1, \ldots, x_n]/(f'_1 x_1 + \ldots + f'_n x_n - 1)$。
我们断言，下列 $A'$-代数同态
$$
\varphi :
A'[y_1, \ldots, y_n]
\longrightarrow
A'[t, x_1, \ldots, x_n, x_{n + 1}]/(f'_1 x_1 + \ldots + f'_n x_n - 1)
$$
把 $y_i$ 送到 $a'_i t + x_i$，并且它是同构。该断言即完成本引理的证明。
$\varphi$ 的逆由下列 $A'$-代数同态给出：
$$
\psi :
A'[t, x_1, \ldots, x_n, x_{n + 1}]/(f'_1 x_1 + \ldots + f'_n x_n - 1)
\longrightarrow
A'[y_1, \ldots, y_n]
$$
它把 $t$ 送到 $-1 + f'_1 y_1 + \ldots + f'_n y_n$，并把 $x_i$ 送到
$y_i + a'_i - a'_i(f'_1 y_1 + \ldots + f'_n y_n)$，其中
$i = 1, \ldots, n$。这是有意义的，因为 $\sum f'_ix_i$ 被映到
$$
\begin{matrix}
\sum f'_i(y_i + a'_i - a'_i(\sum f'_j y_j)) =
(\sum f'_iy_i) + 1 - (\sum f'_j y_j) = 1
\end{matrix}
$$
为看出这两个映射互逆，作如下计算：
$$
\begin{matrix}
\varphi(\psi(t) = \varphi(-1 +  \sum f'_i y_i) =
-1 + \sum f'_i (a'_i t + x_i) = t \\
\varphi(\psi(x_i)) = \varphi(y_i + a'_i - a'_i(\sum f'_j y_j)) =
a'_i t + x_i + a'_i - a'_i(\sum f'_ja'_jt + f'_jx_j) = x_i \\
\psi(\varphi(y_i)) = \psi(a'_i t + x_i) =
a'_i(-1 + \sum f'_j y_j) + y_i + a'_i - a'_i(\sum f'_j y_j) = y_i
\end{matrix}
$$
证明完毕。
\end{proof}

\section{平坦模}
\label{properties-section-flat}

\noindent
在任意环层空间 $(X, \mathcal{O}_X)$ 上，我们知道一个
$\mathcal{O}_X$-模（在一点处）平坦意味着什么；见《模》定义
\ref{modules-definition-flat}
（定义 \ref{modules-definition-flat-at-point}）。
对于仿射概形上的准凝聚层，这与代数章中定义的概念一致。

\begin{lemma}
\label{properties-lemma-flat-module}
\begin{slogan}
模的平坦性与层的平坦性相同。
\end{slogan}
令 $X = \Spec(R)$ 为仿射概形。
令 $\mathcal{F} = \widetilde{M}$，它来自某个 $R$-模 $M$。
准凝聚层 $\mathcal{F}$ 是平坦 $\mathcal{O}_X$-模，当且仅当
$M$ 是平坦 $R$-模。
\end{lemma}

\begin{proof}
$\mathcal{F}$ 的平坦性可在茎上检验；见《模》引理
\ref{modules-lemma-flat-stalks-flat}。
环上模的情形也是如此；见《代数》引理
\ref{algebra-lemma-flat-localization}。
又因为 $\mathcal{F}_x = M_{\mathfrak p}$（当 $x$ 对应于
$\mathfrak p$ 时），本引理成立。
\end{proof}

\section{局部自由模}
\label{properties-section-finite-locally-free}

\noindent
在任意环层空间上，我们知道一个 $\mathcal{O}_X$-模是（有限）局部自由的
意味着什么。在仿射概形上，这与代数章中定义的概念一致。

\begin{lemma}
\label{properties-lemma-locally-free-module}
令 $X = \Spec(R)$ 为仿射概形。
令 $\mathcal{F} = \widetilde{M}$，它来自某个 $R$-模 $M$。
准凝聚层 $\mathcal{F}$ 是（有限）局部自由 $\mathcal{O}_X$-模，当且仅当
$M$ 是（有限）局部自由 $R$-模。
\end{lemma}

\begin{proof}
这由定义得出；见《模》定义
\ref{modules-definition-locally-free}
以及《代数》定义 \ref{algebra-definition-locally-free}。
\end{proof}

\noindent
有限局部自由模可以用许多不同方式刻画。

\begin{lemma}
\label{properties-lemma-finite-locally-free}
令 $X$ 为概形。
令 $\mathcal{F}$ 为准凝聚 $\mathcal{O}_X$-模。
下列条件等价：
\begin{enumerate}
\item $\mathcal{F}$ 是有限表现的平坦 $\mathcal{O}_X$-模；
\item $\mathcal{F}$ 是有限表现 $\mathcal{O}_X$-模，并且对所有
$x \in X$，茎 $\mathcal{F}_x$ 都是自由 $\mathcal{O}_{X, x}$-模；
\item $\mathcal{F}$ 是有限型局部自由 $\mathcal{O}_X$-模；
\item $\mathcal{F}$ 是有限局部自由 $\mathcal{O}_X$-模；
\item $\mathcal{F}$ 是有限型 $\mathcal{O}_X$-模，对每个 $x \in X$，
茎 $\mathcal{F}_x$ 都是自由 $\mathcal{O}_{X, x}$-模，并且函数
$$
\rho_\mathcal{F} : X \to \mathbf{Z}, \quad
x \longmapsto
\dim_{\kappa(x)} \mathcal{F}_x \otimes_{\mathcal{O}_{X, x}} \kappa(x)
$$
在 $X$ 的 Zariski 拓扑中局部常值。
\end{enumerate}
\end{lemma}

\begin{proof}
本引理立即归结为仿射情形。在这种情形下，它是《代数》引理
\ref{algebra-lemma-finite-projective} 的重述。该转译使用引理
\ref{properties-lemma-finite-type-module}、
\ref{properties-lemma-finite-presentation-module}、
\ref{properties-lemma-flat-module} 和
\ref{properties-lemma-locally-free-module}。
\end{proof}

\begin{lemma}
\label{properties-lemma-finite-locally-free-reduced}
令 $X$ 为约化概形。令 $\mathcal{F}$ 为准凝聚
$\mathcal{O}_X$-模。则引理 \ref{properties-lemma-finite-locally-free} 的等价条件
还等价于
\begin{enumerate}
\item[(6)] $\mathcal{F}$ 是有限型 $\mathcal{O}_X$-模，并且函数
$$
\rho_\mathcal{F} : X \to \mathbf{Z}, \quad
x \longmapsto
\dim_{\kappa(x)} \mathcal{F}_x \otimes_{\mathcal{O}_{X, x}} \kappa(x)
$$
在 $X$ 的 Zariski 拓扑中局部常值。
\end{enumerate}
\end{lemma}

\begin{proof}
本引理立即归结为仿射情形。在这种情形下，它是《代数》引理
\ref{algebra-lemma-finite-projective-reduced} 的重述。
\end{proof}

\section{局部投射模}
\label{properties-section-locally-projective}

\noindent
代数章中已有工作的一个推论是，可以如下定义局部投射模。

\begin{definition}
\label{properties-definition-locally-projective}
令 $X$ 为概形。令 $\mathcal{F}$ 为准凝聚
$\mathcal{O}_X$-模。称 $\mathcal{F}$ 是{\it 局部投射的}，如果对每个
仿射开集 $U \subset X$，$\mathcal{O}_X(U)$-模
$\mathcal{F}(U)$ 都是投射的。
\end{definition}

\begin{lemma}
\label{properties-lemma-locally-projective}
令 $X$ 为概形。
令 $\mathcal{F}$ 为准凝聚 $\mathcal{O}_X$-模。
下列条件等价：
\begin{enumerate}
\item $\mathcal{F}$ 局部投射；
\item 存在仿射开覆盖 $X = \bigcup U_i$，使得
$\mathcal{O}_X(U_i)$-模 $\mathcal{F}(U_i)$ 对每个 $i$ 都是投射的。
\end{enumerate}
特别地，若 $X = \Spec(A)$ 且 $\mathcal{F} = \widetilde{M}$，
则 $\mathcal{F}$ 局部投射，当且仅当 $M$ 是投射 $A$-模。
\end{lemma}

\begin{proof}
首先注意，若 $M$ 是投射 $A$-模且 $A \to B$ 是环映射，则
$M \otimes_A B$ 是投射 $B$-模；见《代数》引理
\ref{algebra-lemma-ascend-properties-modules}。
因此，若 $U$ 是仿射开集，且 $\mathcal{F}(U)$ 是投射
$\mathcal{O}_X(U)$-模，则标准开集 $D(f)$ 是仿射开集，并且
$\mathcal{F}(D(f))$ 是投射 $\mathcal{O}_X(D(f))$-模，这对所有
$f \in \mathcal{O}_X(U)$ 都成立。
假设 (2) 成立。令 $U \subset X$ 为任意仿射开集。
可以找到由有限多个标准开集组成的开覆盖
$U = \bigcup_{j = 1, \ldots, m} D(f_j)$；其中这些标准开集
$D(f_j)$ 满足：对每个 $j$，开集 $D(f_j)$ 都是某个 $U_i$ 的标准开集；
见《概形》引理
\ref{schemes-lemma-standard-open-two-affines}。
因此，若置 $A = \mathcal{O}_X(U)$，且 $M$ 是一个 $A$-模，使
$\mathcal{F}|_U$ 对应于 $M$，则 $M_{f_j}$ 是投射
$A_{f_j}$-模。由此，$A \to B = \prod A_{f_j}$ 是忠实平坦环映射，
且 $M \otimes_A B$ 是投射 $B$-模。因此由《代数》定理
\ref{algebra-theorem-ffdescent-projectivity}，$M$ 是投射的。
\end{proof}

\begin{lemma}
\label{properties-lemma-locally-projective-pullback}
令 $f : X \to Y$ 为概形态射。
令 $\mathcal{G}$ 为准凝聚 $\mathcal{O}_Y$-模。
若 $\mathcal{G}$ 在 $Y$ 上局部投射，则 $f^*\mathcal{G}$
在 $X$ 上局部投射。
\end{lemma}

\begin{proof}
这由《代数》引理 \ref{algebra-lemma-ascend-properties-modules}
和引理 \ref{properties-lemma-locally-projective} 得出。
\end{proof}

\section{准凝聚层的延拓}
\label{properties-section-extending-quasi-coherent-sheaves}

\noindent
有时，能够证明开子概形上给定的准凝聚层可延拓到整个概形，是很有用的。

\begin{lemma}
\label{properties-lemma-extend-trivial}
令 $j : U \to X$ 为概形的拟紧开浸入。
\begin{enumerate}
\item $U$ 上任意准凝聚层都可延拓为 $X$ 上的准凝聚层。
\item 令 $\mathcal{F}$ 为 $X$ 上的准凝聚层。令
$\mathcal{G} \subset \mathcal{F}|_U$ 为准凝聚子层。则存在
准凝聚子层 $\mathcal{H}$（属于 $\mathcal{F}$），使得
$\mathcal{H}|_U = \mathcal{G}$，这里的等式是在
$\mathcal{F}|_U$ 的子层意义下成立。
\item 令 $\mathcal{F}$ 为 $X$ 上的准凝聚层。令 $\mathcal{G}$ 为
$U$ 上的准凝聚层。令
$\varphi : \mathcal{G} \to \mathcal{F}|_U$ 为
$\mathcal{O}_U$-模态射。则存在准凝聚层 $\mathcal{H}$（为
$\mathcal{O}_X$-模）和映射 $\psi : \mathcal{H} \to \mathcal{F}$，使得
$\mathcal{H}|_U = \mathcal{G}$ 且 $\psi|_U = \varphi$。
\end{enumerate}
\end{lemma}

\begin{proof}
浸入是分离的（见《概形》引理
\ref{schemes-lemma-immersions-monomorphisms}），并且 $j$ 依假设拟紧。
因此，对任意准凝聚层 $\mathcal{G}$（在 $U$ 上），层
$j_*\mathcal{G}$ 是到
$X$ 上的一个延拓。见《概形》引理
\ref{schemes-lemma-push-forward-quasi-coherent} 和《层》第
\ref{sheaves-section-open-immersions} 节。

\medskip\noindent
假设 $\mathcal{F}$、$\mathcal{G}$ 如 (2) 中。则
$j_*\mathcal{G}$ 是 $X$ 上的准凝聚层（见上文）。
它是 $j_*j^*\mathcal{F}$ 的子层。因此核
$$
\mathcal{H} =
\Ker(\mathcal{F} \oplus j_* \mathcal{G}
\longrightarrow j_*j^*\mathcal{F})
$$
也是准凝聚的；见《概形》第 \ref{schemes-section-quasi-coherent} 节。
形式地检验可知 $\mathcal{H} \subset \mathcal{F}$ 且
$\mathcal{H}|_U = \mathcal{G}$（再次使用《层》第
\ref{sheaves-section-open-immersions} 节的材料）。

\medskip\noindent
(3) 的证明方式与 (2) 相同。只需取
$\mathcal{H} = \Ker(\mathcal{F} \oplus j_* \mathcal{G}
\to j_*j^*\mathcal{F})$，连同它到 $\mathcal{F}$ 的显然映射，
以及它与 $\mathcal{G}$ 的显然等同（在 $U$ 上）。
\end{proof}

\begin{lemma}
\label{properties-lemma-extend}
令 $X$ 为拟紧且拟分离概形。令 $U \subset X$ 为拟紧开集。
令 $\mathcal{F}$ 为准凝聚 $\mathcal{O}_X$-模。
令 $\mathcal{G} \subset \mathcal{F}|_U$ 为有限型准凝聚
$\mathcal{O}_U$-子模。则存在有限型准凝聚子模
$\mathcal{G}' \subset \mathcal{F}$，使得
$\mathcal{G}'|_U = \mathcal{G}$。
\end{lemma}

\begin{proof}
令 $n$ 为仿射开集 $U_i \subset X$ 的最小数目，其中
$i = 1, \ldots , n$ 且 $X = U \cup \bigcup U_i$。
（这里使用 $X$ 拟紧。）假设可证明 $n = 1$ 的情形。于是可依次把
$\mathcal{G}$ 延拓为 $\mathcal{G}_1$（定义在 $U \cup U_1$ 上），
再延拓为 $\mathcal{G}_2$（定义在 $U \cup U_1 \cup U_2$ 上），
再延拓为 $\mathcal{G}_3$（定义在 $U \cup U_1 \cup U_2 \cup U_3$ 上），
依此类推。因此归结为 $n = 1$ 的情形。

\medskip\noindent
所以可假设 $X = U \cup V$，其中 $V$ 仿射。
因为 $X$ 拟分离，并且 $U$、$V$ 是拟紧开集，所以
$U \cap V$ 是拟紧开集。只需对系统
$(V, U \cap V, \mathcal{F}|_V, \mathcal{G}|_{U \cap V})$
证明本引理，因为可把所得的层 $\mathcal{G}'$（在 $V$ 上）与给定的
层 $\mathcal{G}$（在 $U$ 上）沿着它们在 $U \cap V$ 上的公共值粘合。
因此归结为 $X$ 仿射的情形。

\medskip\noindent
假设 $X = \Spec(R)$。写成 $\mathcal{F} = \widetilde M$，其中
它来自某个 $R$-模 $M$。由上面的引理 \ref{properties-lemma-extend-trivial}，可找到
准凝聚子层 $\mathcal{H} \subset \mathcal{F}$，它限制为
$\mathcal{G}$（在 $U$ 上）。写成 $\mathcal{H} = \widetilde N$，其中它来自某个
$R$-模 $N$。对每个 $u \in U$，存在 $f \in R$，使得
$u \in D(f) \subset U$ 且 $N_f$ 有限生成；见引理
\ref{properties-lemma-finite-type-module}。
由于 $U$ 拟紧，可用有限多个 $D(f_i)$ 覆盖它，使得
$N_{f_i}$ 由有限多个元素生成，设这些元素为
$x_{i, 1}/f_i^N, \ldots, x_{i, r_i}/f_i^N$。
令 $N' \subset N$ 为由元素 $x_{i, j}$ 生成的子模。则子层
$\mathcal{G}' = \widetilde{N'} \subset \mathcal{H} \subset \mathcal{F}$
满足要求。
\end{proof}

\begin{lemma}
\label{properties-lemma-quasi-coherent-colimit-finite-type}
令 $X$ 为拟紧且拟分离概形。
任意准凝聚 $\mathcal{O}_X$-模层都是其有限型准凝聚
$\mathcal{O}_X$-子模的有向余极限。
\end{lemma}

\begin{proof}
该余极限是有向的，因为若 $\mathcal{G}_1$、$\mathcal{G}_2$
是有限型准凝聚子层，则
$\mathcal{G}_1 \oplus \mathcal{G}_2 \to \mathcal{F}$ 的像是有限型
准凝聚子模。
令 $U \subset X$ 为任意仿射开集，并令
$s \in \Gamma(U, \mathcal{F})$ 为任意截面。
令 $\mathcal{G} \subset \mathcal{F}|_U$ 为由 $s$ 生成的子层。
则显然 $\mathcal{G}$ 是准凝聚的，并且作为
$\mathcal{O}_U$-模是有限型的。由引理 \ref{properties-lemma-extend} 可知，
$\mathcal{G}$ 是某个有限型准凝聚子层
$\mathcal{G}' \subset \mathcal{F}$ 的限制。由于 $X$ 的拓扑具有由仿射
开集组成的一组基，可知 $\mathcal{F}$ 的每个局部截面都局部地包含在
某个有限型准凝聚子模中。故结论成立。
\end{proof}

\begin{lemma}
\label{properties-lemma-extend-finite-presentation}
令 $X$ 为拟紧且拟分离概形。
令 $\mathcal{F}$ 为准凝聚 $\mathcal{O}_X$-模。
令 $U \subset X$ 为拟紧开集。
令 $\mathcal{G}$ 为有限表现 $\mathcal{O}_U$-模。
令 $\varphi : \mathcal{G} \to \mathcal{F}|_U$ 为
$\mathcal{O}_U$-模态射。
则存在有限表现 $\mathcal{O}_X$-模 $\mathcal{G}'$，以及
$\mathcal{O}_X$-模态射 $\varphi' : \mathcal{G}' \to \mathcal{F}$，
使得 $\mathcal{G}'|_U = \mathcal{G}$ 且
$\varphi'|_U = \varphi$。
\end{lemma}

\begin{proof}
证明的开头重复引理 \ref{properties-lemma-extend} 证明的开头。不过我们仍将它仔细写出。

\medskip\noindent
令 $n$ 为仿射开集 $U_i \subset X$ 的最小数目，其中
$i = 1, \ldots , n$ 且 $X = U \cup \bigcup U_i$。
（这里使用 $X$ 拟紧。）假设可证明 $n = 1$ 的情形。于是可依次把对
$(\mathcal{G}, \varphi)$ 延拓为对
$(\mathcal{G}_1, \varphi_1)$（定义在 $U \cup U_1$ 上），
再延拓为对 $(\mathcal{G}_2, \varphi_2)$（定义在
$U \cup U_1 \cup U_2$ 上），
再延拓为对 $(\mathcal{G}_3, \varphi_3)$（定义在
$U \cup U_1 \cup U_2 \cup U_3$ 上），依此类推。
因此归结为 $n = 1$ 的情形。

\medskip\noindent
所以可假设 $X = U \cup V$，其中 $V$ 仿射。
因为 $X$ 拟分离且 $U$ 拟紧，所以 $U \cap V \subset V$ 拟紧。
假设对系统
$(V, U \cap V, \mathcal{F}|_V, \mathcal{G}|_{U \cap V}, \varphi|_{U \cap V})$
证明了本引理，从而产生
$(\mathcal{G}', \varphi')$（在 $V$ 上）。那么可把
$\mathcal{G}'$（在 $V$ 上）与给定的层 $\mathcal{G}$（在 $U$ 上）
沿着它们在 $U \cap V$ 上的公共值粘合；
类似地，可把映射 $\varphi'$ 与映射 $\varphi$ 沿着它们在
$U \cap V$ 上的公共值粘合。因此归结为 $X$ 仿射的情形。

\medskip\noindent
假设 $X = \Spec(R)$。
由上面的引理 \ref{properties-lemma-extend-trivial}，可找到准凝聚层
$\mathcal{H}$，以及映射
$\psi : \mathcal{H} \to \mathcal{F}$（在 $X$ 上）；它限制为
$\mathcal{G}$ 和 $\varphi$（在 $U$ 上）。
由引理 \ref{properties-lemma-extend}，可找到有限型准凝聚
$\mathcal{O}_X$-子模 $\mathcal{H}' \subset \mathcal{H}$，使得
$\mathcal{H}'|_U = \mathcal{G}$。因此，替换 $\mathcal{H}$，所用的是
$\mathcal{H}'$；并替换 $\psi$，所用的是 $\psi$ 在 $\mathcal{H}'$ 上的限制。此后，
可假设 $\mathcal{H}$ 是有限型的。
由引理 \ref{properties-lemma-finite-presentation-module} 可知
$\mathcal{H} = \widetilde{N}$，其中 $N$ 是有限生成 $R$-模。
因此存在如下准凝聚 $\mathcal{O}_X$-模短正合列中的满射：
$$
0 \to \mathcal{K} \to \mathcal{O}_X^{\oplus n} \to \mathcal{H} \to 0
$$
其中 $\mathcal{K}$ 定义为该核。
由于 $\mathcal{G}$ 是有限表现的，且由《模》引理
\ref{modules-lemma-kernel-surjection-finite-free-onto-finite-presentation}
有 $\mathcal{H}|_U = \mathcal{G}$，故限制
$\mathcal{K}|_U$ 是有限型 $\mathcal{O}_U$-模。因此再次由引理
\ref{properties-lemma-extend}，存在有限型准凝聚 $\mathcal{O}_X$-子模
$\mathcal{K}' \subset \mathcal{K}$，使得
$\mathcal{K}'|_U = \mathcal{K}|_U$。本引理所提问题的解是置
$$
\mathcal{G}' = \mathcal{O}_X^{\oplus n}/\mathcal{K}'
$$
它显然是有限表现的，并且限制后给出 $\mathcal{G}$（在 $U$ 上），
其中 $\varphi'$ 等于复合
$$
\mathcal{G}' = \mathcal{O}_X^{\oplus n}/\mathcal{K}'
\to \mathcal{O}_X^{\oplus n}/\mathcal{K} = \mathcal{H} \xrightarrow{\psi}
\mathcal{F}.
$$
本引理证毕。
\end{proof}

\begin{lemma}
\label{properties-lemma-lift-finite-presentation}
令 $X$ 为拟紧且拟分离概形。令 $U \subset X$ 为拟紧开集。
令 $\mathcal{G}$ 为 $\mathcal{O}_U$-模。
\begin{enumerate}
\item 若 $\mathcal{G}$ 是有限型准凝聚模，则存在有限型准凝聚
$\mathcal{O}_X$-模 $\mathcal{G}'$，使得
$\mathcal{G}'|_U = \mathcal{G}$。
\item 若 $\mathcal{G}$ 是有限表现的，则存在有限表现
$\mathcal{O}_X$-模 $\mathcal{G}'$，使得
$\mathcal{G}'|_U = \mathcal{G}$。
\end{enumerate}
\end{lemma}

\begin{proof}
(2) 是引理 \ref{properties-lemma-extend-finite-presentation} 中
$\mathcal{F} = 0$ 的特殊情形。对于 (1)，先由引理
\ref{properties-lemma-extend-trivial} 把
$\mathcal{G} = \mathcal{F}|_U$ 写为某个准凝聚
$\mathcal{O}_X$-模的限制，然后对 $\mathcal{G} = \mathcal{F}|_U$
应用引理 \ref{properties-lemma-extend}。
\end{proof}

\noindent
下述引理说明，拟紧且拟分离概形上的每个准凝聚层都是有限表现
$\mathcal{O}$-模的滤过余极限。事实上，在下一个引理中，我们将其改述为
（也许更熟悉的）有向集上的有向余极限。

\begin{lemma}
\label{properties-lemma-directed-colimit-diagram-finite-presentation}
\begin{slogan}
拟紧且拟分离概形上的准凝聚模是有限表现模的滤过余极限。
\end{slogan}
令 $X$ 为概形。假设 $X$ 拟紧且拟分离。
令 $\mathcal{F}$ 为准凝聚 $\mathcal{O}_X$-模。
则存在
\begin{enumerate}
\item 一个滤过指标范畴 $\mathcal{I}$（见《范畴》定义
\ref{categories-definition-directed}）；
\item 一个图 $\mathcal{I} \to \textit{Mod}(\mathcal{O}_X)$（见《范畴》第
\ref{categories-section-limits} 节），
$i \mapsto \mathcal{F}_i$；
\item $\mathcal{O}_X$-模态射
$\varphi_i : \mathcal{F}_i \to \mathcal{F}$；
\end{enumerate}
使得每个 $\mathcal{F}_i$ 都是有限表现的，并且态射 $\varphi_i$
诱导同构
$$
\colim_i \mathcal{F}_i
=
\mathcal{F}.
$$
\end{lemma}

\begin{proof}
选取集合 $I$，并对每个 $i \in I$ 选取一个有限表现
$\mathcal{O}_X$-模，以及一个 $\mathcal{O}_X$-模同态
$\varphi_i : \mathcal{F}_i \to \mathcal{F}$，使其具有如下性质：
对任意 $\psi : \mathcal{G} \to \mathcal{F}$，若 $\mathcal{G}$
有限表现，则存在 $i \in I$，以及同构
$\alpha : \mathcal{F}_i \to \mathcal{G}$，使得
$\varphi_i = \psi \circ \alpha$。由《模》引理
\ref{modules-lemma-set-isomorphism-classes-finite-type-modules}
显然存在这样的集合（也见其证明）。
以 $\mathcal{I}$ 表示如下范畴：
$\Ob(\mathcal{I}) = I$，并且给定 $i, i' \in I$，置
$$
\Mor_\mathcal{I}(i, i') =
\{\alpha : \mathcal{F}_i \to \mathcal{F}_{i'} \mid
\alpha \circ \varphi_{i'} = \varphi_i
\}.
$$
我们断言 $\mathcal{I}$ 是滤过范畴，并且
$\mathcal{F} = \colim_i \mathcal{F}_i$。

\medskip\noindent
令 $i, i' \in I$。可以考虑态射
$$
\mathcal{F}_i \oplus \mathcal{F}_{i'} \longrightarrow \mathcal{F}
$$
，它是 $\varphi_i$ 与 $\varphi_{i'}$ 的直和。
因为有限表现 $\mathcal{O}_X$-模的直和仍有限表现，故存在某个
$i'' \in I$，使得
$\varphi_{i''} : \mathcal{F}_{i''} \to \mathcal{F}$
同构于上面陈列的到 $\mathcal{F}$ 的箭头。
由于有交换图
$$
\xymatrix{
\mathcal{F}_i \ar[r] \ar[d] & \mathcal{F} \ar@{=}[d] \\
\mathcal{F}_i \oplus \mathcal{F}_{i'} \ar[r] & \mathcal{F}
}
\quad
\text{且}
\quad
\xymatrix{
\mathcal{F}_{i'} \ar[r] \ar[d] & \mathcal{F} \ar@{=}[d] \\
\mathcal{F}_i \oplus \mathcal{F}_{i'} \ar[r] & \mathcal{F}
}
$$
可知有态射 $i \to i''$ 和 $i' \to i''$，它们属于 $\mathcal{I}$。
接着，假设给定 $i, i' \in I$ 以及态射
$\alpha, \beta : i \to i'$（它们对应于 $\mathcal{O}_X$-模映射
$\alpha, \beta : \mathcal{F}_i \to \mathcal{F}_{i'}$）。
此时考虑余等化子
$$
\mathcal{G} =
\Coker(
\mathcal{F}_i \xrightarrow{\alpha - \beta} \mathcal{F}_{i'}
)
$$
注意，$\mathcal{G}$ 是有限表现 $\mathcal{O}_X$-模。
因为根据范畴 $\mathcal{I}$ 中态射的定义，
$\varphi_{i'} \circ \alpha = \varphi_{i'} \circ \beta$，
所以得到诱导映射 $\psi : \mathcal{G} \to \mathcal{F}$。
因此，对 $(\mathcal{G}, \psi)$ 再次同构于对
$(\mathcal{F}_{i''}, \varphi_{i''})$，其中某个 $i''$。
从而存在态射 $i' \to i''$（在 $\mathcal{I}$ 中），它等化
$\alpha$ 与 $\beta$。故已证明范畴 $\mathcal{I}$ 是滤过的。

\medskip\noindent
还须证明该图的余极限是 $\mathcal{F}$。根据余极限的定义以及我们对范畴
$\mathcal{I}$ 的定义，有典范映射
$$
\varphi :
\colim_i \mathcal{F}_i
\longrightarrow
\mathcal{F}.
$$
选取 $x \in X$。我们来证明 $\varphi_x$ 是同构。回顾
$$
(\colim_i \mathcal{F}_i)_x
=
\colim_i \mathcal{F}_{i, x},
$$
见《层》第 \ref{sheaves-section-limits-sheaves} 节。
先证明映射 $\varphi_x$ 是单射。
假设 $s \in \mathcal{F}_{i, x}$ 为一个元素，并且 $s$ 映到
$\mathcal{F}_x$ 中的零。则存在拟紧开集 $U$，使得 $s$ 来自
$s \in \mathcal{F}_i(U)$，并且
$\varphi_i(s) = 0$ 在 $\mathcal{F}(U)$ 中成立。
由引理 \ref{properties-lemma-extend}，可找到有限型准凝聚子层
$\mathcal{K} \subset \Ker(\varphi_i)$，它限制为准凝聚
$\mathcal{O}_U$-子模（属于 $\mathcal{F}_i$），而该子模由 $s$ 生成：
$\mathcal{K}|_U = \mathcal{O}_U\cdot s \subset \mathcal{F}_i|_U$。
显然，$\mathcal{F}_i/\mathcal{K}$ 是有限表现的，并且映射
$\varphi_i$ 通过商映射
$\mathcal{F}_i \to \mathcal{F}_i/\mathcal{K}$ 分解。因此可找到
$i' \in I$ 以及态射
$\alpha : \mathcal{F}_i \to \mathcal{F}_{i'}$（在 $\mathcal{I}$ 中），
它可与商映射
$\mathcal{F}_i \to \mathcal{F}_i/\mathcal{K}$ 等同。
于是截面 $s$ 映到 $\mathcal{F}_{i'}(U)$ 中的零，特别地，它在
$(\colim_i \mathcal{F}_i)_x =
\colim_i \mathcal{F}_{i, x}$ 中也为零。单射性得证。
最后证明映射 $\varphi_x$ 是满射。选取 $s \in \mathcal{F}_x$。
选取拟紧开邻域 $U \subset X$（它包含 $x$），使得 $s$ 对应于截面
$s \in \mathcal{F}(U)$。考虑映射
$s : \mathcal{O}_U \to \mathcal{F}$（乘以 $s$）。
由引理 \ref{properties-lemma-extend-finite-presentation}，存在有限表现
$\mathcal{O}_X$-模 $\mathcal{G}$，以及 $\mathcal{O}_X$-模映射
$\mathcal{G} \to \mathcal{F}$，使得 $\mathcal{G}|_U \to \mathcal{F}|_U$
可等同于
$s : \mathcal{O}_U \to \mathcal{F}$。
再次根据 $\mathcal{I}$ 的定义，存在 $i \in I$，使得
$\mathcal{G} \to \mathcal{F}$ 同构于
$\varphi_i : \mathcal{F}_i \to \mathcal{F}$。显然存在截面
$s' \in \mathcal{F}_i(U)$，它映到 $s \in \mathcal{F}(U)$。
这证明了满射性，也完成了本引理的证明。
\end{proof}

\begin{lemma}
\label{properties-lemma-directed-colimit-finite-presentation}
令 $X$ 为概形。假设 $X$ 拟紧且拟分离。
令 $\mathcal{F}$ 为准凝聚 $\mathcal{O}_X$-模。
则存在
\begin{enumerate}
\item 一个有向集 $I$（见《范畴》定义
\ref{categories-definition-directed-set}）；
\item 一个系统 $(\mathcal{F}_i, \varphi_{ii'})$，它定义在 $I$ 上并属于
$\textit{Mod}(\mathcal{O}_X)$（见《范畴》定义
\ref{categories-definition-system-over-poset}）；
\item $\mathcal{O}_X$-模态射
$\varphi_i : \mathcal{F}_i \to \mathcal{F}$；
\end{enumerate}
使得每个 $\mathcal{F}_i$ 都是有限表现的，并且态射 $\varphi_i$
诱导同构
$$
\colim_i \mathcal{F}_i
=
\mathcal{F}.
$$
\end{lemma}

\begin{proof}
这是引理 \ref{properties-lemma-directed-colimit-diagram-finite-presentation} 与
《范畴》引理 \ref{categories-lemma-directed-category-system} 的直接推论
（再结合 $\mathcal{O}_X$-模层范畴中余极限存在这一事实；见《层》第
\ref{sheaves-section-limits-sheaves} 节）。
\end{proof}

\begin{lemma}
\label{properties-lemma-finite-directed-colimit-surjective-maps}
令 $X$ 为概形。假设 $X$ 拟紧且拟分离。
令 $\mathcal{F}$ 为有限型准凝聚 $\mathcal{O}_X$-模。
则可写成 $\mathcal{F} = \colim \mathcal{F}_i$，其中
$\mathcal{F}_i$ 是有限表现的，并且所有过渡映射
$\mathcal{F}_i \to \mathcal{F}_{i'}$ 都是满射。
\end{lemma}

\begin{proof}
把 $\mathcal{F} = \colim \mathcal{G}_i$ 写成有限表现
$\mathcal{O}_X$-模的滤过余极限
（引理 \ref{properties-lemma-directed-colimit-finite-presentation}）。
我们断言 $\mathcal{G}_i \to \mathcal{F}$ 对某个 $i$ 是满射。
事实上，选取有限仿射开覆盖 $X = U_1 \cup \ldots \cup U_m$。
选取截面 $s_{jl} \in \mathcal{F}(U_j)$ 来生成
$\mathcal{F}|_{U_j}$；见引理 \ref{properties-lemma-finite-type-module}。
由《层》引理 \ref{sheaves-lemma-directed-colimits-sections}，
$s_{jl}$ 属于 $\mathcal{G}_i \to \mathcal{F}$ 的像，只要 $i$ 足够大。
因此，$\mathcal{G}_i \to \mathcal{F}$ 是满射，只要 $i$ 足够大。
选取这样的 $i$，并令 $\mathcal{K} \subset \mathcal{G}_i$ 为映射
$\mathcal{G}_i \to \mathcal{F}$ 的核。把
$\mathcal{K} = \colim \mathcal{K}_a$ 写成其有限型准凝聚子模的滤过余极限
（引理 \ref{properties-lemma-quasi-coherent-colimit-finite-type}）。于是
$\mathcal{F} = \colim \mathcal{G}_i/\mathcal{K}_a$ 是本引理所提问题的解。
\end{proof}

\begin{lemma}
\label{properties-lemma-application-directed-colimit}
令 $X$ 为拟紧且拟分离概形。
令 $\mathcal{F}$ 为有限型准凝聚 $\mathcal{O}_X$-模。
令 $U \subset X$ 为拟紧开集，并使 $\mathcal{F}|_U$ 是有限表现的。
则存在 $\mathcal{O}_X$-模映射
$\varphi : \mathcal{G} \to \mathcal{F}$，满足
(a) $\mathcal{G}$ 有限表现，
(b) $\varphi$ 为满射，并且
(c) $\varphi|_U$ 为同构。
\end{lemma}

\begin{proof}
把 $\mathcal{F} = \colim \mathcal{F}_i$ 写成有向余极限，其中每个
$\mathcal{F}_i$ 都有限表现；见引理
\ref{properties-lemma-directed-colimit-finite-presentation}。
选取有限仿射开覆盖 $X = \bigcup V_j$，并选取有限多个截面
$s_{jl} \in \mathcal{F}(V_j)$ 生成 $\mathcal{F}|_{V_j}$；见引理
\ref{properties-lemma-finite-type-module}。
由《层》引理 \ref{sheaves-lemma-directed-colimits-sections}，
$s_{jl}$ 属于 $\mathcal{F}_i \to \mathcal{F}$ 的像，只要 $i$ 足够大。
因此，$\mathcal{F}_i \to \mathcal{F}$ 是满射，只要 $i$ 足够大。
选取这样的 $i$，并令 $\mathcal{K} \subset \mathcal{F}_i$ 为映射
$\mathcal{F}_i \to \mathcal{F}$ 的核。由于 $\mathcal{F}_U$ 是有限表现的，
可知 $\mathcal{K}|_U$ 是有限型的；见《模》引理
\ref{modules-lemma-kernel-surjection-finite-free-onto-finite-presentation}。
因此可找到有限型准凝聚子模
$\mathcal{K}' \subset \mathcal{K}$，满足
$\mathcal{K}'|_U = \mathcal{K}|_U$；见引理 \ref{properties-lemma-extend}。于是
$\mathcal{G} = \mathcal{F}_i/\mathcal{K}'$ 连同给定映射
$\mathcal{G} \to \mathcal{F}$ 即为所求。
\end{proof}

\noindent
令 $X$ 为概形。在下述引理中，我们使用{\it 有限表现准凝聚
$\mathcal{O}_X$-代数 $\mathcal{A}$}这一概念。它意味着，对每个仿射开集
$\Spec(R) \subset X$，都有 $\mathcal{A} = \widetilde{A}$，其中 $A$
是一个（交换）$R$-代数，并且作为 $R$-代数有限表现。

\begin{lemma}
\label{properties-lemma-algebra-directed-colimit-finite-presentation}
令 $X$ 为概形。假设 $X$ 拟紧且拟分离。
令 $\mathcal{A}$ 为准凝聚 $\mathcal{O}_X$-代数。
则存在
\begin{enumerate}
\item 一个有向集 $I$（见《范畴》定义
\ref{categories-definition-directed-set}）；
\item 一个系统 $(\mathcal{A}_i, \varphi_{ii'})$，它定义在 $I$ 上并属于
$\mathcal{O}_X$-代数范畴；
\item $\mathcal{O}_X$-代数态射
$\varphi_i : \mathcal{A}_i \to \mathcal{A}$；
\end{enumerate}
使得每个 $\mathcal{A}_i$ 都是有限表现准凝聚
$\mathcal{O}_X$-代数，并且态射 $\varphi_i$ 诱导同构
$$
\colim_i \mathcal{A}_i
=
\mathcal{A}.
$$
\end{lemma}

\begin{proof}
先如引理 \ref{properties-lemma-directed-colimit-finite-presentation}，
把 $\mathcal{A} = \colim_i \mathcal{F}_i$ 写成有限表现准凝聚层的有向余极限。
对每个 $i$，令 $\mathcal{B}_i = \text{Sym}(\mathcal{F}_i)$ 为
$\mathcal{F}_i$ 在 $\mathcal{O}_X$ 上的对称代数。记
$\mathcal{I}_i = \Ker(\mathcal{B}_i \to \mathcal{A})$。写成
$\mathcal{I}_i = \colim_j \mathcal{F}_{i, j}$，其中
$\mathcal{F}_{i, j}$ 是 $\mathcal{I}_i$ 的有限型准凝聚子模；见引理
\ref{properties-lemma-quasi-coherent-colimit-finite-type}。
令 $\mathcal{I}_{i, j} \subset \mathcal{I}_i$ 等于
$\mathcal{B}_i$-理想，该理想由 $\mathcal{F}_{i, j}$ 生成。置
$\mathcal{A}_{i, j} = \mathcal{B}_i/\mathcal{I}_{i, j}$。
则 $\mathcal{A}_{i, j}$ 是有限表现准凝聚
$\mathcal{O}_X$-代数。定义 $(i, j) \leq (i', j')$ 当且仅当
$i \leq i'$，并且映射 $\mathcal{B}_i \to \mathcal{B}_{i'}$
把理想 $\mathcal{I}_{i, j}$ 映入理想 $\mathcal{I}_{i', j'}$。
于是显然 $\mathcal{A} = \colim_{i, j} \mathcal{A}_{i, j}$。
\end{proof}

\noindent
令 $X$ 为概形。在下述引理中，我们使用{\it 有限型准凝聚
$\mathcal{O}_X$-代数 $\mathcal{A}$}这一概念。它意味着，对每个仿射开集
$\Spec(R) \subset X$，都有 $\mathcal{A} = \widetilde{A}$，其中 $A$
是一个（交换）$R$-代数，并且作为 $R$-代数有限型。

\begin{lemma}
\label{properties-lemma-algebra-directed-colimit-finite-type}
令 $X$ 为概形。假设 $X$ 拟紧且拟分离。
令 $\mathcal{A}$ 为准凝聚 $\mathcal{O}_X$-代数。
则 $\mathcal{A}$ 是其有限型准凝聚 $\mathcal{O}_X$-子代数的有向余极限。
\end{lemma}

\begin{proof}
若 $\mathcal{A}_1, \mathcal{A}_2 \subset \mathcal{A}$ 是有限型准凝聚
$\mathcal{O}_X$-子代数，则
$\mathcal{A}_1 \otimes_{\mathcal{O}_X} \mathcal{A}_2 \to \mathcal{A}$
的像也是有限型准凝聚 $\mathcal{O}_X$-子代数（略去若干细节），并且包含
$\mathcal{A}_1$ 和 $\mathcal{A}_2$。由此可见该系统是有向的。
为证明 $\mathcal{A}$ 是该系统的余极限，如引理
\ref{properties-lemma-algebra-directed-colimit-finite-presentation}，
把 $\mathcal{A} = \colim_i \mathcal{A}_i$ 写成有限表现准凝聚
$\mathcal{O}_X$-代数的有向余极限。则像
$\mathcal{A}'_i = \Im(\mathcal{A}_i \to \mathcal{A})$
是 $\mathcal{A}$ 的有限型准凝聚子代数。由于
$\mathcal{A}$ 是这些子代数的余极限，结论得证。
\end{proof}
\noindent
令 $X$ 为概形。在下述引理中，我们使用{\it 有限（分别地，整的）准凝聚
$\mathcal{O}_X$-代数 $\mathcal{A}$}这一概念。它意味着，对每个仿射开集
$\Spec(R) \subset X$，都有 $\mathcal{A} = \widetilde{A}$，其中 $A$
是一个（交换）$R$-代数，并且作为 $R$-代数是有限的（分别地，是整的）。

\begin{lemma}
\label{properties-lemma-finite-algebra-directed-colimit-finite-finitely-presented}
令 $X$ 为概形。假设 $X$ 拟紧且拟分离。
令 $\mathcal{A}$ 为有限准凝聚 $\mathcal{O}_X$-代数。
则 $\mathcal{A} = \colim \mathcal{A}_i$ 是有限且有限表现的准凝聚
$\mathcal{O}_X$-代数的有向余极限，并且所有过渡映射
$\mathcal{A}_{i'} \to \mathcal{A}_i$ 都是满射。
\end{lemma}

\begin{proof}
由引理 \ref{properties-lemma-finite-directed-colimit-surjective-maps}，存在有限表现的
$\mathcal{O}_X$-模 $\mathcal{F}$ 以及满射
$\mathcal{F} \to \mathcal{A}$。利用代数结构，得到满射
$$
\text{Sym}^*_{\mathcal{O}_X}(\mathcal{F}) \longrightarrow \mathcal{A}
$$
以 $\mathcal{J}$ 表示其核。把 $\mathcal{J} = \colim \mathcal{E}_i$
写成有限型 $\mathcal{O}_X$-子模 $\mathcal{E}_i$ 的滤过余极限
（引理 \ref{properties-lemma-quasi-coherent-colimit-finite-type}）。置
$$
\mathcal{A}_i = \text{Sym}^*_{\mathcal{O}_X}(\mathcal{F})/(\mathcal{E}_i)
$$
其中 $(\mathcal{E}_i)$ 表示由映射
$\mathcal{E}_i \to \text{Sym}^*_{\mathcal{O}_X}(\mathcal{F})$ 的像生成的理想层。
于是每个 $\mathcal{A}_i$ 都是有限表现的 $\mathcal{O}_X$-代数，
过渡映射都是满射，并且 $\mathcal{A} = \colim \mathcal{A}_i$。
为完成证明，还须说明 $\mathcal{A}_i$ 是有限
$\mathcal{O}_X$-代数，只要 $i$ 足够大。为此，选取仿射开覆盖
$X = U_1 \cup \ldots \cup U_m$。选取生成元
$f_{j, 1}, \ldots, f_{j, N_j} \in \Gamma(U_i, \mathcal{F})$。
因为 $\mathcal{A}(U_j)$ 是有限 $\mathcal{O}_X(U_j)$-代数，所以对每个
$k$，都存在首一多项式 $P_{j, k} \in \mathcal{O}(U_j)[T]$，使得
$P_{j, k}(f_{j, k})$ 在 $\mathcal{A}(U_j)$ 中为零。由于依构造有
$\mathcal{A} = \colim \mathcal{A}_i$，故有
$P_{j, k}(f_{j, k}) = 0$ 在 $\mathcal{A}_i(U_j)$ 中成立，只要 $i$ 足够大。对这样的
$i$，代数 $\mathcal{A}_i$ 是有限的。
\end{proof}

\begin{lemma}
\label{properties-lemma-integral-algebra-directed-colimit-finite}
令 $X$ 为概形。假设 $X$ 拟紧且拟分离。
令 $\mathcal{A}$ 为整的准凝聚 $\mathcal{O}_X$-代数。则
\begin{enumerate}
\item $\mathcal{A}$ 是其有限准凝聚 $\mathcal{O}_X$-子代数的有向余极限；
\item $\mathcal{A}$ 是有限且有限表现的准凝聚 $\mathcal{O}_X$-代数的直余极限。
\end{enumerate}
\end{lemma}

\begin{proof}
由引理 \ref{properties-lemma-algebra-directed-colimit-finite-type}，有
$\mathcal{A} = \colim \mathcal{A}_i$，其中
$\mathcal{A}_i \subset \mathcal{A}$ 遍历有限型准凝聚
$\mathcal{O}_X$-代数。任意有限型准凝聚 $\mathcal{O}_X$-子代数（属于
$\mathcal{A}$）都是有限的（把《代数》引理
\ref{algebra-lemma-characterize-finite-in-terms-of-integral}
应用于 $\mathcal{A}_i(U) \subset \mathcal{A}(U)$，其中 $U$ 是 $X$
中的仿射开集）。这证明了 (1)。

\medskip\noindent
为证明 (2)，利用引理 \ref{properties-lemma-directed-colimit-finite-presentation}，
把 $\mathcal{A} = \colim \mathcal{F}_i$ 写成有限表现
$\mathcal{O}_X$-模的余极限。对每个 $i$，令 $\mathcal{J}_i$ 为映射
$$
\text{Sym}^*_{\mathcal{O}_X}(\mathcal{F}_i) \longrightarrow \mathcal{A}
$$
的核。若 $i' \geq i$，则有诱导映射
$\mathcal{J}_i \to \mathcal{J}_{i'}$，并且
$\mathcal{A} =
\colim \text{Sym}^*_{\mathcal{O}_X}(\mathcal{F}_i)/\mathcal{J}_i$。
此外，准凝聚 $\mathcal{O}_X$-代数
$\text{Sym}^*_{\mathcal{O}_X}(\mathcal{F}_i)/\mathcal{J}_i$ 是有限的
（见上文）。把 $\mathcal{J}_i = \colim \mathcal{E}_{ik}$ 写成有限表现
$\mathcal{O}_X$-模的余极限。给定 $i' \geq i$ 与 $k$，存在 $k'$，使得有映射
$\mathcal{E}_{ik} \to \mathcal{E}_{i'k'}$，并使图
$$
\xymatrix{
\mathcal{J}_i \ar[r] & \mathcal{J}_{i'} \\
\mathcal{E}_{ik} \ar[u] \ar[r] & \mathcal{E}_{i'k'} \ar[u]
}
$$
交换。这由《模》引理
\ref{modules-lemma-finite-presentation-quasi-compact-colimit} 得出。
它诱导映射
$$
\mathcal{A}_{ik} =
\text{Sym}^*_{\mathcal{O}_X}(\mathcal{F}_i)/(\mathcal{E}_{ik})
\longrightarrow
\text{Sym}^*_{\mathcal{O}_X}(\mathcal{F}_{i'})/(\mathcal{E}_{i'k'}) =
\mathcal{A}_{i'k'}
$$
其中 $(\mathcal{E}_{ik})$ 表示由 $\mathcal{E}_{ik}$ 生成的理想。
准凝聚 $\mathcal{O}_X$-代数 $\mathcal{A}_{ki}$ 当 $k$ 足够大时
既有限表现又有限（见引理
\ref{properties-lemma-finite-algebra-directed-colimit-finite-finitely-presented} 的证明）。
最后，有
$$
\colim \mathcal{A}_{ik} = \colim \mathcal{A}_i = \mathcal{A}
$$
事实上，第一个等式已在引理
\ref{properties-lemma-finite-algebra-directed-colimit-finite-finitely-presented}
的证明中说明；第二个等式成立，是因为 $\mathcal{A}$ 是模
$\mathcal{F}_i$ 的余极限。
\end{proof}





\section{Gabber 的结果}
\label{properties-section-gabber}

\noindent
本节证明 Gabber 的一个结果：对每个概形，都存在一个基数 $\kappa$，
使得每个准凝聚模 $\mathcal{F}$ 都是它的 $\kappa$-生成准凝聚子层之并。
由此可知，概形上的准凝聚层范畴是具有极限和充分多内射对象的
Grothendieck 阿贝尔范畴\footnote{Akhil Mathew 在一篇
\href{https://amathew.wordpress.com/2011/07/30/quasi-coherent-sheaves-presentable-categories-and-a-result-of-gabber/}{博客文章}
中给出了很好的解释。}。

\begin{definition}
\label{properties-definition-kappa-generated}
令 $(X, \mathcal{O}_X)$ 为环层空间，令 $\kappa$ 为无限基数。
称 $\mathcal{O}_X$-模层 $\mathcal{F}$ 是{\it $\kappa$-生成的}，
如果存在开覆盖 $X = \bigcup U_i$，使得 $\mathcal{F}|_{U_i}$ 由子集
$R_i \subset \mathcal{F}(U_i)$ 生成，并且该子集的基数至多为
$\kappa$。
\end{definition}

\noindent
注意，至多 $\kappa$ 个 $\kappa$-生成模的直和仍是
$\kappa$-生成的，因为 $\kappa \otimes \kappa = \kappa$；见《集合》第
\ref{sets-section-cardinals} 节。
特别地，两个 $\kappa$-生成模的直和具有此性质。
此外，$\kappa$-生成层的商层仍是 $\kappa$-生成的。
（但对子模未必如此。）

\begin{lemma}
\label{properties-lemma-set-of-iso-classes}
令 $(X, \mathcal{O}_X)$ 为环层空间，令 $\kappa$ 为基数。
则存在集合 $T$ 以及族 $(\mathcal{F}_t)_{t \in T}$，其成员是
$\kappa$-生成 $\mathcal{O}_X$-模，使得每个 $\kappa$-生成
$\mathcal{O}_X$-模都同构于某个 $\mathcal{F}_t$。
\end{lemma}

\begin{proof}
$X$ 的覆盖组成一个集合（只要不允许重复）。
假设 $X = \bigcup U_i$ 是一个覆盖，并且 $\mathcal{F}_i$
是 $\mathcal{O}_{U_i}$-模。则具有性质
$\mathcal{O}_X$-模 $\mathcal{F}$ 的同构类组成一个集合，其中
$\mathcal{F}|_{U_i} \cong \mathcal{F}_i$，因为粘合映射组成一个集合。
这样，问题化为证明：商映射
$\oplus_{k \in \kappa} \mathcal{O}_X \to \mathcal{F}$ 的（同构类）
对任意环层空间 $X$ 都组成一个集合。这是显然的。
\end{proof}

\noindent
下面就是本节标题所指的结果。

\begin{lemma}
\label{properties-lemma-colimit-kappa}
令 $X$ 为概形。则存在一个基数 $\kappa$，使得每个准凝聚模
$\mathcal{F}$ 都是它的 $\kappa$-生成准凝聚子模的有向余极限。
\end{lemma}

\begin{proof}
选取仿射开覆盖 $X = \bigcup_{i \in I} U_i$。对每一对
$i, j$，选取仿射开覆盖
$U_i \cap U_j = \bigcup_{k \in I_{ij}} U_{ijk}$。
写成 $U_i = \Spec(A_i)$ 和 $U_{ijk} = \Spec(A_{ijk})$。
令 $\kappa$ 为任意无限基数，并且它 $\geq$ 各集合
$I$、$I_{ij}$ 的基数。

\medskip\noindent
令 $\mathcal{F}$ 为准凝聚层。置 $M_i = \mathcal{F}(U_i)$ 和
$M_{ijk} = \mathcal{F}(U_{ijk})$。注意
$$
M_i \otimes_{A_i} A_{ijk} = M_{ijk} = M_j \otimes_{A_j} A_{ijk}.
$$
见《概形》引理 \ref{schemes-lemma-widetilde-pullback}。
利用选择公理，选取映射
$$
(i, j, k, m) \mapsto S(i, j, k, m)
$$
它对每个 $i, j \in I$、$k \in I_{ij}$ 和 $m \in M_i$，
指定有限子集 $S(i, j, k, m) \subset M_j$，使得
$$
m \otimes 1 = \sum\nolimits_{m' \in S(i, j, k, m)} m' \otimes a_{m'}
$$
在 $M_{ijk}$ 中成立，其中某些 $a_{m'} \in A_{ijk}$。
此外，约定 $S(i, i, k, m) = \{m\}$ 对上面所有
$i, j = i, k, m$ 都成立。固定这样的映射。

\medskip\noindent
给定子集族 $\mathcal{S} = (S_i)_{i \in I}$，其中
$S_i \subset M_i$ 的基数至多为 $\kappa$，置
$\mathcal{S}' = (S'_i)$，其中
$$
S'_j = \bigcup\nolimits_{(i, k, m)\text{ 使得 }m \in S_i}
S(i, j, k, m)
$$
注意 $S_i \subset S'_i$。又注意，$S'_i$ 的基数至多为
$\kappa$，因为它是以一个基数至多为 $\kappa$ 的集合为指标的有限集之并。
置 $\mathcal{S}^{(0)} = \mathcal{S}$、
$\mathcal{S}^{(1)} = \mathcal{S}'$，并归纳地置
$\mathcal{S}^{(n + 1)} = (\mathcal{S}^{(n)})'$。再置
$\mathcal{S}^{(\infty)} = \bigcup_{n \geq 0} \mathcal{S}^{(n)}$。
写 $\mathcal{S}^{(\infty)} = (S^{(\infty)}_i)$，则对任意元素
$m \in S^{(\infty)}_i$，$m$ 在 $M_{ijk}$ 中的像都可写成有限和
$\sum m' \otimes a_{m'}$，其中 $m' \in S_j^{(\infty)}$。
由此可见，若置
$$
N_i = A_i\text{-子模，属于 }M_i\text{ 由下列对象生成： }S^{(\infty)}_i
$$
则有
$$
N_i \otimes_{A_i} A_{ijk} = N_j \otimes_{A_j} A_{ijk}.
$$
把两边视为 $M_{ijk}$ 的子模。因此，存在准凝聚子层
$\mathcal{G} \subset \mathcal{F}$，使得 $\mathcal{G}(U_i) = N_i$。
此外，依构造，层 $\mathcal{G}$ 是 $\kappa$-生成的。

\medskip\noindent
令 $\{\mathcal{G}_t\}_{t \in T}$ 为所有 $\kappa$-生成准凝聚子层组成的集合。
若 $t, t' \in T$，则 $\mathcal{G}_t + \mathcal{G}_{t'}$
也是 $\kappa$-生成准凝聚子层，因为它是映射
$\mathcal{G}_t \oplus \mathcal{G}_{t'} \to \mathcal{F}$ 的像。
因此，该系统（按包含关系排序）是有向的。
上述论证表明，$\mathcal{F}$ 在 $U_i$ 上的每个截面都属于某个
$\mathcal{G}_t$（因为可以从 $\mathcal{S}$ 开始，并使给定截面属于
$S_i$）。所以
$\colim_t \mathcal{G}_t \to \mathcal{F}$ 既单又满，正是所需。
\end{proof}

\begin{proposition}
\label{properties-proposition-coherator}
令 $X$ 为概形。
\begin{enumerate}
\item 范畴 $\QCoh(\mathcal{O}_X)$ 是 Grothendieck 阿贝尔范畴。
因此，$\QCoh(\mathcal{O}_X)$ 具有充分多内射对象和所有极限。
\item 包含函子
$\QCoh(\mathcal{O}_X) \to \textit{Mod}(\mathcal{O}_X)$
具有右伴随\footnote{这个函子有时称为{\it 凝聚化函子}。}
$$
Q : \textit{Mod}(\mathcal{O}_X) \longrightarrow \QCoh(\mathcal{O}_X)
$$
使得对每个准凝聚层 $\mathcal{F}$，伴随映射
$Q(\mathcal{F}) \to \mathcal{F}$ 都是同构。
\end{enumerate}
\end{proposition}

\begin{proof}
(1) 意味着 $\QCoh(\mathcal{O}_X)$ 满足：(a) 具有所有余极限，
(b) 滤过余极限正合，(c) 具有生成对象；见《内射对象》第
\ref{injectives-section-grothendieck-conditions} 节。
由《概形》第 \ref{schemes-section-quasi-coherent} 节，
$\QCoh(\mathcal{O}_X)$ 中的余极限存在，并且与
$\textit{Mod}(\mathcal{O}_X)$ 中的余极限一致。由《模》引理
\ref{modules-lemma-limits-colimits}，滤过余极限正合。
因此 (a) 和 (b) 成立。为构造生成对象 $U$，选取引理
\ref{properties-lemma-colimit-kappa} 中的基数 $\kappa$。选取引理
\ref{properties-lemma-set-of-iso-classes} 中的族
$(\mathcal{F}_t)_{t \in T}$，它由 $\kappa$-生成准凝聚层组成。置
$U = \bigoplus_{t \in T} \mathcal{F}_t$。由于
$\QCoh(\mathcal{O}_X)$ 的每个对象都是 $\kappa$-生成准凝聚模的滤过余极限，
即同构于 $\mathcal{F}_t$ 的对象的滤过余极限，显然 $U$ 是生成对象。
极限和内射对象的断言在任意 Grothendieck 阿贝尔范畴中都成立；见《内射对象》
定理 \ref{injectives-theorem-injective-embedding-grothendieck} 和引理
\ref{injectives-lemma-grothendieck-products}。

\medskip\noindent
证明 (2)。为构造 $Q$，使用下述一般步骤。
给定对象 $\mathcal{F}$，它属于 $\textit{Mod}(\mathcal{O}_X)$，考虑函子
$$
\QCoh(\mathcal{O}_X)^{opp} \longrightarrow \textit{Sets},\quad
\mathcal{G} \longmapsto \Hom_X(\mathcal{G}, \mathcal{F})
$$
该函子把余极限变为极限，因而可表；见《内射对象》引理
\ref{injectives-lemma-grothendieck-brown}。
因此存在准凝聚层 $Q(\mathcal{F})$ 以及函子性同构
$\Hom_X(\mathcal{G}, \mathcal{F}) = \Hom_X(\mathcal{G}, Q(\mathcal{F}))$，
其中 $\mathcal{G}$ 属于 $\QCoh(\mathcal{O}_X)$。由 Yoneda 引理
（《范畴》引理 \ref{categories-lemma-yoneda}），构造
$\mathcal{F} \leadsto Q(\mathcal{F})$ 关于 $\mathcal{F}$ 是函子性的。
依构造，$Q$ 是包含函子的右伴随。
映射 $Q(\mathcal{F}) \to \mathcal{F}$ 是同构，只要 $\mathcal{F}$ 准凝聚；
这是包含函子
$\QCoh(\mathcal{O}_X) \to \textit{Mod}(\mathcal{O}_X)$
为全忠实函子这一事实的形式推论。
\end{proof}





\section{支集包含于闭子集的截面}
\label{properties-section-sections-with-support-in-closed}

\noindent
给定任意拓扑空间 $X$、闭子集 $Z \subset X$ 和阿贝尔层
$\mathcal{F}$，可以取由支集包含于 $Z$ 的截面组成的子层。
若 $X$ 为概形，$Z$ 为闭子概形，而 $\mathcal{F}$ 为准凝聚模，
则还有一种变体：取概形论意义下支集位于 $Z$ 上的截面。
不过，在概形情形中必须谨慎，因为所得 $\mathcal{O}_X$-模未必准凝聚。

\begin{lemma}
\label{properties-lemma-quasi-coherent-finite-type-ideals}
令 $X$ 为拟紧且拟分离概形。
令 $U \subset X$ 为开子概形。下列条件等价：
\begin{enumerate}
\item $U$ 在 $X$ 中逆拟紧；
\item $U$ 拟紧；
\item $U$ 是有限多个仿射开集之并；
\item 存在有限型准凝聚理想层
$\mathcal{I} \subset \mathcal{O}_X$，使得 $X \setminus U = V(\mathcal{I})$
（作为集合）。
\end{enumerate}
\end{lemma}

\begin{proof}
(1)、(2)、(3) 的等价性由引理
\ref{properties-lemma-quasi-separated-quasi-compact-open-retrocompact} 得出。
假设 (1)、(2)、(3) 成立。令 $T = X \setminus U$。由《概形》引理
\ref{schemes-lemma-reduced-closed-subscheme}，存在唯一准凝聚理想层
$\mathcal{J}$，它在 $T$ 上切出约化诱导闭子概形结构。
注意 $\mathcal{J}|_U = \mathcal{O}_U$，而后者是有限型
$\mathcal{O}_U$-模。由引理 \ref{properties-lemma-extend}，存在有限型准凝聚子层
$\mathcal{I} \subset \mathcal{J}$，并且它满足
$\mathcal{I}|_U = \mathcal{J}|_U$。于是
$X \setminus U = V(\mathcal{I})$，从而得到 (4)。反之，若
$\mathcal{I}$ 如 (4)，并且 $W = \Spec(R) \subset X$ 为仿射开集，则
$\mathcal{I}|_W = \widetilde{I}$，其中某个有限生成理想
$I \subset R$；见引理 \ref{properties-lemma-finite-type-module}。
因此 $U \cap W = \Spec(R) \setminus V(I)$ 拟紧；见《代数》引理
\ref{algebra-lemma-qc-open}。于是 $U \subset X$ 由引理
\ref{properties-lemma-retrocompact} 可知是逆拟紧的。
\end{proof}

\begin{lemma}
\label{properties-lemma-sections-annihilated-by-ideal}
令 $X$ 为概形。
令 $\mathcal{I} \subset \mathcal{O}_X$ 为准凝聚理想层。
令 $\mathcal{F}$ 为准凝聚 $\mathcal{O}_X$-模。
考虑 $\mathcal{O}_X$-模层 $\mathcal{F}'$，它对每个开集
$U \subset X$ 赋予
$$
\mathcal{F}'(U)
=
\{s \in \mathcal{F}(U) \mid
\mathcal{I}s = 0\}
$$
假设 $\mathcal{I}$ 有限型。则
\begin{enumerate}
\item $\mathcal{F}'$ 是准凝聚 $\mathcal{O}_X$-模层；
\item 对任意仿射开集 $U \subset X$，有
$\mathcal{F}'(U) = \{s \in \mathcal{F}(U) \mid \mathcal{I}(U)s = 0\}$；
\item $\mathcal{F}'_x = \{s \in \mathcal{F}_x \mid \mathcal{I}_x s = 0\}$。
\end{enumerate}
\end{lemma}

\begin{proof}
显然，定义 $\mathcal{F}'$ 的规则给出 $\mathcal{F}$ 的子层
（层条件容易检验）。因此可以在 $X$ 上局部地验证其他断言。
换言之，可以假设
$X = \Spec(A)$、$\mathcal{F} = \widetilde{M}$ 且
$\mathcal{I} = \widetilde{I}$。在此情形下，显然
$\mathcal{F}'(U) = \{x \in M \mid Ix = 0\} =: M'$，因为
$\widetilde{I}$ 由它的整体截面 $I$ 生成；这证明了 (2)。
为证明 $\mathcal{F}'$ 准凝聚，只须说明对每个 $f \in A$ 都有
$\{x \in M_f \mid I_f x = 0\} = (M')_f$。
写 $I = (g_1, \ldots, g_t)$；这是可行的，因为 $\mathcal{I}$ 有限型；
见引理 \ref{properties-lemma-finite-type-module}。
若 $x = y/f^n$ 且 $I_fx = 0$，则这意味着对每个 $i$ 都存在
$m \geq 0$，使得 $f^mg_ix = 0$。
可以选取一个 $m$，它对所有 $i$ 都适用（这里正用到了 $I$ 有限生成）。
于是 $f^mx \in M'$，并且
$x/f^n = f^mx/f^{n + m}$ 在 $(M')_f$ 中成立，正是所需。
(3) 的证明类似，故略去。
\end{proof}

\begin{definition}
\label{properties-definition-subsheaf-sections-annihilated-by-ideal}
令 $X$ 为概形。
令 $\mathcal{I} \subset \mathcal{O}_X$ 为有限型准凝聚理想层。
令 $\mathcal{F}$ 为准凝聚 $\mathcal{O}_X$-模。
把上面引理 \ref{properties-lemma-sections-annihilated-by-ideal} 定义的子层
$\mathcal{F}' \subset \mathcal{F}$ 称为{\it 被 $\mathcal{I}$ 零化的截面子层}。
\end{definition}

\begin{lemma}
\label{properties-lemma-push-sections-annihilated-by-ideal}
令 $f : X \to Y$ 为概形的拟紧拟分离态射。
令 $\mathcal{I} \subset \mathcal{O}_Y$ 为有限型准凝聚理想层。
令 $\mathcal{F}$ 为准凝聚 $\mathcal{O}_X$-模。令
$\mathcal{F}' \subset \mathcal{F}$ 为被
$f^{-1}\mathcal{I}\mathcal{O}_X$ 零化的截面子层。
则 $f_*\mathcal{F}' \subset f_*\mathcal{F}$ 是被
$\mathcal{I}$ 零化的截面子层。
\end{lemma}

\begin{proof}
略。（提示：$f$ 拟紧且拟分离这一假设蕴含
$f_*\mathcal{F}$ 准凝聚，因而引理
\ref{properties-lemma-sections-annihilated-by-ideal} 可应用于
$\mathcal{I}$ 和 $f_*\mathcal{F}$。）
\end{proof}

\noindent
对于拓扑空间上的阿贝尔层，我们已在《模》评注
\ref{modules-remark-sections-support-in-closed} 中讨论过
支集包含于闭子集的截面子层。对于准凝聚模，这个子模未必准凝聚；
但若该闭子集的补集逆拟紧，则它准凝聚。

\begin{lemma}
\label{properties-lemma-sections-supported-on-closed-subset}
令 $X$ 为概形。令 $Z \subset X$ 为闭子集。
令 $\mathcal{F}$ 为准凝聚 $\mathcal{O}_X$-模。
考虑 $\mathcal{O}_X$-模层 $\mathcal{F}'$，它对每个开集
$U \subset X$ 赋予
$$
\mathcal{F}'(U)
=
\{s \in \mathcal{F}(U) \mid
\text{下列对象的支集： }s\text{ 包含于 }Z \cap U\}
$$
若 $X \setminus Z$ 是 $X$ 的逆拟紧开集，则
\begin{enumerate}
\item 对仿射开集 $U \subset X$，存在有限生成理想
$I \subset \mathcal{O}_X(U)$，使得 $Z \cap U = V(I)$；
\item 对 (1) 中的 $U$ 和 $I$，有
$\mathcal{F}'(U) = \{x \in \mathcal{F}(U) \mid
I^nx = 0 \text{ 对某个 } n\}$；
\item $\mathcal{F}'$ 是准凝聚 $\mathcal{O}_X$-模层。
\end{enumerate}
\end{lemma}

\begin{proof}
(1) 是《代数》引理 \ref{algebra-lemma-qc-open}。
令 $U = \Spec(A)$，并令 $I$ 如 (1)。
则 $\mathcal{F}|_U$ 是与某个 $A$-模 $M$ 相伴的准凝聚层。我们有
$$
\mathcal{F}'(U) = \{x \in M \mid x = 0\text{ 于 }M_\mathfrak p
\text{ 对所有 }\mathfrak p \not \in Z\}.
$$
这是由《模》定义 \ref{modules-definition-support} 得出的。因此
$x \in \mathcal{F}'(U)$ 当且仅当
$V(\text{Ann}(x)) \subset V(I)$；见《代数》引理
\ref{algebra-lemma-support-element}。由于 $I$ 有限生成，这等价于
$I^n x = 0$ 对某个 $n$ 成立。这证明了 (2)。

\medskip\noindent
证明 (3)。注意，对给定开集 $U \subset X$，有正合序列
$$
0 \to \mathcal{F}'(U) \to \mathcal{F}(U) \to \mathcal{F}(U \setminus Z)
$$
若以 $j : X \setminus Z \to X$ 表示包含态射，则注意
$\mathcal{F}(U \setminus Z)$ 是模
$j_*(\mathcal{F}|_{X \setminus Z})$ 在 $U$ 上的截面。因此有正合序列
$$
0 \to \mathcal{F}' \to \mathcal{F} \to j_*(\mathcal{F}|_{X \setminus Z})
$$
限制 $\mathcal{F}|_{X \setminus Z}$ 是准凝聚的。
所以 $j_*(\mathcal{F}|_{X \setminus Z})$ 准凝聚；这是由《概形》引理
\ref{schemes-lemma-push-forward-quasi-coherent}
及关于 $j$ 拟紧的假设（任意开浸入都分离）得出的。
于是 $\mathcal{F}'$ 作为准凝聚模间映射的核是准凝聚的；见《概形》第
\ref{schemes-section-quasi-coherent} 节。
\end{proof}

\begin{definition}
\label{properties-definition-subsheaf-sections-supported-on-closed}
令 $X$ 为概形。
令 $T \subset X$ 为闭子集，并使其补集在 $X$ 中逆拟紧。
令 $\mathcal{F}$ 为准凝聚 $\mathcal{O}_X$-模。
把引理 \ref{properties-lemma-sections-supported-on-closed-subset} 定义的准凝聚子层
$\mathcal{F}' \subset \mathcal{F}$ 称为{\it 支集位于 $T$ 上的截面子层}。
\end{definition}

\begin{lemma}
\label{properties-lemma-push-sections-supported-on-closed-subset}
令 $f : X \to Y$ 为概形的拟紧拟分离态射。
令 $Z \subset Y$ 为闭子集，使得 $Y \setminus Z$ 在 $Y$ 中逆拟紧。
令 $\mathcal{F}$ 为准凝聚 $\mathcal{O}_X$-模。令
$\mathcal{F}' \subset \mathcal{F}$ 为支集位于 $f^{-1}Z$ 上的截面子层。
则 $f_*\mathcal{F}' \subset f_*\mathcal{F}$ 是支集位于
$Z$ 上的截面子层。
\end{lemma}

\begin{proof}
略。（提示：先证明 $X \setminus f^{-1}Z$ 在 $X$ 中逆拟紧，因为
$Y \setminus Z$ 在 $Y$ 中逆拟紧。因此引理
\ref{properties-lemma-sections-supported-on-closed-subset} 可应用于
$f^{-1}Z$ 和 $\mathcal{F}$。由于 $f$ 拟紧且拟分离，可知
$f_*\mathcal{F}$ 准凝聚。因此引理
\ref{properties-lemma-sections-supported-on-closed-subset} 可应用于
$Z$ 和 $f_*\mathcal{F}$。最后，直接对照这些层。）
\end{proof}





\section{准凝聚层的截面}
\label{properties-section-sections-quasi-coherent}

\noindent
下面计算仿射谱的拟紧开集上准凝聚层的截面。

\begin{lemma}
\label{properties-lemma-sections-over-quasi-compact-open-in-affine}
令 $A$ 为环。
令 $I \subset A$ 为有限生成理想。
令 $M$ 为 $A$-模。
则存在典范映射
$$
\colim_n \Hom_A(I^n, M)
\longrightarrow
\Gamma(\Spec(A) \setminus V(I), \widetilde{M}).
$$
这个映射总是单射。
若对所有 $x \in M$ 都有 $Ix = 0 \Rightarrow x = 0$，
则该映射是同构。一般地，置
$M_n = \{x \in M \mid I^nx = 0\}$，则存在同构
$$
\colim_n \Hom_A(I^n, M/M_n)
\longrightarrow
\Gamma(\Spec(A) \setminus V(I), \widetilde{M}).
$$
\end{lemma}

\begin{proof}
由于 $I^{n + 1} \subset I^n$ 且 $M_n \subset M_{n + 1}$，
可以通过与这些映射复合，得到典范 $A$-模映射
$$
\Hom_A(I^n, M)
\longrightarrow
\Hom_A(I^{n + 1}, M)
$$
以及
$$
\Hom_A(I^n, M/M_n)
\longrightarrow
\Hom_A(I^{n + 1}, M/M_{n + 1})
$$
并把它们用作这些系统中的过渡映射。给定
$A$-模映射 $\varphi : I^n \to M$，得到层映射
$\widetilde{\varphi} : \widetilde{I^n} \to \widetilde{M}$，
可以把它限制到开集 $\Spec(A) \setminus V(I)$。
由于 $\widetilde{I^n}$ 限制到该开集后给出结构层，便得到
$\Gamma(\Spec(A) \setminus V(I), \widetilde{M})$ 的一个元素。
略去它与系统 $\Hom_A(I^n, M)$ 中过渡映射相容的验证。
这给出第一条箭头。为得到第二条箭头，注意
$\widetilde{M}$ 和 $\widetilde{M/M_n}$ 在开集
$\Spec(A) \setminus V(I)$ 上一致，因为层 $\widetilde{M_n}$
显然支集位于 $V(I)$ 上。因此可以使用与前面相同的机制。

\medskip\noindent
下面用代数语言具体说明如何定义这条箭头。
设 $I = (f_1, \ldots, f_t)$。则
$\Spec(A) \setminus V(I) = \bigcup_{i = 1, \ldots, t} D(f_i)$。
因此
$$
0 \to
\Gamma(\Spec(A) \setminus V(I), \widetilde{M}) \to
\bigoplus\nolimits_i M_{f_i} \to
\bigoplus\nolimits_{i, j} M_{f_if_j}
$$
正合。假设 $\varphi : I^n \to M$ 是 $A$-模映射。
考虑元素向量 $\varphi(f_i^n)/f_i^n \in M_{f_i}$。
容易看出，这个向量在上面正合序列的第二个直和中映为零。
因而得到 $\Gamma(\Spec(A) \setminus V(I), \widetilde{M})$ 的一个元素。
略去此描述与上面描述一致的验证。

\medskip\noindent
利用这一描述来证明第一条箭头是单射。
事实上，若 $\varphi$ 映为零，则对每个 $i$，元素
$\varphi(f_i^n)/f_i^n$ 在 $M_{f_i}$ 中为零。换言之，对每个
$i$，都有 $f_i^m\varphi(f_i^n) = 0$ 对某个 $m \geq 0$ 成立。
可以选取一个 $m$，它对所有 $i$ 都适用。
于是 $\varphi(f_i^{n + m}) = 0$ 对所有 $i$ 都成立。
容易看出，这意味着
$\varphi|_{I^{t(n + m - 1) + 1}} = 0$；换言之，$\varphi$
在余极限的第 $t(n + m - 1) + 1$ 项中映为零。
单射性由此得出。

\medskip\noindent
注意，每个 $M_n = 0$，只要 $Ix = 0 \Rightarrow x = 0$
对 $x \in M$ 成立。因此，为完成本引理的证明，只须说明第二条箭头是同构。

\medskip\noindent
尝试构造本引理第二个映射的逆。令
$s \in \Gamma(\Spec(A) \setminus V(I), \widetilde{M})$。
它对应于上面正合序列第一个直和中的向量
$x_i/f_i^n$，其中 $x_i \in M$。
因此，对每个 $i, j$，都存在 $m \geq 0$，使得
$f_i^m f_j^m (f_j^n x_i - f_i^n x_j) = 0$ 在 $M$ 中成立。
可以选取一个 $m$，它对所有有序对 $i, j$ 都适用。
把 $x_i$ 替换为 $f_i^mx_i$，并把 $n$ 替换为 $n + m$ 后，
可知得到 $f_j^nx_i = f_i^nx_j$，它在 $M$ 中对所有 $i, j$ 都成立。
引入
$$
K_n = \{x \in M \mid f_1^nx = \ldots = f_t^nx = 0\}
$$
我们断言存在 $A$-模映射
$$
\varphi :
I^{t(n - 1) + 1}
\longrightarrow
M/K_n
$$
它把单项式
$f_1^{e_1} \ldots f_t^{e_t}$（其中 $\sum e_i = t(n - 1) + 1$）
映到模 $K_n$ 的类，该类由表达式
$f_1^{e_1} \ldots f_i^{e_i - n} \ldots f_t^{e_t}x_i$
给出，其中选取 $i$ 使得 $e_i \geq n$
（注意，至少存在一个这样的 $i$）。
为说明这一点，假设
$$
\sum\nolimits_{E = (e_1, \ldots, e_t), |E| = t(n - 1) + 1}
a_E f_1^{e_1} \ldots f_t^{e_t} = 0
$$
是这些单项式间的关系，其系数 $a_E$ 属于 $A$。
则会把它映到
$$
z =
\sum\nolimits_{E = (e_1, \ldots, e_t), |E| = t(n - 1) + 1}
a_E f_1^{e_1} \ldots f_{i(E)}^{e_{i(E)} - n} \ldots f_t^{e_t}x_{i(E)}
$$
其中对每个多重指标 $E$，选定一个特定的 $i(E)$，使得
$e_{i(E)} \geq n$。
注意，若把它乘以 $f_j^n$，其中 $j$ 任意，便得到零，因为利用上面的关系
$f_j^nx_i = f_i^nx_j$ 可得
\begin{align*}
f_j^nz & = \sum\nolimits_{E = (e_1, \ldots, e_t), |E| = t(n - 1) + 1}
a_E f_1^{e_1} \ldots f_j^{e_j + n}
\ldots f_{i(E)}^{e_{i(E)} - n} \ldots f_t^{e_t}x_{i(E)} \\
& =
\sum\nolimits_{E = (e_1, \ldots, e_t), |E| = t(n - 1) + 1}
a_E f_1^{e_1} \ldots f_t^{e_t}x_j
= 0.
\end{align*}
因此 $z \in K_n$，从而每个关系在 $M/K_n$ 中都映为零。
这证明了断言。

\medskip\noindent
注意 $K_n \subset M_{t(n - 1) + 1}$。因此映射 $\varphi$
特别给出 $A$-模映射
$I^{t(n - 1) + 1} \to M/M_{t(n - 1) + 1}$。
这证明了本引理的第二条箭头是满射。略去单射性的证明。
\end{proof}

\begin{example}
\label{properties-example-does-not-work-in-general}
下面给出两个例子，说明引理
\ref{properties-lemma-sections-over-quasi-compact-open-in-affine}
的第一条陈列映射并非同构。

\medskip\noindent
令 $k$ 为域。考虑环
$$
A = k[x, y, z_1, z_2, \ldots]/(x^nz_n).
$$
置 $I = (x)$ 并令 $M = A$。则元素 $y/x$ 定义了 $\Spec(A)$
的结构层在 $D(x) = \Spec(A) \setminus V(I)$ 上的一个截面。
我们断言，$y/x$ 不在典范映射
$\colim \Hom_A(I^n, A) \to A_x = \mathcal{O}(D(x))$ 的像中。
事实上，若在像中，则它会来自同态
$\varphi : I^n \to A$，其中某个 $n$。置 $a = \varphi(x^n)$。
于是会有 $x^m(xa - x^ny) = 0$，其中某个 $m > 0$。
这意味着 $x^{m + 1}a = x^{m + n}y$，继而意味着
$\varphi(x^{n + m + 1}) = x^{m + n}y$。这导致矛盾，因为它将蕴含
$$
0 = \varphi(0) = \varphi(z_{n + m + 1} x^{n + m + 1}) =
x^{m + n}y z_{n + m + 1}
$$
而这在环 $A$ 中不成立。

\medskip\noindent
令 $k$ 为域。考虑环
$$
A = k[f, g, x, y, \{a_n, b_n\}_{n \geq 1}]/
(fy - gx, \{a_nf^n + b_ng^n\}_{n \geq 1}).
$$
置 $I = (f, g)$ 并令 $M = A$。则 $x/f \in A_f$ 和
$y/g \in A_g$ 在 $A_{fg}$ 中映到同一个元素。因此，它们定义了截面
$s$；这是 $\Spec(A)$ 的结构层在
$D(f) \cup D(g) = \Spec(A) \setminus V(I)$ 上的截面。
但是，不存在 $n \geq 0$，使得 $s$ 来自如引理
\ref{properties-lemma-sections-over-quasi-compact-open-in-affine}
第一条陈列箭头之源中的 $A$-模映射
$\varphi : I^n \to A$。事实上，给定这样的模映射，置
$x_n = \varphi(f^n)$ 和 $y_n = \varphi(g^n)$。
则有 $f^mx_n = f^{n + m - 1}x$ 以及
$g^my_n = g^{n + m - 1}y$，其中某个 $m \geq 0$
（见本引理的证明）。但此时会有
$0 = \varphi(0) =
\varphi(a_{n + m}f^{n + m} + b_{n + m}g^{n + m}) =
a_{n + m}f^{n + m - 1}x + b_{n + m}g^{n + m - 1}y$，
而这在环 $A$ 中不成立。
\end{example}

\noindent
在 Noether 情形中，我们将改进下述引理；见《概形上同调》引理
\ref{coherent-lemma-homs-over-open}。

\begin{lemma}
\label{properties-lemma-sections-over-quasi-compact-open}
令 $X$ 为拟紧概形。
令 $\mathcal{I} \subset \mathcal{O}_X$ 为有限型准凝聚理想层。
令 $Z \subset X$ 为由 $\mathcal{I}$ 定义的闭子概形，并置
$U = X \setminus Z$。
令 $\mathcal{F}$ 为准凝聚 $\mathcal{O}_X$-模。
典范映射
$$
\colim_n \Hom_{\mathcal{O}_X}(\mathcal{I}^n,
\mathcal{F})
\longrightarrow
\Gamma(U, \mathcal{F})
$$
是单射。进一步假设 $X$ 拟分离。
令 $\mathcal{F}_n \subset \mathcal{F}$ 为被
$\mathcal{I}^n$ 零化的截面子层。典范映射
$$
\colim_n \Hom_{\mathcal{O}_X}(\mathcal{I}^n,
\mathcal{F}/\mathcal{F}_n)
\longrightarrow
\Gamma(U, \mathcal{F})
$$
是同构。
\end{lemma}

\begin{proof}
令 $\Spec(A) = W \subset X$ 为仿射开集。写
$\mathcal{F}|_W = \widetilde{M}$，其中某个 $A$-模 $M$，并写
$\mathcal{I}|_W = \widetilde{I}$，其中某个有限型理想
$I \subset A$。把本引理的第一条陈列映射限制到 $W$，得到引理
\ref{properties-lemma-sections-over-quasi-compact-open-in-affine}
的第一条陈列映射。由于可以用有限多个仿射开集覆盖 $X$，
这证明了本引理的第一条陈列映射是单射。

\medskip\noindent
有 $\mathcal{F}_n|_W = \widetilde{M_n}$，其中
$M_n \subset M$ 如引理
\ref{properties-lemma-sections-over-quasi-compact-open-in-affine}
中定义（略去细节）。该引理保证存在双射
$$
\colim_n  \Hom_{\mathcal{O}_W}(
\mathcal{I}^n|_W, (\mathcal{F}/\mathcal{F}_n)|_W)
\longrightarrow
\Gamma(U \cap W, \mathcal{F})
$$
这对任意这样的仿射开集 $W$ 成立。

\medskip\noindent
为证明本引理的第二条陈列箭头是双射，选取有限仿射开覆盖
$X = \bigcup_{j = 1, \ldots, m} W_j$。单射性立即由上文和覆盖的有限性得出。
若 $X$ 拟分离，则对每一对 $j, j'$，选取有限仿射开覆盖
$$
W_j \cap W_{j'} = \bigcup\nolimits_{k = 1, \ldots, m_{jj'}} W_{jj'k}.
$$
令 $s \in \Gamma(U, \mathcal{F})$。如上所见，对每个 $j$，存在
$n_j$ 以及映射
$\varphi_j : \mathcal{I}^{n_j}|_{W_j} \to
(\mathcal{F}/\mathcal{F}_{n_j})|_{W_j}$，
它对应于 $s|_{U \cap W_j}$。
同理，对每个三元组 $(j, j', k)$，存在整数 $n_{jj'k}$，
使得 $\varphi_j$ 与 $\varphi_{j'}$ 作为映射
$\mathcal{I}^{n_{jj'k}} \to \mathcal{F}/\mathcal{F}_{n_{jj'k}}$
限制到 $W_{jj'k}$ 上时一致。令
$n = \max\{n_j, n_{jj'k}\}$，则 $\varphi_j$ 可粘合为映射
$\mathcal{I}^n \to \mathcal{F}/\mathcal{F}_n$，该映射定义在 $X$ 上。
这证明了该映射的满射性。
\end{proof}





\section{丰沛可逆层}
\label{properties-section-ample}

\noindent
回顾《模》引理 \ref{modules-lemma-s-open}：给定可逆层
$\mathcal{L}$，它在局部环层空间 $X$ 上；并给定整体截面
$s$，它属于 $\mathcal{L}$，则集合
$X_s = \{x \in X \mid s \not \in \mathfrak m_x\mathcal{L}_x\}$
是开的。一般地，有
$X_s \cap X_{s'} = X_{ss'}$，其中 $ss'$ 表示截面
$s \otimes s' \in \Gamma(X, \mathcal{L} \otimes \mathcal{L}')$。

\begin{definition}
\label{properties-definition-ample}
\begin{reference}
\cite[II Definition 4.5.3]{EGA}
\end{reference}
令 $X$ 为概形。
令 $\mathcal{L}$ 为可逆 $\mathcal{O}_X$-模。称 $\mathcal{L}$ {\it 丰沛}，若
\begin{enumerate}
\item $X$ 拟紧；
\item 对每个 $x \in X$，都存在 $n \geq 1$ 和
$s \in \Gamma(X, \mathcal{L}^{\otimes n})$，使得
$x \in X_s$ 且 $X_s$ 仿射，
\end{enumerate}
\end{definition}

\begin{lemma}
\label{properties-lemma-ample-power-ample}
\begin{reference}
\cite[II Proposition 4.5.6(i)]{EGA}
\end{reference}
令 $X$ 为概形。令 $\mathcal{L}$ 为可逆 $\mathcal{O}_X$-模。
令 $n \geq 1$。则 $\mathcal{L}$ 丰沛，当且仅当
$\mathcal{L}^{\otimes n}$ 丰沛。
\end{lemma}

\begin{proof}
这由 $X_{s^n} = X_s$ 这一事实得出。
\end{proof}

\begin{lemma}
\label{properties-lemma-ample-on-closed}
令 $X$ 为概形。
令 $\mathcal{L}$ 为丰沛可逆 $\mathcal{O}_X$-模。
对任意闭子概形 $Z \subset X$，$\mathcal{L}$ 到 $Z$ 的限制丰沛。
\end{lemma}

\begin{proof}
这是显然的，因为拟紧空间的闭子集拟紧，而仿射概形的闭子概形仿射
（见《概形》引理 \ref{schemes-lemma-closed-immersion-affine-case}）。
\end{proof}

\begin{lemma}
\label{properties-lemma-affine-cap-s-open}
令 $X$ 为概形。令 $\mathcal{L}$ 为可逆 $\mathcal{O}_X$-模。
令 $s \in \Gamma(X, \mathcal{L})$。对任意仿射集
$U \subset X$，交 $U \cap X_s$ 是仿射的。
\end{lemma}

\begin{proof}
这转化为下述代数问题。
令 $R$ 为环。令 $N$ 为可逆 $R$-模
（即秩为 1 的局部自由模）。令 $s \in N$ 为元素。
则 $V = \{\mathfrak p \mid s \not \in \mathfrak p N\}$ 是
$\Spec(R)$ 的仿射开子集。

\medskip\noindent
令 $A = \bigoplus_{n \geq 0} A_n$ 为 $N$ 的对称代数
（它是交换的），并把 $s$ 看作 $A_1$ 的元素。
置 $B = A/(s - 1)A$。这是一个 $R$-代数，其构造与任意基变换
$R \to R'$ 交换。因此，
$B' = B \otimes_R R'$ 是零环，只要 $s$ 在
$N' = N \otimes_R R'$ 中映为零。
由此，若 $x \in \Spec(R) \setminus V$，则
$B \otimes_R \kappa(x) = 0$。
我们断定 $\Spec(B) \to \Spec(R)$ 通过 $V$ 分解，因为在
$x \not \in V$ 上的纤维为空。另一方面，若
$\Spec(R') \subset V$ 为仿射开集，则 $s$ 映到 $N'$ 的一个基元素，
并且可知 $B' = R'[s]/(s - 1) \cong R'$。
因此 $\Spec(B) \to V$ 是同构，而 $V$ 的确仿射。
\end{proof}

\begin{lemma}
\label{properties-lemma-ample-tensor-globally-generated}
\begin{reference}
\cite[II Proposition 4.5.6(ii)]{EGA}
\end{reference}
令 $X$ 为概形。令 $\mathcal{L}$ 和 $\mathcal{M}$
为可逆 $\mathcal{O}_X$-模。若
\begin{enumerate}
\item $\mathcal{L}$ 丰沛；
\item 开集 $X_t$，其中
$t \in \Gamma(X, \mathcal{M}^{\otimes m})$ 且 $m > 0$，覆盖 $X$，
\end{enumerate}
则 $\mathcal{L} \otimes \mathcal{M}$ 丰沛。
\end{lemma}

\begin{proof}
检验定义 \ref{properties-definition-ample} 的条件。
由于 $\mathcal{L}$ 丰沛，可知 $X$ 拟紧。
令 $x \in X$。选取 $n \geq 1$、$m \geq 1$、
$s \in \Gamma(X, \mathcal{L}^{\otimes n})$ 和
$t \in \Gamma(X, \mathcal{M}^{\otimes m})$，使得
$x \in X_s$、$x \in X_t$ 且 $X_s$ 仿射。
则
$s^mt^n \in \Gamma(X, (\mathcal{L} \otimes \mathcal{M})^{\otimes nm})$、
$x \in X_{s^mt^n}$，并且由引理
\ref{properties-lemma-affine-cap-s-open}，$X_{s^mt^n}$ 仿射。
\end{proof}

\begin{lemma}
\label{properties-lemma-affine-s-opens}
令 $X$ 为概形。令 $\mathcal{L}$ 为可逆 $\mathcal{O}_X$-模。
假设开集 $X_s$，其中
$s \in \Gamma(X, \mathcal{L}^{\otimes n})$ 且 $n \geq 1$，
构成 $X$ 的拓扑基。
则在这些开集之中，仿射的开集 $X_s$ 构成 $X$ 的拓扑基。
\end{lemma}

\begin{proof}
令 $x \in X$。选取仿射开邻域
$\Spec(R) = U \subset X$，它含有 $x$。由假设，存在
$n \geq 1$ 和 $s \in \Gamma(X, \mathcal{L}^{\otimes n})$，
使得 $X_s \subset U$。由上面的引理
\ref{properties-lemma-affine-cap-s-open}，交 $X_s = U \cap X_s$ 仿射。
由于 $U$ 可任意缩小，结论得证。
\end{proof}

\begin{lemma}
\label{properties-lemma-affine-s-opens-cover-quasi-separated}
令 $X$ 为概形，$\mathcal{L}$ 为可逆 $\mathcal{O}_X$-模。
假设对每个点 $x$（属于 $X$），都存在 $n \geq 1$ 和
$s \in \Gamma(X, \mathcal{L}^{\otimes n})$，使得
$x \in X_s$ 且 $X_s$ 仿射。则 $X$ 分离。
\end{lemma}

\begin{proof}
先证明 $X$ 拟分离。由假设，可以找到 $X$ 的覆盖，该覆盖由形如
$X_s$ 的仿射开集组成。由引理
\ref{properties-lemma-affine-cap-s-open}，任意两个这样的集合之交仿射，
所以《概形》引理
\ref{schemes-lemma-characterize-quasi-separated} 蕴含 $X$ 拟分离。

\medskip\noindent
为证明 $X$ 分离，可以使用赋值判据，即《概形》引理
\ref{schemes-lemma-valuative-criterion-separatedness}。
因此，令 $A$ 为分式域为 $K$ 的赋值环，并考虑两个态射
$f, g : \Spec(A) \to X$，使得两个复合
$\Spec(K) \to \Spec(A) \to X$ 相同。
由于 $A$ 是局部环，存在 $p, q \ge 1$、
$s \in \Gamma(X,
\mathcal{L}^{\otimes p})$ 和
$t \in \Gamma(X,
\mathcal{L}^{\otimes q})$，使得
$X_s$ 与 $X_t$ 仿射，
$f(\Spec A) \subseteq X_s$ 且 $g(\Spec A)
\subseteq X_t$。
现在把 $s$ 替换为 $s^q$，把 $t$ 替换为 $t^p$，
并把 $\mathcal{L}$ 替换为 $\mathcal{L}^{\otimes pq}$。
这是无害的，因为 $X_s = X_{s^q}$ 且 $X_t = X_{t^p}$，
而现在 $s$ 和 $t$ 都是同一个层 $\mathcal{L}$ 的截面。

\medskip\noindent
准凝聚模 $f^*\mathcal{L}$ 对应于 $A$-模 $M$，而
$g^*\mathcal{L}$ 对应于 $A$-模 $N$；这是由仿射概形上准凝聚模的分类
（《概形》引理 \ref{schemes-lemma-quasi-coherent-affine}）得出的。
$A$-模 $M$ 与 $N$ 都是秩为 $1$ 的局部自由模
（引理 \ref{properties-lemma-locally-free-module}），而由于 $A$ 局部，它们自由
（《代数》引理 \ref{algebra-lemma-K0-local}）。因此，可以把
$M$ 和 $N$ 等同为 $A$-子模，它们分别属于 $M \otimes_A K$ 和 $N
\otimes_A K$。等式
$f|_{\Spec(K)} = g|_{\Spec(K)}$ 确定一个同构
$\phi \colon M \otimes_A K \to N
\otimes_A K$。

\medskip\noindent
令 $x \in M$ 和 $y \in N$ 分别为把 $s$ 沿 $f$ 与 $g$ 拉回所得的元素。
它们满足 $\phi(x \otimes 1) = y \otimes 1$。
$f$ 的像包含于 $X_s$，故 $x \not\in \mathfrak{m}_A M$，
即 $x$ 生成 $M$。因此，$\phi$ 确定 $M$ 到 $N$ 中由
$y$ 生成的子模的同构。对称地使用 $t$，$\phi^{-1}$
确定 $N$ 到 $M$ 的一个子模的同构。因此，$\phi$ 限制为
$M$ 与 $N$ 之间的同构。由于 $x$ 生成 $M$，它的像 $y$ 生成
$N$，从而 $y \not\in \mathfrak{m}_A N$。所以
$g(\Spec(A)) \subseteq X_s$。因为 $X_s$ 仿射，由《概形》引理
\ref{schemes-lemma-affine-separated}，它分离，因而得出 $f = g$。
\end{proof}

\begin{lemma}
\label{properties-lemma-ample-separated}
令 $X$ 为概形。若 $X$ 上存在丰沛可逆层，则 $X$ 分离。
\end{lemma}

\begin{proof}
立即由引理 \ref{properties-lemma-affine-s-opens-cover-quasi-separated} 和定义
\ref{properties-definition-ample} 得出。
\end{proof}

\begin{lemma}
\label{properties-lemma-map-into-proj}
令 $X$ 为概形。
令 $\mathcal{L}$ 为可逆 $\mathcal{O}_X$-模。
把 $S = \Gamma_*(X, \mathcal{L})$ 视为分次环。
若 $X$ 的每个点都包含于某个开子概形 $X_s$，其中某个齐次
$s \in S_{+}$，则存在到 $S$ 的齐次谱（见《构造》第
\ref{constructions-section-proj} 节）的典范概形态射
$$
f : X \longrightarrow Y = \text{Proj}(S),
$$
它具有下列性质：
\begin{enumerate}
\item 有 $f^{-1}(D_{+}(s)) = X_s$，对任意齐次
$s \in S_{+}$；
\item 存在与乘法映射相容的 $\mathcal{O}_X$-模映射
$f^*\mathcal{O}_Y(n) \to \mathcal{L}^{\otimes n}$；见《构造》方程
(\ref{constructions-equation-multiply})；
\item 复合
$S_n \to \Gamma(Y, \mathcal{O}_Y(n)) \to \Gamma(X, \mathcal{L}^{\otimes n})$
是恒等映射；
\item 对每个 $x \in X$，都存在整数 $d \geq 1$ 以及开邻域
$U \subset X$，它包含 $x$，使得
$f^*\mathcal{O}_Y(dn)|_U \to \mathcal{L}^{\otimes dn}|_U$
都是同构，对所有 $n \in \mathbf{Z}$。
\end{enumerate}
\end{lemma}

\begin{proof}
以 $\psi : S \to \Gamma_*(X, \mathcal{L})$ 表示恒等映射。
将使用《构造》引理
\ref{constructions-lemma-invertible-map-into-proj} 的三元组
$(U(\psi), r_{\mathcal{L}, \psi}, \theta)$。
由假设，开子概形 $U(\psi)$ 的等于 $X$。因此
$r_{\mathcal{L}, \psi} : U(\psi) \to Y$ 在整个 $X$ 上有定义。
置 $f = r_{\mathcal{L}, \psi}$。
(2) 中的映射是 $\theta$ 的分量。
(3) 由上引引理中的条件 (2) 得出。
(1) 由 (3) 与上引引理中的条件 (1) 结合得出。
(4) 由《构造》引理
\ref{constructions-lemma-invertible-map-into-proj} 的最后一个陈述得出，
因为其中提到的映射 $\alpha$ 是同构。
\end{proof}

\begin{lemma}
\label{properties-lemma-map-into-proj-quasi-compact}
令 $X$ 为概形。令 $\mathcal{L}$ 为可逆 $\mathcal{O}_X$-模。
置 $S = \Gamma_*(X, \mathcal{L})$。
假设 (a) $X$ 的每个点都包含于某个开子概形 $X_s$，其中某个齐次
$s \in S_{+}$，并且 (b) $X$ 拟紧。则上面引理
\ref{properties-lemma-map-into-proj} 的典范概形态射
$f : X \longrightarrow \text{Proj}(S)$ 拟紧且像稠密。
\end{lemma}

\begin{proof}
为证明 $f$ 拟紧，只须说明 $f^{-1}(D_{+}(s))$ 拟紧，对任意齐次
$s \in S_{+}$。把
$X = \bigcup_{i = 1, \ldots, n} X_i$ 写成有限多个仿射开集之并。
由引理 \ref{properties-lemma-affine-cap-s-open}，每个交
$X_s \cap X_i$ 都仿射。因此
$X_s = \bigcup_{i = 1, \ldots, n} X_s \cap X_i$ 拟紧。
假设 $f$ 的像不稠密以导出矛盾。由于开集
$D_+(s)$（其中 $s \in S_+$ 齐次）构成
$\text{Proj}(S)$ 的拓扑基，可以找到这样的 $s$，使得
$D_+(s) \not = \emptyset$ 且 $f(X) \cap D_+(s) = \emptyset$。
由引理 \ref{properties-lemma-map-into-proj}，这意味着 $X_s = \emptyset$。
由引理 \ref{properties-lemma-invert-s-sections}，这意味着某个幂 $s^n$
是 $\mathcal{L}^{\otimes n\deg(s)}$ 的零截面。
继而这意味着 $D_+(s) = \emptyset$，即所需矛盾。
\end{proof}

\begin{lemma}
\label{properties-lemma-ample-immersion-into-proj}
令 $X$ 为概形。令 $\mathcal{L}$ 为可逆 $\mathcal{O}_X$-模。
置 $S = \Gamma_*(X, \mathcal{L})$。
假设 $\mathcal{L}$ 丰沛。则引理 \ref{properties-lemma-map-into-proj} 的典范概形态射
$f : X \longrightarrow \text{Proj}(S)$ 是像稠密的开浸入。
\end{lemma}

\begin{proof}
由引理 \ref{properties-lemma-affine-s-opens-cover-quasi-separated}，可知
$X$ 拟分离。选取有限多个齐次
$s_1, \ldots, s_n \in S_{+}$，使得 $X_{s_i}$ 仿射并且
$X = \bigcup X_{s_i}$。
设 $s_i$ 的次数为 $d_i$。$D_{+}(s_i)$ 在 $f$ 下的逆像是
$X_{s_i}$；见引理 \ref{properties-lemma-map-into-proj}。
由引理 \ref{properties-lemma-invert-s-sections}，环映射
$$
(S^{(d_i)})_{(s_i)} = \Gamma(D_{+}(s_i), \mathcal{O}_{\text{Proj}(S)})
\longrightarrow
\Gamma(X_{s_i}, \mathcal{O}_X)
$$
是同构。因此 $f$ 诱导同构
$X_{s_i} \to D_{+}(s_i)$。从而 $f$ 是 $X$ 到开子概形
$\bigcup_{i = 1, \ldots, n} D_{+}(s_i)$ 上的同构；该开子概形属于
$\text{Proj}(S)$。
像由引理 \ref{properties-lemma-map-into-proj-quasi-compact} 可知稠密。
\end{proof}

\begin{lemma}
\label{properties-lemma-open-in-proj-ample}
令 $X$ 为概形。
令 $S$ 为分次环。假设 $X$ 拟紧，并且存在开浸入
$$
j : X \longrightarrow Y = \text{Proj}(S).
$$
则 $j^*\mathcal{O}_Y(d)$ 是丰沛可逆层，其中某个 $d > 0$。
\end{lemma}

\begin{proof}
这是《构造》引理 \ref{constructions-lemma-ample-on-proj}。
\end{proof}

\begin{proposition}
\label{properties-proposition-characterize-ample}
令 $X$ 为拟紧概形。
令 $\mathcal{L}$ 为 $X$ 上的可逆层。
置 $S = \Gamma_*(X, \mathcal{L})$。
下列条件等价：
\begin{enumerate}
\item
\label{properties-item-ample}
$\mathcal{L}$ 丰沛；
\item
\label{properties-item-immersion}
开集 $X_s$（其中 $s \in S_{+}$ 齐次）覆盖 $X$，且相伴态射
$X \to \text{Proj}(S)$ 是开浸入；
\item
\label{properties-item-s-basis}
开集 $X_s$（其中 $s \in S_{+}$ 齐次）构成 $X$ 的拓扑基；
\item
\label{properties-item-s-affine-basis}
仿射的开集 $X_s$（其中 $s \in S_{+}$ 齐次）构成 $X$ 的拓扑基；
\item
\label{properties-item-qc-gg}
对每个准凝聚层 $\mathcal{F}$（在 $X$ 上），典范映射
$$
\Gamma(X, \mathcal{F} \otimes_{\mathcal{O}_X} \mathcal{L}^{\otimes n})
\otimes_{\mathbf{Z}} \mathcal{L}^{\otimes -n}
\longrightarrow
\mathcal{F}
$$
之像的和（其中 $n \geq 1$）等于 $\mathcal{F}$；
\item
\label{properties-item-qc-i-gg}
与 (\ref{properties-item-qc-gg}) 相同的性质，其中 $\mathcal{F}$
遍历所有准凝聚理想层；
\item
\label{properties-item-c-gg}
$X$ 拟分离，并且对每个有限型准凝聚层
$\mathcal{F}$（在 $X$ 上），都存在整数 $n_0$，使得
$\mathcal{F} \otimes_{\mathcal{O}_X} \mathcal{L}^{\otimes n}$
对所有 $n \geq n_0$ 都全局生成；
\item
\label{properties-item-c-q}
$X$ 拟分离，并且对每个有限型准凝聚层
$\mathcal{F}$（在 $X$ 上），都存在整数 $n > 0$、$k \geq 0$，使得
$\mathcal{F}$ 是 $k$ 个 $\mathcal{L}^{\otimes - n}$ 的直和的商；
\item
\label{properties-item-c-i-q}
与 (\ref{properties-item-c-q}) 相同，其中 $\mathcal{F}$ 遍历
$X$ 上的所有有限型理想层。
\end{enumerate}
\end{proposition}

\begin{proof}
引理 \ref{properties-lemma-ample-immersion-into-proj} 给出
(\ref{properties-item-ample}) $\Rightarrow$ (\ref{properties-item-immersion})。
引理 \ref{properties-lemma-ample-power-ample} 与
\ref{properties-lemma-open-in-proj-ample} 给出蕴含
(\ref{properties-item-ample}) $\Leftarrow$ (\ref{properties-item-immersion})。
蕴含
(\ref{properties-item-immersion}) $\Rightarrow$
(\ref{properties-item-s-affine-basis}) $\Rightarrow$ (\ref{properties-item-s-basis})
由《构造》第 \ref{constructions-section-proj} 节是显然的。
引理 \ref{properties-lemma-affine-s-opens} 给出
(\ref{properties-item-s-basis}) $\Rightarrow$ (\ref{properties-item-ample})。
因此，前四个条件全都等价。

\medskip\noindent
假设等价条件 (1)--(4) 成立。
特别地，注意 $X$ 分离（因为它是分离概形
$\text{Proj}(S)$ 的开子概形）。
令 $\mathcal{F}$ 为 $X$ 上的准凝聚层。
选取齐次 $s \in S_{+}$，使得 $X_s$ 仿射。
我们断言，任意截面 $m \in \Gamma(X_s, \mathcal{F})$
都属于上面 (\ref{properties-item-qc-gg}) 所陈列的某个映射之像。
由于这些仿射集 $X_s$ 覆盖 $X$，这将蕴含
(\ref{properties-item-qc-gg})。
事实上，由引理 \ref{properties-lemma-invert-s-sections}，可以把
$m$ 写成 $m' \otimes s^{-n}$ 的像，其中某个
$n \geq 1$ 和某个
$m' \in \Gamma(X, \mathcal{F} \otimes \mathcal{L}^{\otimes n})$。
这证明了断言。

\medskip\noindent
显然，(\ref{properties-item-qc-gg}) $\Rightarrow$ (\ref{properties-item-qc-i-gg})。
假设 (\ref{properties-item-qc-i-gg})，并证明 $\mathcal{L}$ 丰沛。
选取 $x \in X$。令 $U \subset X$ 为包含 $x$ 的仿射开集。
置 $Z = X \setminus U$。可把 $Z$ 看作约化闭子概形；见《概形》第
\ref{schemes-section-reduced} 节。
令 $\mathcal{I} \subset \mathcal{O}_X$ 为与闭子概形 $Z$
对应的准凝聚理想层。由假设 (\ref{properties-item-qc-i-gg})，存在
$n \geq 1$ 以及截面
$s \in \Gamma(X, \mathcal{I} \otimes \mathcal{L}^{\otimes n})$，
使得 $s$ 在 $x$ 处不消失（更准确地说，
$s \not \in \mathfrak m_x \mathcal{I}_x \otimes \mathcal{L}_x^{\otimes n}$）。
可把 $s$ 看作 $\mathcal{L}^{\otimes n}$ 的截面。
由于它显然沿 $Z$ 消失，可知 $X_s \subset U$。
因此 $X_s$ 仿射；见引理 \ref{properties-lemma-affine-cap-s-open}。
这证明了 $\mathcal{L}$ 丰沛。至此已经证明 (1)--(6) 等价。

\medskip\noindent
假设等价条件 (1)--(6) 成立。下文将不再另行说明地使用如下事实：
两个全局生成模层的张量积仍全局生成（见《模》引理
\ref{modules-lemma-tensor-product-globally-generated}）。
由 (1)，可找到元素 $s_i \in S_{d_i}$，其中 $d_i \geq 1$，使得
$X = \bigcup_{i = 1, \ldots, n} X_{s_i}$。
置 $d = d_1\ldots d_n$。由此，$\mathcal{L}^{\otimes d}$ 由
$$
s_1^{d/d_1}, \ldots, s_n^{d/d_n}.
$$
全局生成。这意味着，若 $\mathcal{L}^{\otimes j}$ 全局生成，
则 $\mathcal{L}^{\otimes j + dn}$ 也全局生成，对所有 $n \geq 0$。
固定 $j \in \{0, \ldots, d - 1\}$。对任意点 $x \in X$，存在
$n \geq 1$ 以及整体截面 $s$，它属于 $\mathcal{L}^{j + dn}$，
它在 $x$ 处不消失；这是把 (\ref{properties-item-qc-gg}) 应用于
$\mathcal{F} = \mathcal{L}^{\otimes j}$ 和丰沛可逆层
$\mathcal{L}^{\otimes d}$ 得出的。由于 $X$ 拟紧，可以找到有限列整数
$n_i$ 以及整体截面 $s_i$，它们属于
$\mathcal{L}^{\otimes j + dn_i}$，并且在 $X$ 的每一点并非全都消失。
由于 $\mathcal{L}^{\otimes d}$ 全局生成，这意味着
$\mathcal{L}^{\otimes j + dn}$ 全局生成，其中
$n = \max\{n_i\}$。因为对模 $d$ 的每个同余类都证明了这一点，
所以存在 $n_0 = n_0(\mathcal{L})$，使得
$\mathcal{L}^{\otimes n}$ 对所有 $n \geq n_0$ 都全局生成。
至此可知，若 $\mathcal{F}$ 全局生成，则
$\mathcal{F} \otimes \mathcal{L}^{\otimes n}$ 对所有
$n \geq n_0$ 也全局生成。

\medskip\noindent
继续假设等价条件 (1)--(6) 成立。
令 $\mathcal{F}$ 为有限型准凝聚 $\mathcal{O}_X$-模层。
以 $\mathcal{F}_n \subset \mathcal{F}$ 表示
(\ref{properties-item-qc-gg}) 中典范映射的像。依构造，
$\mathcal{F}_n \otimes \mathcal{L}^{\otimes n}$ 全局生成。
由 (\ref{properties-item-qc-gg})，$\mathcal{F}$ 是子层
$\mathcal{F}_n$（$n \geq 1$）之和。由《模》引理
\ref{modules-lemma-finite-type-quasi-compact-colimit}，可知
$\mathcal{F} = \sum_{n = 1, \ldots, N} \mathcal{F}_n$，其中某个
$N \geq 1$。因此，
$\mathcal{F} \otimes \mathcal{L}^{\otimes n}$ 就全局生成，只要
$n \geq N + n_0(\mathcal{L})$，其中 $n_0(\mathcal{L})$
如上。由此断定 (1)--(6) 蕴含 (\ref{properties-item-c-gg})。

\medskip\noindent
假设 (\ref{properties-item-c-gg})。令 $\mathcal{F}$ 为有限型准凝聚
$\mathcal{O}_X$-模层。由 (\ref{properties-item-c-gg})，存在整数
$n \geq 1$，使得典范映射
$$
\Gamma(X, \mathcal{F} \otimes_{\mathcal{O}_X} \mathcal{L}^{\otimes n})
\otimes_{\mathbf{Z}} \mathcal{L}^{\otimes -n}
\longrightarrow
\mathcal{F}
$$
满射。令 $I$ 为
$\Gamma(X, \mathcal{F} \otimes_{\mathcal{O}_X} \mathcal{L}^{\otimes n})$
的有限子集组成的集合，并按包含关系偏序。则 $I$ 是有向偏序集。
对 $i = \{s_1, \ldots, s_{r(i)}\}$，令
$\mathcal{F}_i \subset \mathcal{F}$ 为映射
$$
\bigoplus\nolimits_{j = 1, \ldots, r(i)} \mathcal{L}^{\otimes -n}
\longrightarrow
\mathcal{F}
$$
的像；在该映射中，乘以 $s_j$ 这一作用施加在第 $j$ 个因子上。
上面的满射性蕴含
$\mathcal{F} = \colim_{i \in I} \mathcal{F}_i$。
因此可应用《模》引理
\ref{modules-lemma-finite-type-quasi-compact-colimit}，并得出
$\mathcal{F} = \mathcal{F}_i$，其中某个 $i$。
所以已经证明 (\ref{properties-item-c-q})。换言之，
(\ref{properties-item-c-gg}) $\Rightarrow$ (\ref{properties-item-c-q})。

\medskip\noindent
蕴含 (\ref{properties-item-c-q}) $\Rightarrow$ (\ref{properties-item-c-i-q}) 是平凡的。

\medskip\noindent
最后，假设 (\ref{properties-item-c-i-q})。
令 $\mathcal{I} \subset \mathcal{O}_X$ 为准凝聚理想层。
由引理 \ref{properties-lemma-quasi-coherent-colimit-finite-type}
（这里使用了 $X$ 拟分离这一条件），可知
$\mathcal{I} = \colim_\alpha I_\alpha$，其中每个
$I_\alpha$ 都是有限型准凝聚层。由于依假设，每个 $I_\alpha$
都是 $\mathcal{L}$ 的负张量幂的商，所以对 $\mathcal{I}$
也得到同样的结论（当然，不再有幂的有限性或有界性）。
因此断定 (\ref{properties-item-c-i-q}) 蕴含 (\ref{properties-item-qc-i-gg})。
命题证毕。
\end{proof}

\begin{lemma}
\label{properties-lemma-ample-on-locally-closed}
设 $X$ 为概形。设 $\mathcal{L}$ 为丰沛可逆
$\mathcal{O}_X$-模。设 $i : X' \to X$ 为概形态射。
假设下列条件中至少有一个成立：
\begin{enumerate}
\item $i$ 是拟紧浸入；
\item $X'$ 拟紧，且 $i$ 是浸入；
\item $i$ 拟紧，并在 $X'$ 与 $i(X')$ 之间诱导同胚；
\item $X'$ 拟紧，并且 $i$ 在 $X'$ 与 $i(X')$ 之间诱导同胚。
\end{enumerate}
则 $i^*\mathcal{L}$ 在 $X'$ 上丰沛。
\end{lemma}

\begin{proof}
注意，在情形 (1) 与 (3) 中，概形 $X'$ 拟紧，
因为依定义 \ref{properties-definition-ample}，$X$ 拟紧。
所以只需证明 (2) 与 (4)。由于 (2) 是 (4) 的特殊情形，
只需证明 (4)。

\medskip\noindent
假设条件 (4) 成立。对于 $s \in \Gamma(X, \mathcal{L}^{\otimes d})$，
以 $s' = i^*s$ 表示 $s$ 在 $X'$ 上的拉回。注意，
$s'$ 是 $(i^*\mathcal{L})^{\otimes d}$ 的截面。由命题
\ref{properties-proposition-characterize-ample}，开集 $X_s$（其中
$s \in \Gamma(X, \mathcal{L}^{\otimes d})$）构成 $X$ 的拓扑基。
由《模》备注 \ref{modules-remark-pullback-s-open}，
$X'_{s'} = i^{-1}(X_s)$；又因为 $X' \to i(X')$ 是同胚，
所以开集 $X'_{s'}$ 构成 $X'$ 的拓扑基。因此，由命题
\ref{properties-proposition-characterize-ample}，$i^*\mathcal{L}$ 丰沛。
\end{proof}

\begin{lemma}
\label{properties-lemma-ample-on-product}
设 $S$ 为拟分离概形。设 $X$、$Y$ 为 $S$ 上的概形。
设 $\mathcal{L}$ 为丰沛可逆 $\mathcal{O}_X$-模，
并设 $\mathcal{N}$ 为丰沛可逆 $\mathcal{O}_Y$-模。
则 $\mathcal{M} = \text{pr}_1^*\mathcal{L}
\otimes_{\mathcal{O}_{X \times_S Y}} \text{pr}_2^*\mathcal{N}$
是 $X \times_S Y$ 上的丰沛可逆层。
\end{lemma}

\begin{proof}
态射 $i : X \times_S Y \to X \times Y$ 是拟紧浸入；见《概形》引理
\ref{schemes-lemma-fibre-product-after-map}。另一方面，$\mathcal{M}$
是沿 $i$ 拉回 $X \times Y$ 上相应可逆模层所得。
由引理 \ref{properties-lemma-ample-on-locally-closed}，只需对 $X \times Y$
证明该引理。我们检验定义 \ref{properties-definition-ample} 中
$\mathcal{M}$ 在 $X \times Y$ 上满足 (1) 与 (2)。

\medskip\noindent
由于 $X$ 和 $Y$ 都拟紧，$X \times Y$ 也拟紧。
令 $z \in X \times Y$ 为一点。令 $x \in X$ 和 $y \in Y$
为其投影。选取 $n > 0$ 以及
$s \in \Gamma(X, \mathcal{L}^{\otimes n})$，
使得 $X_s$ 是 $x$ 的仿射开邻域。选取 $m > 0$ 以及
$t \in \Gamma(Y, \mathcal{N}^{\otimes m})$，
使得 $Y_t$ 是 $y$ 的仿射开邻域。于是
$r = \text{pr}_1^*s \otimes \text{pr}_2^*t$ 是 $\mathcal{M}$ 的截面，且
$(X \times Y)_r = X_s \times Y_t$。这是 $z$ 的仿射开邻域，
证明完成。
\end{proof}







\section{仿射概形与拟仿射概形}
\label{properties-section-affine-quasi-affine}

\begin{lemma}
\label{properties-lemma-quasi-affine-O-ample}
设 $X$ 为概形。
则 $X$ 拟仿射当且仅当 $\mathcal{O}_X$ 丰沛。
\end{lemma}

\begin{proof}
假设 $X$ 拟仿射。置 $A = \Gamma(X, \mathcal{O}_X)$。
考虑由引理 \ref{properties-lemma-quasi-affine} 给出的如下开浸入：
$$
j : X \longrightarrow \Spec(A)
$$
注意，$\Spec(A) = \text{Proj}(A[T])$；见《构造》例
\ref{constructions-example-trivial-proj}。
所以可应用引理 \ref{properties-lemma-open-in-proj-ample}，断定
$\mathcal{O}_X$ 丰沛。

\medskip\noindent
假设 $\mathcal{O}_X$ 丰沛。注意，作为分次环，
$\Gamma_*(X, \mathcal{O}_X) \cong A[T]$。
所以，由引理 \ref{properties-lemma-ample-immersion-into-proj} 与
\ref{properties-lemma-quasi-affine}，并结合上面已见的等式
$\Spec(A) = \text{Proj}(A[T])$ 对任意环 $A$ 都成立这一事实，
即可得到结论。
\end{proof}

\begin{lemma}
\label{properties-lemma-quasi-affine-locally-closed}
设 $X$ 为拟仿射概形。对任意拟紧浸入
$i : X' \to X$，概形 $X'$ 都拟仿射。
\end{lemma}

\begin{proof}
这也可不使用丰沛可逆层的材料而直接证明；我们建议读者在餐巾纸上
自行完成。由于 $X$ 拟仿射，由引理
\ref{properties-lemma-quasi-affine-O-ample}，$\mathcal{O}_X$ 丰沛。
于是，由引理 \ref{properties-lemma-ample-on-locally-closed}，
$\mathcal{O}_{X'}$ 丰沛。再由引理
\ref{properties-lemma-quasi-affine-O-ample}，$X'$ 拟仿射。
\end{proof}

\begin{lemma}
\label{properties-lemma-characterize-affine}
设 $X$ 为概形。假设存在有限多个元素
$f_1, \ldots, f_n \in \Gamma(X, \mathcal{O}_X)$，使得
\begin{enumerate}
\item 每个 $X_{f_i}$ 都是 $X$ 的仿射开集；并且
\item $f_1, \ldots, f_n$ 在
$\Gamma(X, \mathcal{O}_X)$ 中生成的理想等于单位理想。
\end{enumerate}
则 $X$ 仿射。
\end{lemma}

\begin{proof}
假设有引理所述的 $f_1, \ldots, f_n$。
可写成 $1 = \sum g_i f_i$，其中某些
$g_j \in \Gamma(X, \mathcal{O}_X)$，所以显然
$X = \bigcup X_{f_i}$。
（$f_i$ 不可能在某点全都消失。）由于每个 $X_{f_i}$
都拟紧（因为它仿射），可知 $X$ 拟紧。所以，由上面的引理
\ref{properties-lemma-quasi-affine-O-ample}，$X$ 拟仿射。
考虑开浸入
$$
j : X \to \Spec(\Gamma(X, \mathcal{O}_X)),
$$
见引理 \ref{properties-lemma-quasi-affine}。右端标准开集
$D(f_i)$ 的逆像等于左端的 $X_{f_i}$，并且态射 $j$ 诱导同构
$X_{f_i} \cong D(f_i)$；见引理
\ref{properties-lemma-invert-f-affine}。由于 $f_i$ 生成单位理想，可知
$\Spec(\Gamma(X, \mathcal{O}_X))
= \bigcup_{i = 1, \ldots, n} D(f_i)$。所以 $j$ 是同构。
\end{proof}







\section{准凝聚层与丰沛可逆层}
\label{properties-section-ample-quasi-coherent}

\noindent
本节的主题是：在有丰沛可逆层时，每个准凝聚层都来自某个分次模。

\begin{situation}
\label{properties-situation-ample}
设 $X$ 为概形。
设 $\mathcal{L}$ 为 $X$ 上的丰沛可逆层。
置 $S = \Gamma_*(X, \mathcal{L})$，视为分次环。
置 $Y = \text{Proj}(S)$。
设 $f : X \to Y$ 为引理 \ref{properties-lemma-map-into-proj} 的典范态射。
它配备一个 $\mathbf{Z}$-分次 $\mathcal{O}_X$-代数映射
$\bigoplus f^*\mathcal{O}_Y(n) \to \bigoplus \mathcal{L}^{\otimes n}$。
\end{situation}

\noindent
下一个引理其实是再下一个引理的特殊情形，不过先明确指出它成立
似乎更妥当。

\begin{lemma}
\label{properties-lemma-ample-gcd-is-one}
在情形 \ref{properties-situation-ample} 中，典范态射 $f : X \to Y$
把 $X$ 映入开子概形 $W = W_1 \subset Y$；在其上，
$\mathcal{O}_Y(1)$ 可逆，并且所有乘法映射
$\mathcal{O}_Y(n) \otimes_{\mathcal{O}_Y} \mathcal{O}_Y(m) \to
\mathcal{O}_Y(n + m)$
都是同构（见《构造》引理
\ref{constructions-lemma-where-invertible}）。此外，所有映射
$f^*\mathcal{O}_Y(n) \to \mathcal{L}^{\otimes n}$ 都是同构。
\end{lemma}

\begin{proof}
由命题 \ref{properties-proposition-characterize-ample}，存在整数 $n_0$，使得
$\mathcal{L}^{\otimes n}$ 对所有 $n \geq n_0$ 都全局生成。
令 $x \in X$ 为一点。由上文，可找到
$a \in S_{n_0}$ 与 $b \in S_{n_0 + 1}$，使得
$a$ 和 $b$ 在 $x$ 处都不消失。因此
$f(x) \in D_{+}(a) \cap D_{+}(b) = D_{+}(ab)$。由《构造》引理
\ref{constructions-lemma-where-invertible}，可知
$f(x) \in W_1$，如所需。构造映射 $f$ 时曾使用《构造》引理
\ref{constructions-lemma-invertible-map-into-proj}；由该引理，映射
$f^*\mathcal{O}_Y(n_0) \to \mathcal{L}^{\otimes n_0}$ 与
$f^*\mathcal{O}_Y(n_0 + 1) \to \mathcal{L}^{\otimes n_0 + 1}$
在 $x$ 的某个邻域上是同构。结合代数结构的相容性以及 $f$ 映入 $W$
这一事实，可断定所有映射
$f^*\mathcal{O}_Y(n) \to \mathcal{L}^{\otimes n}$
在 $x$ 的某个邻域上都是同构。结论得证。
\end{proof}

\noindent
回顾《模》定义 \ref{modules-definition-gamma-star}：给定局部环层空间
$X$、可逆层 $\mathcal{L}$ 以及 $\mathcal{O}_X$-模 $\mathcal{F}$，
我们有分次 $\Gamma_*(X, \mathcal{L})$-模
$$
\Gamma_*(X, \mathcal{L}, \mathcal{F}) =
\bigoplus\nolimits_{n \in \mathbf{Z}}
\Gamma(X, \mathcal{F} \otimes_{\mathcal{O}_X} \mathcal{L}^{\otimes n}).
$$
下一个引理说明，在情形 \ref{properties-situation-ample} 中，可以从这个分次模
恢复准凝聚 $\mathcal{O}_X$-模 $\mathcal{F}$。还可参看《构造》引理
\ref{constructions-lemma-quasi-coherent-projective-space}，其中我们在
特殊情形 $X = \mathbf{P}^n_R$ 下证明了这个引理。

\begin{lemma}
\label{properties-lemma-ample-quasi-coherent}
在情形 \ref{properties-situation-ample} 中，设 $\mathcal{F}$ 为 $X$ 上的准凝聚层。
置 $M = \Gamma_*(X, \mathcal{L}, \mathcal{F})$，视为分次 $S$-模。
存在同构
$$
f^*\widetilde{M} \longrightarrow \mathcal{F}
$$
且关于 $\mathcal{F}$ 函子性成立，并使得
$M_0 \to \Gamma(\text{Proj}(S), \widetilde{M}) \to \Gamma(X, \mathcal{F})$
为恒等映射。
\end{lemma}

\begin{proof}
令 $s \in S_{+}$ 为齐次元素，使得 $X_s$ 是 $X$ 中的仿射开集。
回顾，$\widetilde{M}|_{D_{+}(s)}$ 对应于
$S_{(s)}$-模 $M_{(s)}$；见《构造》引理
\ref{constructions-lemma-proj-sheaves}。又回顾，
$f^{-1}(D_{+}(s)) = X_s$。由于 $X$ 带有丰沛可逆层，
它拟紧且拟分离；见第 \ref{properties-section-ample} 节。
由引理 \ref{properties-lemma-invert-s-sections}，存在典范同构
$M_{(s)} = \Gamma_*(X, \mathcal{L}, \mathcal{F})_{(s)} \to
\Gamma(X_s, \mathcal{F})$。由于 $\mathcal{F}$ 准凝聚，这给出典范同构
$$
f^*\widetilde{M}|_{X_s} \to \mathcal{F}|_{X_s}
$$
由于 $\mathcal{L}$ 在 $X$ 上丰沛，可知 $X$ 被形如 $X_s$ 的仿射开集
覆盖。因此，只需证明所显示的映射可在交叠上粘合。证明从略。
\end{proof}

\begin{remark}
\label{properties-remark-neurotic}
采用引理 \ref{properties-lemma-ample-quasi-coherent} 的假设与记号。
把该引理所显示的映射记为 $\theta_\mathcal{F}$。注意，引理
\ref{properties-lemma-ample-gcd-is-one} 的同构
$f^*\mathcal{O}_Y(n) \to \mathcal{L}^{\otimes n}$ 正是
$\theta_{\mathcal{L}^{\otimes n}}$。考虑乘法映射
$$
\widetilde{M} \otimes_{\mathcal{O}_Y} \mathcal{O}_Y(n)
\longrightarrow
\widetilde{M(n)}
$$
见《构造》方程
(\ref{constructions-equation-multiply-more-generally})。
把它拉回到 $X$，并考虑
$$
\xymatrix{
f^*\widetilde{M} \otimes_{\mathcal{O}_X} f^*\mathcal{O}_Y(n)
\ar[r]
\ar[d]_{\theta_\mathcal{F} \otimes \theta_{\mathcal{L}^{\otimes n}}}
&
f^*\widetilde{M(n)}
\ar[d]^{\theta_{\mathcal{F} \otimes \mathcal{L}^{\otimes n}}}
\\
\mathcal{F} \otimes \mathcal{L}^{\otimes n} \ar[r]^{\text{id}} &
\mathcal{F} \otimes \mathcal{L}^{\otimes n}
}
$$
这里使用了显然的等同
$M(n) = \Gamma_*(X, \mathcal{L}, \mathcal{F} \otimes \mathcal{L}^{\otimes n})$。
该图交换。证明从略。
\end{remark}

\noindent
下面的引理应当可以从引理 \ref{properties-lemma-ample-quasi-coherent}
推出（反之亦然），但直接重复证明似乎更简单。

\begin{lemma}
\label{properties-lemma-proj-quasi-coherent}
设 $S$ 为分次环，并且 $X = \text{Proj}(S)$ 拟紧。
设 $\mathcal{F}$ 为准凝聚 $\mathcal{O}_X$-模。置
$M = \bigoplus_{n \in \mathbf{Z}} \Gamma(X, \mathcal{F}(n))$，
视为分次 $S$-模；见《构造》第
\ref{constructions-section-invertible-on-proj} 节。《构造》引理
\ref{constructions-lemma-comparison-proj-quasi-coherent} 的映射
$$
\widetilde{M} \longrightarrow \mathcal{F}
$$
是同构。如果 $X$ 被标准开集 $D_+(f)$ 覆盖，其中 $f$ 的次数为 $1$，
则所诱导的映射
$M_n \to \Gamma(X, \mathcal{F}(n))$ 是恒等映射。
\end{lemma}

\begin{proof}
由于 $X$ 拟紧，可以找到正次数的齐次元素
$f_1, \ldots, f_n \in S$，使得
$X = D_+(f_1) \cup \ldots \cup D_+(f_n)$。令 $d$ 为
$f_1, \ldots, f_n$ 各次数的最小公倍数。用 $f_i$ 的某次幂替换它后，
可假设每个 $f_i$ 的次数都是 $d$。于是可知
$\mathcal{L} = \mathcal{O}_X(d)$ 可逆，乘法映射
$\mathcal{O}_X(ad) \otimes \mathcal{O}_X(bd) \to \mathcal{O}_X((a + b)d)$
都是同构，并且每个 $f_i$ 决定一个整体截面 $s_i$，它属于 $\mathcal{L}$，
使得 $X_{s_i} = D_+(f_i)$；见《构造》引理
\ref{constructions-lemma-where-invertible} 与
\ref{constructions-lemma-principal-open}。因此
$\Gamma(X, \mathcal{F}(ad)) =
\Gamma(X, \mathcal{F} \otimes \mathcal{L}^{\otimes a})$。
回顾，$\widetilde{M}|_{D_{+}(f_i)}$ 对应于
$S_{(f_i)}$-模 $M_{(f_i)}$；见《构造》引理
\ref{constructions-lemma-proj-sheaves}。由于 $f_i$ 的次数是 $d$，
$M_{(f_i)}$ 的同构类只依赖于 $M$ 中次数可被 $d$ 整除的齐次分量。
更准确地说，$M_{(f_i)}$ 的同构类只依赖于分次
$\Gamma_*(X, \mathcal{L})$-模
$\Gamma_*(X, \mathcal{L}, \mathcal{F})$，以及像 $s_i$，即 $f_i$ 在
$\Gamma_*(X, \mathcal{L})$ 中的像。概形 $X$ 依假设拟紧，
并且由《构造》引理
\ref{constructions-lemma-proj-separated} 分离。由引理
\ref{properties-lemma-invert-s-sections}，存在典范同构
$$
M_{(f_i)} = \Gamma_*(X, \mathcal{L}, \mathcal{F})_{(s_i)} \to
\Gamma(X_{s_i}, \mathcal{F}).
$$
于是，《构造》引理
\ref{constructions-lemma-comparison-proj-quasi-coherent}
中该映射的构造表明，它在 $D_+(f_i)$ 上是同构；
由于这些开集覆盖 $X$，它就是同构。最后一个陈述的证明从略。
\end{proof}




\section{寻找适当的仿射开集}
\label{properties-section-finding-affine-opens}

\noindent
本节汇集在一般程度各异的情形下有关仿射开集存在性的一些结果。

\begin{lemma}
\label{properties-lemma-maximal-points-affine}
设 $X$ 为拟分离概形。设 $Z_1, \ldots, Z_n$ 为 $X$
的两两不同的不可约分支；见《拓扑》第
\ref{topology-section-irreducible-components} 节。设
$\eta_i \in Z_i$ 为它们的泛点；见《概形》引理
\ref{schemes-lemma-scheme-sober}。存在仿射开邻域
$\eta_i \in U_i$，使得 $U_i \cap U_j = \emptyset$
对所有 $i \not = j$ 都成立。特别地，
$U = U_1 \cup \ldots \cup U_n$ 是包含所有点
$\eta_1, \ldots, \eta_n$ 的仿射开集。
\end{lemma}

\begin{proof}
令 $V_i$ 为任意仿射开集，它包含 $\eta_i$，并且与闭集
$Z_1 \cup \ldots \hat Z_i \ldots \cup Z_n$ 不交。
由于 $X$ 拟分离，对每个 $i$，并集
$W_i = \bigcup_{j, j \not = i} V_i \cap V_j$ 是 $V_i$
中不包含 $\eta_i$ 的拟紧开集。由《代数》引理
\ref{algebra-lemma-standard-open-containing-maximal-point}，
可找到开邻域 $U_i \subset V_i$，它包含 $\eta_i$ 且与 $W_i$ 不交。
最后，$U$ 仿射，因为它是环
$R_1 \times \ldots \times R_n$ 的谱，其中
$R_i = \mathcal{O}_X(U_i)$；见《概形》引理
\ref{schemes-lemma-disjoint-union-affines}。
\end{proof}

\begin{remark}
\label{properties-remark-maximal-points-affine}
若 $X$ 不拟分离，则上面的引理
\ref{properties-lemma-maximal-points-affine} 不成立。下面给出一例。取
$R = \mathbf{Q}[x, y_1, y_2, \ldots]/((x-i)y_i)$。
考虑极小素理想 $\mathfrak p = (y_1, y_2, \ldots)$，它是 $R$ 的理想。
把 $\Spec(R)$ 的两个副本沿着（非拟紧）开集
$\Spec(R) \setminus V(\mathfrak p)$ 粘合，
得到概形 $X$（按《概形》例
\ref{schemes-example-affine-space-zero-doubled} 的方式粘合）。
那么，$X$ 中对应于 $\mathfrak p$ 的两个极大点不包含在同一个
仿射开集中。理由是，$\Spec(R)$ 中任何包含 $\mathfrak p$ 的开集
都包含无穷多条“直线” $x = i$、$y_j = 0$，其中
$j \not = i$，其参数为 $y_i$。细节从略。
\end{remark}

\noindent
尽管有上面的例子，对于“大多数”有限族不可约闭子集，至少在
$X$ 拟紧时，仍可应用上面的引理
\ref{properties-lemma-maximal-points-affine}。这是因为 $X$ 含有稠密的分离开集。

\begin{lemma}
\label{properties-lemma-quasi-compact-dense-open-separated}
设 $X$ 为拟紧概形。存在分离的稠密开集 $V \subset X$。
\end{lemma}

\begin{proof}
设 $X = \bigcup_{i = 1, \ldots, n} U_i$ 是 $n$ 个仿射开子概形之并。
我们对 $n$ 归纳证明。$n = 1$ 时显然。设
$V' \subset \bigcup_{i = 1, \ldots, n - 1} U_i$ 为分离的稠密开子概形；
归纳假设保证它存在。考虑
$$
V = V' \amalg (U_n \setminus \overline{V'}).
$$
显然，$V$ 是 $X$ 的分离稠密开子概形。
\end{proof}

\noindent
事实表明，即使 $X$ 既拟分离又拟紧，也未必存在分离、拟紧的
稠密开集；见《例》引理
\ref{examples-lemma-no-dense-separated-quasi-compact-open-in-qcqs}。
下面稍微加强上面的引理 \ref{properties-lemma-maximal-points-affine}。

\begin{lemma}
\label{properties-lemma-point-and-maximal-points-affine}
设 $X$ 为拟分离概形。设 $Z_1, \ldots, Z_n$ 为 $X$
的两两不同的不可约分支。设 $\eta_i \in Z_i$ 为其泛点。
任取 $x \in X$。存在仿射开集 $U \subset X$，包含
$x$ 以及所有 $\eta_i$。
\end{lemma}

\begin{proof}
假设 $x \in Z_1 \cap \ldots \cap Z_r$，且
$x \not \in Z_{r + 1}, \ldots, Z_n$。于是可选取仿射开集
$W \subset X$，使得 $x \in W$，并且
$W \cap Z_i = \emptyset$ 对 $i = r + 1, \ldots, n$ 成立。
注意，显然 $\eta_i \in W$ 对 $i = 1, \ldots, r$ 成立。
由引理 \ref{properties-lemma-maximal-points-affine}，可选取两两不交的仿射开集
$U_i \subset X$，使得 $\eta_i \in U_i$ 对
$i = r + 1, \ldots, n$ 成立。由于 $X$ 拟分离，开集
$W \cap U_i$ 拟紧，而且不包含 $\eta_i$，其中
$i = r + 1, \ldots, n$。所以，由《代数》引理
\ref{algebra-lemma-standard-open-containing-maximal-point}，可缩小 $U_i$，
使得 $W \cap U_i = \emptyset$ 对
$i = r + 1, \ldots, n$ 成立。
于是并集
$U = W \cup \bigcup_{i = r + 1, \ldots, n} U_i$ 是不交并，因而
（由《概形》引理 \ref{schemes-lemma-disjoint-union-affines}）
是所需的仿射开集。
\end{proof}

\begin{lemma}
\label{properties-lemma-ample-finite-set-in-affine}
设 $X$ 为概形。假设下列条件之一成立：
\begin{enumerate}
\item 概形 $X$ 拟仿射；
\item 概形 $X$ 同构于某个仿射概形的局部闭子概形；
\item $X$ 上存在丰沛可逆层；
\item 概形 $X$ 同构于 $\text{Proj}(S)$ 的局部闭子概形，
其中 $S$ 为某个分次环。
\end{enumerate}
则对任意有限子集 $E \subset X$，存在仿射开集
$U \subset X$，使得 $E \subset U$。
\end{lemma}

\begin{proof}
由《性质》定义 \ref{properties-definition-quasi-affine}，拟仿射概形是某个仿射概形的
拟紧开子概形。任意仿射概形 $\Spec(R)$ 都同构于
$\text{Proj}(R[X])$，其中 $R[X]$ 以 $\deg(X) = 1$ 分次。
由命题 \ref{properties-proposition-characterize-ample}，若 $X$ 有丰沛可逆层，
则 $X$ 同构于 $\text{Proj}(S)$ 的开子概形，其中 $S$ 为某个分次环。
所以，只需在情形 (4) 中证明引理。
（我们建议读者自行直接证明情形 (2)。）

\medskip\noindent
因此，假设 $X \subset \text{Proj}(S)$ 是局部闭子概形，其中 $S$
是某个分次环。令 $T = \overline{X} \setminus X$。
回顾，标准开集 $D_{+}(f)$ 构成 $\text{Proj}(S)$ 的拓扑基。
由于 $E$ 有限，可选取有限多个齐次元素 $f_i \in S_{+}$，使得
$$
E \subset
D_{+}(f_1) \cup \ldots \cup D_{+}(f_n) \subset
\text{Proj}(S) \setminus T
$$
假设 $E = \{\mathfrak p_1, \ldots, \mathfrak p_m\}$
是 $\text{Proj}(S)$ 的子集。
考虑理想 $I = (f_1, \ldots, f_n) \subset S$。
由于 $I \not \subset \mathfrak p_j$ 对所有
$j = 1, \ldots, m$ 都成立，由《代数》引理
\ref{algebra-lemma-graded-silly} 可知，存在齐次元素
$f \in I$，并且 $f \not \in \mathfrak p_j$ 对所有
$j = 1, \ldots, m$ 都成立。于是 $E \subset D_{+}(f) \subset
D_{+}(f_1) \cup \ldots \cup D_{+}(f_n)$。由于 $D_{+}(f)$ 与
$T$ 不交，可知 $X \cap D_{+}(f)$ 是仿射概形
$D_{+}(f)$ 的闭子概形，因而是 $X$ 的仿射开集，如所需。
\end{proof}

\begin{lemma}
\label{properties-lemma-ample-finite-set-in-principal-affine}
设 $X$ 为概形。设 $\mathcal{L}$ 为 $X$ 上的丰沛可逆层。设
$$
E \subset W \subset X
$$
其中 $E$ 有限，且 $W$ 是 $X$ 中的开集。则存在 $n > 0$
以及截面 $s \in \Gamma(X, \mathcal{L}^{\otimes n})$，使得
$X_s$ 仿射，并且 $E \subset X_s \subset W$。
\end{lemma}

\begin{proof}
读者可修改引理 \ref{properties-lemma-ample-finite-set-in-affine} 的证明来证明本引理；
我们改从该引理推出它。由引理
\ref{properties-lemma-ample-finite-set-in-affine}，可选取仿射开集
$U \subset W$，使得 $E \subset U$。考虑分次环
$S = \Gamma_*(X, \mathcal{L}) =
\bigoplus_{n \geq 0} \Gamma(X, \mathcal{L}^{\otimes n})$。
对每个 $x \in E$，令 $\mathfrak p_x \subset S$ 为由在 $x$ 处消失的
截面组成的分次理想。显然，$\mathfrak p_x$ 是素理想；又因为
$\mathcal{L}$ 的某个幂全局生成，显然
$S_{+} \not \subset \mathfrak p_x$。令 $I \subset S$ 为由在
$X \setminus U$ 的所有点处消失的截面组成的分次理想。
由于集合 $X_s$ 构成拓扑基，可知
$I \not \subset \mathfrak p_x$ 对所有 $x \in E$ 都成立。
由（分次）素理想回避（《代数》引理
\ref{algebra-lemma-graded-silly}），可找到齐次 $s \in I$，
使得 $s \not \in \mathfrak p_x$ 对所有 $x \in E$ 都成立。
于是 $E \subset X_s \subset U$，并且由引理
\ref{properties-lemma-affine-cap-s-open}，$X_s$ 仿射。
\end{proof}

\begin{lemma}
\label{properties-lemma-quasi-affine-invertible-nonvanishing-section}
设 $X$ 为拟仿射概形。设 $\mathcal{L}$ 为可逆
$\mathcal{O}_X$-模。设 $E \subset W \subset X$，其中 $E$ 有限且
$W$ 开。则存在 $s \in \Gamma(X, \mathcal{L})$，使得
$X_s$ 仿射，并且 $E \subset X_s \subset W$。
\end{lemma}

\begin{proof}
本引理的证明与《代数》引理 \ref{algebra-lemma-silly} 的证明有许多
共同之处。设 $E = \{x_1, \ldots, x_n\}$。若
$E = W = \emptyset$，则 $s = 0$ 可用。若 $W \not = \emptyset$，
则必要时添加一点后，可假设 $E \not = \emptyset$。所以可假设
$n \geq 1$。我们对 $n$ 归纳证明本引理。

\medskip\noindent
基础情形：$n = 1$。用 $W$ 中 $x_1$ 的一个仿射开邻域替换 $W$ 后，
可假设 $W$ 仿射。结合引理 \ref{properties-lemma-quasi-affine-O-ample} 与命题
\ref{properties-proposition-characterize-ample}，可知每个准凝聚
$\mathcal{O}_X$-模都全局生成。因此，存在整体截面 $s$，
它属于 $\mathcal{L}$，且在 $x_1$ 处不消失。另一方面，令 $Z \subset X$ 为
$X \setminus W$ 上的约化诱导闭子概形。把全局生成性应用于
准凝聚理想层 $\mathcal{I}$，即 $Z$ 的理想层，可找到整体截面 $f$，
它属于 $\mathcal{I}$，
它在 $x_1$ 处不消失。于是 $s' = fs$ 是 $\mathcal{L}$ 的整体截面，
它在 $x_1$ 处不消失，并且 $X_{s'} \subset W$。由引理
\ref{properties-lemma-affine-cap-s-open}，$X_{s'}$ 仿射。

\medskip\noindent
归纳步骤：$n > 1$。若存在特化关系 $x_i \leadsto x_j$，其中
$i \not = j$，则只需对
$\{x_1, \ldots, x_n\} \setminus \{x_i\}$ 证明引理，归纳即得结论。
因此可假设各 $x_i$ 之间没有特化关系。由引理
\ref{properties-lemma-ample-finite-set-in-affine} 或引理
\ref{properties-lemma-ample-finite-set-in-principal-affine}，可假设 $W$ 仿射。
由归纳，可找到整体截面 $s$，它属于 $\mathcal{L}$，并使得
$X_s \subset W$ 仿射，并包含 $x_1, \ldots, x_{n - 1}$。
若 $x_n \in X_s$，则已经完成。假设 $s$ 在 $x_n$ 处为零。
由 $n = 1$ 的情形，可找到整体截面 $s'$，它属于 $\mathcal{L}$，并使得
$\{x_n\} \subset X_{s'} \subset
W \setminus \overline{\{x_1, \ldots, x_{n - 1}\}}$。
这里用到了 $x_n$ 不是 $x_1, \ldots, x_{n - 1}$ 的特化。
于是 $s + s'$ 是 $\mathcal{L}$ 的整体截面，它在
$x_1, \ldots, x_n$ 处都不消失，并且
$X_{s + s'} \subset W$；如前即可得出结论。
\end{proof}

\begin{lemma}
\label{properties-lemma-ring-affine-open-injective-local-ring}
设 $X$ 为概形，且 $x \in X$ 为一点。在下列每种情形中，都存在
仿射开邻域 $U \subset X$，它包含 $x$，并使得典范映射
$\mathcal{O}_X(U) \to \mathcal{O}_{X, x}$ 为单射：
\begin{enumerate}
\item $X$ 整；
\item $X$ 局部诺特；
\item $X$ 约化且只有有限多个不可约分支。
\end{enumerate}
\end{lemma}

\begin{proof}
翻译成代数后，这由《代数》引理
\ref{algebra-lemma-subring-of-local-ring} 得出。
\end{proof}

\begin{lemma}
\label{properties-lemma-standard-open-two-affines}
设 $U$、$V$ 为仿射概形，并设 $W \to U$ 与 $W \to V$
为开浸入。对任意 $w \in W$，存在仿射开邻域
$W' \subset W$，它包含 $w$，并使得 $W'$ 映至 $U$ 和 $V$ 中的标准开集。
\end{lemma}

\begin{proof}
令 $X$ 为把 $U$ 与 $V$ 沿 $W$ 粘合所得的概形，并应用《概形》引理
\ref{schemes-lemma-standard-open-two-affines}。
\end{proof}
